% Generated from paper_draft.tex by make_arxiv.py (review notes and
% comments stripped for arXiv). Do not edit; regenerate instead.
%
\documentclass{article}

\usepackage{hyperref}
\usepackage{amsmath}
\usepackage{mathtools}
\usepackage{amssymb}
\usepackage{amsthm}
\usepackage{dsfont}
\usepackage{stmaryrd}
\usepackage{xcolor}
\usepackage{listings}
\usepackage{booktabs}
\usepackage{microtype}
\usepackage{fullpage}
\usepackage{authblk}
\usepackage{aliascnt}
\usepackage{tikz}
\usepackage{needspace}

\setlength{\belowcaptionskip}{6pt}

\newtheorem{theorem}{Theorem}[section]

\newaliascnt{lemma}{theorem}
\newtheorem{lemma}[lemma]{Lemma}
\aliascntresetthe{lemma}

\newaliascnt{corollary}{theorem}
\newtheorem{corollary}[corollary]{Corollary}
\aliascntresetthe{corollary}

\newaliascnt{definition}{theorem}
\newtheorem{definition}[definition]{Definition}
\aliascntresetthe{definition}

\newaliascnt{example}{theorem}
\newtheorem{example}[example]{Example}
\aliascntresetthe{example}

\newaliascnt{remark}{theorem}
\newtheorem{remark}[remark]{Remark}
\aliascntresetthe{remark}

\usepackage[nameinlink]{cleveref}
\crefname{theorem}{Theorem}{Theorems}
\Crefname{theorem}{Theorem}{Theorems}
\crefname{lemma}{Lemma}{Lemmas}
\Crefname{lemma}{Lemma}{Lemmas}
\crefname{corollary}{Corollary}{Corollaries}
\Crefname{corollary}{Corollary}{Corollaries}
\crefname{definition}{Definition}{Definitions}
\Crefname{definition}{Definition}{Definitions}
\crefname{example}{Example}{Examples}
\Crefname{example}{Example}{Examples}
\crefname{remark}{Remark}{Remarks}
\Crefname{remark}{Remark}{Remarks}

\newcommand{\R}{\mathbb{R}}
\newcommand{\Q}{\mathbb{Q}}
\newcommand{\N}{\mathbb{N}}
\newcommand{\Fin}[1]{\mathsf{Fin}(#1)}
\newcommand{\Flag}[1]{\mathcal{F}^{#1}}
\newcommand{\FlagAlg}[1]{\mathcal{A}^{#1}}
\newcommand{\poshom}[1]{\mathrm{Hom}^+(\mathcal{A}^{#1}, \mathbb{R})}
\newcommand{\den}[2]{p(#1;\,#2)}
\newcommand{\tdensity}[2]{\pi(#1;\,#2)}
\newcommand{\lean}[1]{\texttt{#1}}
\newcommand{\emptyt}{\emptyset_{\mathrm{t}}}
\newcommand{\leansmul}{\mathbin{\vcenter{\hbox{$\scriptstyle\bullet$}}}}
\newcommand{\Ext}{\operatorname{Ext}}

\newcommand{\commentout}[1]{}
\newcommand\DOI[1]{DOI: \href{https://doi.org/#1}{#1}}

\tikzset{
  graph vertex/.style={circle, draw, fill=white, inner sep=0pt, minimum size=4.5pt},
  graph root/.style={circle, draw, fill=black, inner sep=0pt, minimum size=4.5pt},
  graph edge/.style={line width=0.45pt},
  graph nonedge/.style={line width=0.45pt, dashed},
  inline/.style={scale=0.20, baseline=-0.5ex},
  fv/.style={circle, fill=white, draw, inner sep=0pt, minimum size=3pt},
  fr/.style={circle, fill=black, inner sep=0pt, minimum size=3pt}
}

\newcommand{\fKtwobull}{%
  \begin{tikzpicture}[inline]
    \draw (-1,0) -- (1,0);
    \node[fr] at (-1,0) {};
    \node[fv] at (1,0) {};
  \end{tikzpicture}%
}
\newcommand{\fbarKtwobull}{%
  \begin{tikzpicture}[inline]
    \node[fr] at (-1,0) {};
    \node[fv] at (1,0) {};
  \end{tikzpicture}%
}
\newcommand{\fbarKthree}{%
  \begin{tikzpicture}[inline]
    \node[fv] at (90:1) {};
    \node[fv] at (210:1) {};
    \node[fv] at (330:1) {};
  \end{tikzpicture}%
}
\newcommand{\fbarPthree}{%
  \begin{tikzpicture}[inline]
    \draw (210:1) -- (330:1);
    \node[fv] at (90:1) {};
    \node[fv] at (210:1) {};
    \node[fv] at (330:1) {};
  \end{tikzpicture}%
}
\newcommand{\fPthree}{%
  \begin{tikzpicture}[inline]
    \draw (210:1) -- (90:1) -- (330:1);
    \node[fv] at (90:1) {};
    \node[fv] at (210:1) {};
    \node[fv] at (330:1) {};
  \end{tikzpicture}%
}
\newcommand{\fKthree}{%
  \begin{tikzpicture}[inline]
    \draw (210:1) -- (90:1) -- (330:1) -- (210:1);
    \node[fv] at (90:1) {};
    \node[fv] at (210:1) {};
    \node[fv] at (330:1) {};
  \end{tikzpicture}%
}
\newcommand{\fbarKthreebull}{%
  \begin{tikzpicture}[inline]
    \node[fr] at (90:1) {};
    \node[fv] at (210:1) {};
    \node[fv] at (330:1) {};
  \end{tikzpicture}%
}
\newcommand{\fEthreebull}{%
  \begin{tikzpicture}[inline]
    \draw (210:1) -- (330:1);
    \node[fr] at (90:1) {};
    \node[fv] at (210:1) {};
    \node[fv] at (330:1) {};
  \end{tikzpicture}%
}
\newcommand{\fbarPthreebull}{%
  \begin{tikzpicture}[inline]
    \draw (90:1) -- (330:1);
    \node[fr] at (90:1) {};
    \node[fv] at (210:1) {};
    \node[fv] at (330:1) {};
  \end{tikzpicture}%
}
\newcommand{\fPthreebull}{%
  \begin{tikzpicture}[inline]
    \draw (90:1) -- (330:1) -- (210:1);
    \node[fr] at (90:1) {};
    \node[fv] at (210:1) {};
    \node[fv] at (330:1) {};
  \end{tikzpicture}%
}
\newcommand{\fKonetwobull}{%
  \begin{tikzpicture}[inline]
    \draw (90:1) -- (210:1);
    \draw (90:1) -- (330:1);
    \node[fr] at (90:1) {};
    \node[fv] at (210:1) {};
    \node[fv] at (330:1) {};
  \end{tikzpicture}%
}
\newcommand{\fCfive}{%
  \begin{tikzpicture}[inline]
    \draw (90:1) -- (18:1) -- (-54:1) -- (-126:1) -- (162:1) -- cycle;
    \node[fv] at (90:1) {};
    \node[fv] at (18:1) {};
    \node[fv] at (-54:1) {};
    \node[fv] at (-126:1) {};
    \node[fv] at (162:1) {};
  \end{tikzpicture}%
}
\newcommand{\fCfivebull}{%
  \begin{tikzpicture}[inline]
    \draw (90:1) -- (18:1) -- (-54:1) -- (-126:1) -- (162:1) -- cycle;
    \node[fr] at (90:1) {};
    \node[fv] at (18:1) {};
    \node[fv] at (-54:1) {};
    \node[fv] at (-126:1) {};
    \node[fv] at (162:1) {};
  \end{tikzpicture}%
}

\lstdefinelanguage{Lean4}{
  keywords={def,theorem,lemma,instance,structure,class,import,open,namespace,
             end,where,fun,let,have,show,by,match,with,if,then,else,
             return,do,for,in,noncomputable,abbrev,variable,section,
             native_decide,decide,norm_num,simp,ring,linarith,omega,
             exact,apply,rw,calc,intro,constructor,use,refine,obtain,
             rcases,cases,induction,assumption,contradiction,trivial,rfl,
             macro,elab,notation,example,
             infix,infixl,infixr,prefix,postfix},
  keywordstyle=\color{blue!70!black}\bfseries,
  comment=[l]{--},
  morecomment=[l]{/--},
  commentstyle=\color{gray}\itshape,
  stringstyle=\color{orange!80!black},
  morestring=[b]",
  showstringspaces=false,
  sensitive=true,
  basicstyle=\ttfamily\small,
  breaklines=true,
  columns=fullflexible,
  escapeinside={(*@}{@*)},
  literate=
    {->}{{$\to$}}2
    {<-}{{$\leftarrow$}}2
    {<=}{{$\leq$}}1
    {>=}{{$\geq$}}1
    {alpha}{{$\alpha$}}1
    {forall}{{$\forall$}}1
    {exists}{{$\exists$}}1
}
\lstset{language=Lean4, frame=single, framesep=4pt,
        rulecolor=\color{black!20}, backgroundcolor=\color{gray!5}}

\title{Formalizing Flag Algebras in Lean}
\date{\today}

\author[1,4]{Gyeongwon Jeong\thanks{These authors are joint first authors of this work.}}
\author[1,4]{Seonghun Park\protect\footnotemark[1]}
\author[1]{Jihoon Hyun}
\author[2,3]{Sang-il Oum}
\author[4]{Hongseok~Yang}
\affil[1]{\small School of Computing, KAIST, Daejeon, Korea}
\affil[2]{\small Discrete Mathematics Group, Institute for Basic Science (IBS), Daejeon, Korea}
\affil[3]{\small Department of Mathematical Sciences, KAIST, Daejeon, Korea}
\affil[4]{\small School of Computational Sciences, Korea Institute for Advanced Study (KIAS), Seoul,~Korea}
\affil[ ]{E-mail addresses: \texttt{jgyw0910@kaist.ac.kr},
\texttt{hun57@kaist.ac.kr},
\texttt{qawbecrdtey@kaist.ac.kr},
\texttt{sagil@ibs.re.kr},
\texttt{hongseokyang@kias.re.kr}}

\begin{document}
\maketitle

\begin{abstract}
Razborov's flag algebra method is a powerful tool for proving asymptotic
inequalities in extremal graph theory, often reducing the task to finding a
finite certificate by semidefinite programming.  The method is
particularly well suited to constrained problems, where one seeks bounds for
graphs that avoid specified finite patterns, such as triangles.  We present a
machine-checked formalization of the method for finite simple graphs, together
with a certificate-to-proof compiler that turns externally generated
certificate data into algebraic proofs checked by Lean.

The formalization covers the mathematical foundations of the method: partially
labeled graphs, their densities in large graphs, the quotient algebra of
density expressions, graph-limit semantics through positive homomorphisms, and
the downward operators used to average out labels.  Our compiler treats the
output of the external semidefinite programming search as candidate
certificate data rather than as trusted input: lists of flags,
target-density data, and rational
positive-semidefinite matrices.  It then directs Lean to compute the required
density and multiplication facts independently and check those computations
against the mathematical definitions, verify positive semidefiniteness exactly
over $\mathbb{Q}$, and carry out the algebraic normalization steps that occur in
flag-algebra proofs.  In our case studies, the compiler yields machine-checked
proofs of seven Tur\'an-type upper bounds.  These include the upper bounds in
Mantel's theorem and the Erd\H{o}s pentagon theorem; no proof of the latter
that avoids flag algebras is known.  They also include a $C_4$-density bound
for triangle-free graphs and edge-density bounds for $K_4$-free, $K_5$-free,
and $C_5$-free graphs.  We use the underlying formalization directly as well,
without the compiler, to give a separate proof of Mantel's upper bound,
formalize the matching constructions that complete both exact
Tur\'an-density results, and prove two inequalities of Goodman.

The design of our constrained semantics also prompted a separate
meta-theoretic investigation comparing two ways of imposing graph constraints:
building a hereditary constraint into the flag algebra from the start, or
testing inequalities afterward on constrained graph limits with labels chosen
at random.  We briefly state the resulting root-plantability criterion
that characterizes when
these approaches agree; a forthcoming paper will present the complete
mathematical and formal account.
\end{abstract}

\section{Introduction}
\label{sec:intro}

\paragraph{Why flag algebras matter.}
Many central questions in extremal combinatorics ask how often one finite
pattern can occur when another pattern is forbidden.  Let $F$ be a finite
simple graph and $\mathcal{H}$ a finite collection of finite simple graphs.  We write
\[
  \mathrm{ex}(n,\,F;\,\mathcal{H})
  :=
  \max\bigl\{
    \text{number of induced copies of }F\text{ in }G
    \mid |V(G)|=n,\; G\text{ is }\mathcal{H}\text{-free}
  \bigr\},
\]
where an \emph{induced copy} of $F$ in $G$ is a vertex subset
$S\subseteq V(G)$ with $|S|=|V(F)|$ and $G[S]\cong F$, and
\emph{\(\mathcal{H}\)-free} means that $G$ contains no copy of any
member of $\mathcal{H}$ as a (not necessarily induced) subgraph.  We obtain the
corresponding asymptotic quantity by normalizing
$\mathrm{ex}(n,\,F;\,\mathcal{H})$ by the total number
$\binom{n}{|V(F)|}$ of $|V(F)|$-vertex subsets and taking the limit; we call the
result the \emph{Tur\'an density} of $F$ with respect to $\mathcal{H}$:
\[
  \tdensity{F}{\mathcal{H}}
  :=
  \lim_{n\to\infty}
  \frac{\mathrm{ex}(n,\,F;\,\mathcal{H})}{\binom{n}{|V(F)|}}.
\]
For example, Mantel's theorem~\cite{mantel} states that a triangle-free graph on $n$
vertices has at most $\lfloor n^2/4\rfloor$ edges.  Equivalently,
$\mathrm{ex}(n,\,K_2;\,K_3)=\lfloor n^2/4\rfloor$, and the bound is achieved by the
complete bipartite graph $K_{\lfloor n/2\rfloor,\lceil n/2\rceil}$.  The
asymptotic value
\[
  \tdensity{K_2}{K_3}
  = \lim_{n\to\infty}\frac{\lfloor n^2/4\rfloor}{\binom{n}{2}}
  = \frac{1}{2}
\]
records that in any sufficiently large triangle-free graph, at most half of
all vertex pairs are edges.  (Whenever $\mathcal{H}=\{H\}$ is a singleton, we
write $H$ in place of $\mathcal{H}$; so $K_3$ above abbreviates $\{K_3\}$.)

Determining $\tdensity{F}{\mathcal{H}}$ is hard because it requires controlling
\emph{all} sufficiently large \(\mathcal{H}\)-free graphs: a valid upper
bound must guarantee that, asymptotically, the proportion of induced copies of
$F$ never exceeds the claimed value, regardless of the size of the host graph.  No
finite collection of examples suffices, and the extremal graphs achieving the
maximum may change unpredictably as $n$ grows.

Razborov's flag algebra method~\cite{razborov2007flag} provides a proof technique
for exactly this setting: it reduces the task of proving an asymptotic bound
for an infinite class of graphs to checking finitely many algebraic conditions.
A semidefinite programming (SDP) solver proposes a finite package of data
intended to satisfy these conditions; we call this package a candidate
certificate.  The flag algebra
framework explains why a candidate certificate that passes the checks
establishes the desired bound for the whole class.

This combination of mathematical abstraction and computer search has been
unusually effective.  It proved the Erd\H{o}s pentagon
theorem~\cite{hatami2012,grzesik2012}, concerning the maximum asymptotic
density of pentagons in triangle-free graphs, and determined the minimum
triangle density for a given edge density~\cite{razborov2008triangles}.  It
determined the maximum induced $C_5$ density~\cite{balogh2016induced5} and
settled the rainbow-triangle density problem~\cite{balogh2017rainbow}.  It has
also produced the best known upper bound on Tur\'an's tetrahedron
problem~\cite{razborov2010tetrahedron} and several exact
Tur\'an-density results for 3-uniform hypergraphs~\cite{baber2012turan}.

\paragraph{Formalizing the method.}
We first formalize in Lean the mathematical foundations of Razborov's flag
algebra method for finite simple graphs.  This development includes partially
labeled graphs and their densities in large graphs, an algebra of density
expressions modulo the identities that relate different graph sizes, positive
homomorphisms representing graph limits, and the downward operator that turns
a labeled density expression into an unlabeled one.  We call this part of the
development the \emph{specification layer} (\Cref{sec:abstract}).  Its
definitions follow the standard mathematics closely: flags are isomorphism
classes of labeled graphs, densities are sampling probabilities, and algebraic
expressions are identified when expansion identities give them the same
asymptotic meaning.  This layer provides the mathematical interface needed to
express flag-algebra proofs in Lean much as they appear on paper.

\paragraph{From certificates to formal proofs.}
Stating the method is not the same as checking a concrete application.  Each
application is supported by a candidate certificate that is compact to state
but potentially costly to audit: the Erd\H{o}s pentagon certificate relies
on thousands of exact density values and on rational matrices whose
positive semidefiniteness must be established exactly.  Such data are
impractical to verify by hand, yet a formal proof should neither trust them nor
require thousands of individually written proofs.  We close this gap with two
further layers.  The
\emph{reflection layer} (\Cref{sec:reflection}) supplies executable graph
representations, together with adequacy theorems proving that the computations
they support agree with the specification.  Generation commands then
specialize these theorems at elaboration time to produce the verified density,
downward, and multiplication identities required by a certificate.  The
\emph{automation layer} (\Cref{sec:flagmatic}) contains the
certificate-to-proof compiler.  It checks exact rational $LDL^\top$
factorizations to prove positive semidefiniteness over $\mathbb{R}$, expands
the target when necessary, and uses the reflected identities and custom
tactics to normalize the certificate calculation under the forbidden-graph
assumption and assemble the final bound.

\paragraph{Constrained semantics.}
Our formalization also makes a choice about how to impose a forbidden-graph
condition.  In the usual built-in approach, one constructs the flag algebra
using only \(H\)-free graphs.  Our implementation instead retains the ambient
algebra of all finite graphs and imposes \(H\)-freeness when an inequality is
interpreted.  For a typed inequality, \lean{forbidLE} considers \(H\)-free
empty-type positive homomorphisms, chooses an occurrence of the flag type
uniformly at random, uses its vertices as labels, and requires the inequality
to hold with probability one.  \Cref{sec:metatheory} calls this ensemble
semantics and compares it with the built-in (or quotient) semantics.  Every
quotient-valid inequality is ensemble-valid.  The two orders agree at the
empty flag type, but they can differ at a nonempty flag type; they agree for all
inequalities at a fixed flag type exactly when the constrained class is
root-plantable at that flag type.  We prove that blow-up closure guarantees this
agreement and give a \(C_4\)-free counterexample showing that heredity alone
does not.

\paragraph{Contributions.}
The paper makes four main contributions.  First, it formalizes in Lean the core
mathematical theory of flag algebras for finite simple graphs: flags and their
densities, the quotient algebra, positive homomorphisms and semantic order,
and the downward operator and random-extension machinery.  Second, it develops
an executable reflection layer, proves that its computations agree with the
mathematical specification, and uses it to build a certificate-to-proof compiler
that translates externally generated rational certificate data into algebraic
proofs checked by Lean.  Third, it evaluates the compiler on seven upper-bound
certificates, including the bounds needed for Mantel's theorem and the
Erd\H{o}s pentagon theorem, as well as a bound on the \(C_4\)-density of
triangle-free graphs.  A separate direct proof of Mantel's theorem shows how
to work with the flag-algebra formalization without the compiler.  Matching
lower bounds for Mantel's theorem and the Erd\H{o}s pentagon theorem complete
both Tur\'an-density results, while two inequalities of Goodman further
illustrate uses of the same infrastructure outside the certificate pipeline.
Fourth, it compares the
semantics obtained by building a hereditary constraint into the flag algebra
with the ensemble semantics used by our implementation.  This comparison
yields the root-plantability criterion, a sufficient condition based on
blow-up closure, and a \(C_4\)-free counterexample.  The meta-theory is
formalized in a separate, automatically generated Lean development that
extends the manually written flag-algebra library.  The complete Lean
development, including the certificate-to-proof compiler and this
autoformalized meta-theory, is publicly available at
\url{https://github.com/taeyool/lean-flag-algebras-release}.

\paragraph{Organization.}
\Cref{sec:background} reviews the mathematical background of flag algebras and
establishes the notation used throughout the paper.  \Cref{sec:abstract}
presents the specification layer, including its semantic treatment of
forbidden-graph assumptions.  \Cref{sec:reflection} introduces executable
graph representations, proves that their computations agree with the
specification, and explains how elaboration-time commands generate flags and
their theorems.  \Cref{sec:flagmatic} describes the
certificate-to-proof compiler, its exact checks of positive semidefiniteness
and proof automation, and its evaluation on seven certificates.
\Cref{sec:lowerbounds} illustrates direct uses of the underlying formalization
beyond certificate compilation: a direct proof of Mantel's theorem, matching
lower-bound constructions for Mantel's theorem and the Erd\H{o}s pentagon
theorem, and two inequalities of Goodman.  \Cref{sec:obstacles} discusses
the main engineering obstacles and a recurring bijective proof pattern for
finite counting identities.  \Cref{sec:metatheory} presents the core results
of our meta-theory of constrained semantics and describes its separate
autoformalization, while
\Cref{sec:unsure} records the remaining design questions.  Finally,
\Cref{sec:related} discusses related work, \Cref{sec:conclusion} presents our
conclusions, and \Cref{sec:compile-times} reports compilation times for the seven
compiler examples.

\section{Background: Flag Algebras}
\label{sec:background}

This section recalls the mathematical definitions underlying flag algebras,
following Razborov~\cite{razborov2007flag}, and establishes the notation used
later in the paper.  No background in flag algebras or extremal combinatorics
is assumed.  For a non-negative integer $n$, we write $[n]$ for the set of
positive integers at most $n$.  We consider only undirected, finite, and simple
graphs; for a graph $G$, its vertex and edge sets are denoted by $V(G)$ and $E(G)$,
respectively.
We use the standard graph names: $K_n$ is the complete graph on $n$ vertices,
$P_n$ the path and $C_n$ the cycle on $n$ vertices, $\overline{G}$ the
complement of $G$, and $K_{a,b}$ the complete bipartite graph with parts of
sizes $a$ and $b$.

We fix a (possibly empty) finite set $\mathcal{H}$ of forbidden graphs.  In this
paper, an unqualified \emph{subgraph} is an ordinary, not necessarily induced,
subgraph.  A graph is \emph{\(\mathcal{H}\)-free} if it contains no subgraph
isomorphic to any member of $\mathcal{H}$.  All flag types, flags, and flag
algebras defined below are relative to this fixed $\mathcal{H}$.

\subsection{Flag Types and Flags}

\paragraph{Flag Types.}
A flag type specifies the labeled part shared by all flags of that flag type.
Formally, a \emph{flag type} of size $k$ is an \(\mathcal{H}\)-free graph
\(\sigma\) with vertex set $[k]$.  The \emph{empty flag type} $\emptyset$, of
size zero, is the unique graph on the empty vertex set.
We use the term \emph{flag type} rather than Razborov's original \emph{type}
to avoid confusion with Lean's type system in later sections.

\paragraph{Flags.}
For a graph $G$ and a set $S\subseteq V(G)$, let $G[S]$ denote the
\emph{induced subgraph} of $G$ on $S$: its vertex set is $S$, and its edges are
exactly the edges of $G$ with both endpoints in $S$.
Fix a flag type $\sigma$ of size $k$.
A \emph{$\sigma$-flag} is a pair $F = (M, \theta_F)$ consisting of an
\(\mathcal{H}\)-free finite graph $M$ and an injective map
$\theta_F : [k] \hookrightarrow V(M)$ that induces a graph isomorphism from
$\sigma$ to $M[\operatorname{Im}(\theta_F)]$.  The \emph{size} of $F$ is
$|V(M)|$, and we write $V(F)$ for $V(M)$.
For each $i\in[k]$, the vertex $\theta_F(i)$ carries label $i$.  The vertices
in~$\operatorname{Im}(\theta_F)$ are therefore called the
\emph{labeled vertices} of $F$; they are also often called its \emph{roots}.
All remaining vertices are \emph{unlabeled}.
For a set $S \subseteq V(F)$ with $\operatorname{Im}(\theta_F) \subseteq S$,
we write $F[S]$ for the $\sigma$-flag $(M[S], \theta_F)$ and call it the
\emph{$\sigma$-flag induced by $F$ on $S$}.
Flags over the empty flag type $\emptyset$ are simply ordinary
\(\mathcal{H}\)-free graphs with no labeled vertices.
\Cref{fig:flag-examples} illustrates several flags of different flag types.

Two $\sigma$-flags $(M_1,\theta_1)$ and $(M_2,\theta_2)$ are
\emph{isomorphic}, denoted by $(M_1,\theta_1) \cong_\sigma (M_2,\theta_2)$,
if there is a graph isomorphism $\alpha: V(M_1)\to V(M_2)$
with $\alpha \circ \theta_1 = \theta_2$; this equation says that $\alpha$
preserves every label.  The set of isomorphism classes of
$\sigma$-flags of size $n$ is written $\mathcal{F}^\sigma_n$, and
$\mathcal{F}^\sigma = \bigcup_{n \geq k} \mathcal{F}^\sigma_n$.
Each $\mathcal{F}^\sigma_n$ is finite, although
$\mathcal{F}^\sigma$ itself is infinite.
When a flag type~$\sigma$ is fixed, we consistently call an individual flag
of that flag type a $\sigma$-flag.  We reserve the unqualified term \emph{flag} for
statements in which the flag type is unspecified, varies, or is irrelevant.

\begin{figure}[t]
  \centering
  \setlength{\tabcolsep}{1.2em}
  \renewcommand{\arraystretch}{1.3}
  \begin{tabular}{cccccc}
    \begin{tikzpicture}[scale=0.8, baseline=-0.5ex]
      \node[graph vertex] (a) at (0,0) {};
      \node[graph vertex] (b) at (0.8,0) {};
      \draw[graph edge] (a) -- (b);
    \end{tikzpicture}
    &
    \begin{tikzpicture}[scale=0.8, baseline=-0.5ex]
      \node[graph vertex] (a) at (90:0.65) {};
      \node[graph vertex] (b) at (210:0.65) {};
      \node[graph vertex] (c) at (330:0.65) {};
    \end{tikzpicture}
    &
    \begin{tikzpicture}[scale=0.8, baseline=-0.5ex]
      \node[graph vertex] (a) at (90:0.65) {};
      \node[graph vertex] (b) at (210:0.65) {};
      \node[graph vertex] (c) at (330:0.65) {};
      \draw[graph edge] (b) -- (c);
    \end{tikzpicture}
    &
    \begin{tikzpicture}[scale=0.8, baseline=-0.5ex]
      \node[graph vertex] (a) at (90:0.65) {};
      \node[graph vertex] (b) at (210:0.65) {};
      \node[graph vertex] (c) at (330:0.65) {};
      \draw[graph edge] (a) -- (b);
      \draw[graph edge] (a) -- (c);
    \end{tikzpicture}
    &
    \begin{tikzpicture}[scale=0.8, baseline=-0.5ex]
      \node[graph vertex] (a) at (90:0.65) {};
      \node[graph vertex] (b) at (210:0.65) {};
      \node[graph vertex] (c) at (330:0.65) {};
      \draw[graph edge] (a) -- (b);
      \draw[graph edge] (a) -- (c);
      \draw[graph edge] (b) -- (c);
    \end{tikzpicture}
    &
    \begin{tikzpicture}[scale=0.7, baseline=-0.5ex]
      \node[graph vertex] (v1) at (90:0.65) {};
      \node[graph vertex] (v2) at (18:0.65) {};
      \node[graph vertex] (v3) at (-54:0.65) {};
      \node[graph vertex] (v4) at (-126:0.65) {};
      \node[graph vertex] (v5) at (162:0.65) {};
      \draw[graph edge] (v1) -- (v2) -- (v3) -- (v4) -- (v5) -- (v1);
    \end{tikzpicture}
    \\
    $K_2$ & $\overline{K_3}$ & $\overline{P_3}$ & $P_3$ & $K_3$ & $C_5$ \\[0.8em]
    \begin{tikzpicture}[scale=0.8, baseline=-0.5ex]
      \node[graph root] (a) at (0,0) {};
      \node[graph vertex] (b) at (0.8,0) {};
      \draw[graph edge] (a) -- (b);
    \end{tikzpicture}
    &
    \begin{tikzpicture}[scale=0.8, baseline=-0.5ex]
      \node[graph root] (a) at (0,0) {};
      \node[graph vertex] (b) at (0.8,0) {};
    \end{tikzpicture}
    &
    \begin{tikzpicture}[scale=0.8, baseline=-0.5ex]
      \node[graph vertex] (a) at (90:0.65) {};
      \node[graph root] (b) at (210:0.65) {};
      \node[graph vertex] (c) at (330:0.65) {};
      \draw[graph edge] (b) -- (c);
    \end{tikzpicture}
    &
    \begin{tikzpicture}[scale=0.8, baseline=-0.5ex]
      \node[graph vertex] (a) at (90:0.65) {};
      \node[graph root] (b) at (210:0.65) {};
      \node[graph vertex] (c) at (330:0.65) {};
      \draw[graph edge] (b) -- (a);
      \draw[graph edge] (a) -- (c);
    \end{tikzpicture}
    &
    \begin{tikzpicture}[scale=0.8, baseline=-0.5ex]
      \node[graph vertex] (a) at (90:0.65) {};
      \node[graph root] (b) at (210:0.65) {};
      \node[graph vertex] (c) at (330:0.65) {};
      \draw[graph edge] (b) -- (a);
      \draw[graph edge] (b) -- (c);
    \end{tikzpicture}
    &
    \begin{tikzpicture}[scale=0.7, baseline=-0.5ex]
      \node[graph root] (v1) at (90:0.65) {};
      \node[graph vertex] (v2) at (18:0.65) {};
      \node[graph vertex] (v3) at (-54:0.65) {};
      \node[graph vertex] (v4) at (-126:0.65) {};
      \node[graph vertex] (v5) at (162:0.65) {};
      \draw[graph edge] (v1) -- (v2) -- (v3) -- (v4) -- (v5) -- (v1);
    \end{tikzpicture}
    \\
    $K_2^\bullet$ & $\overline{K_2}^\bullet$ & $\overline{P_3}^\bullet$ & $P_3^\bullet$
    & ${P_3'}^\bullet$ & $C_5^\bullet$ \\
  \end{tabular}
  \caption{Examples of flags when no graphs are forbidden
  ($\mathcal{H}=\varnothing$).  \emph{Top row}: $\emptyset$-flags, equivalently
  ordinary graphs, of sizes 2, 3, and 5.  \emph{Bottom row}: $\sigma$-flags of the
  unique one-vertex flag type $\sigma$.  The filled vertex carries label~1,
  and the open vertices are unlabeled.  The $\sigma$-flags $P_3^\bullet$ and
  ${P_3'}^\bullet$ have the same underlying path; $P_3^\bullet$ labels an
  endpoint, whereas ${P_3'}^\bullet$ labels the middle vertex.}
  \label{fig:flag-examples}
\end{figure}

\subsection{Subflag Densities}

In a density calculation, flags play two roles.  The \emph{pattern flag} is the
small local configuration whose occurrence is being measured, while the
\emph{host flag} is the larger flag in which we sample vertices.  These are
roles in the calculation, not different kinds of flags.  We use \(F\), or
\(F_i\), for pattern flags, \(G\) for host flags, and \(G'\) for an intermediate
host flag when a third flag variable is needed.

Let $\sigma$ be a flag type of size $k$.  

\paragraph{Single-flag density.}
For a $\sigma$-flag $F$ of size $m$
(the pattern $\sigma$-flag) and a $\sigma$-flag $G = (M, \theta_G)$ of size $n \geq m$
(the host $\sigma$-flag), the \emph{density of $F$ in $G$}, written $\den{F}{G}$, is
the probability that a uniformly random $(m-k)$-element subset of
$V(G) \setminus \operatorname{Im}(\theta_G)$, together with the $k$ labeled
vertices of $G$, induces a $\sigma$-flag \(F'\) with \(F' \cong_\sigma F\).
Explicitly:
\[
  \den{F}{G}
  \;:=\;
  \frac{1}{\binom{n-k}{m-k}}
    \cdot {\Bigl|\Big\{S \subseteq V(G)\setminus \operatorname{Im}(\theta_G)
                  ~\Big|~ |S|=m-k,\;
                  G[S \cup \operatorname{Im}(\theta_G)]
                  \cong_\sigma F \Big\}\Bigr|}.
\]
This value lies in $[0,1]$ and is invariant under
isomorphism of both $F$ and $G$.
When $n < m$, we set $\den{F}{G} := 0$ by convention.

\begin{example}[Subflag densities in $C_5^\bullet$]
  Among the four unlabeled vertices of
  $\fCfivebull$, exactly two are adjacent to the labeled one.  Hence
  \[
    \den{\fKtwobull}{\fCfivebull} = \frac{2}{4} = \frac12,
    \qquad
    \den{\fbarKtwobull}{\fCfivebull} = \frac{2}{4} = \frac12.
  \]
\end{example}

\paragraph{Chain rule.}
The density $\den{F}{G}$ can be computed in two equivalent ways: by sampling
$m-k$ unlabeled vertices directly from $G$, or by first sampling an
intermediate $\sigma$-flag $G'$ of some size $n'$ with $m \leq n' \leq n$ inside $G$ and
then sampling~$F$ inside $G'$.
These two procedures describe the same sampling experiment.  The resulting
identity is the single-flag case of Razborov's chain
rule~\cite[Lemma~2.2]{razborov2007flag}.

\begin{lemma}[Chain rule]\label{lem:chain-rule}
  For $F \in \mathcal{F}^\sigma_m$, $G \in \mathcal{F}^\sigma_n$, and $n'$
  with $m \leq n' \leq n$,
  \[
    \den{F}{G}
    \;=\;
    \sum_{G' \in \mathcal{F}^\sigma_{n'}} \den{F}{G'}\,\den{G'}{G}.
  \]
\end{lemma}

\paragraph{Multi-flag density.}
For pattern $\sigma$-flags $F_1, \ldots, F_t$ of sizes $m_1, \ldots, m_t$ and
a host $\sigma$-flag $G = (M, \theta_G)$ of size $n$, we say $F_1, \ldots, F_t$
\emph{fit} in $G$ if
\[
  n - k \;\geq\; (m_1 - k) + \cdots + (m_t - k).
\]
When $F_1, \ldots, F_t$ fit in $G$, let $(S_1, \ldots, S_t)$ be a uniformly
random tuple of pairwise-disjoint subsets $S_i \subseteq V(G) \setminus
\operatorname{Im}(\theta_G)$ with $|S_i| = m_i - k$.
The \emph{density of $F_1, \ldots, F_t$ in $G$}, written $\den{F_1,\ldots,F_t}{G}$, is the
probability that $G[S_i \cup \operatorname{Im}(\theta_G)] \cong_\sigma F_i$
for every $i \in [t]$.
When $F_1, \ldots, F_t$ do not fit in $G$, we set
$\den{F_1,\ldots,F_t}{G} := 0$ by convention.

\subsection{The Flag Algebra}

Fix a flag type $\sigma$ of size $k$.
Informally, the flag algebra $\mathcal{A}^\sigma$ is an algebra\footnote{Here
``algebra'' means a vector space whose elements can also be multiplied.  In
this paper, the scalar coefficients are real numbers, multiplication is
bilinear and associative, there is a multiplicative unit, and the order of
multiplication does not matter: $fg=gf$.} of density expressions generated
by $\sigma$-flags.  Each $\sigma$-flag is a symbolic building block: it stands for the
density of that local labeled pattern in a large host graph.  The quotient
below records the identity of Lemma~\ref{lem:chain-rule} among these building
blocks.
Formally, let
$\mathbb{R}[\mathcal{F}^\sigma]$ be the free real vector space on
$\mathcal{F}^\sigma$, that is, the vector space of finite formal real linear
combinations of $\sigma$-flags.  For a $\sigma$-flag $F$, write $\mathbf e_F$
for its corresponding basis vector.
Define the \emph{zero space} $\mathcal{Z}^\sigma$ as the subspace generated
by all elements of the form
\begin{equation}
  \label{eq:zero-element}
  \mathbf e_F - \sum_{G \in \mathcal{F}^\sigma_n} \den{F}{G} \cdot \mathbf e_G,
  \qquad \text{for}\ F \in \mathcal{F}^\sigma_m \ \text{and}\ n \geq m.
\end{equation}
The \emph{flag algebra} is the quotient vector space
\[
  \mathcal{A}^\sigma \;:=\; \mathbb{R}[\mathcal{F}^\sigma] \,/\, \mathcal{Z}^\sigma.
\]
We write $[F]$ for the class of $\mathbf e_F$ in $\mathcal{A}^\sigma$.
Throughout the paper, capital letters $F,G$ denote finite flags, while
lowercase letters $f,g,h$ denote arbitrary flag-algebra elements.
For familiar named or pictorial flags, such as $K_2$ or
$\fKtwobull$, we retain the standard shorthand of omitting the brackets in
algebraic formulas; there the surrounding algebraic context makes the
canonical image unambiguous.

We call $\sum_{G \in \mathcal{F}^\sigma_n} \den{F}{G} \cdot \mathbf e_G$ the
\emph{expansion of $F$ at size $n$}.
By Lemma~\ref{lem:chain-rule}, a $\sigma$-flag $F \in \mathcal{F}^\sigma_m$ and
its expansion at any size $n \geq m$ have the same density in every host
$\sigma$-flag of size $\geq n$.
The quotient by $\mathcal{Z}^\sigma$ records this expansion identity as the
equality
\begin{equation}
  \label{eq:expansion-identity}
  [F]
  =
  \sum_{G \in \mathcal{F}^\sigma_n} \den{F}{G}\,[G]
  \qquad\text{in }\mathcal{A}^\sigma .
\end{equation}
We may therefore freely replace $[F]$ by its expansion at any size $n \geq m$.

\begin{example}[$K_2$ expanded at size $3$ in the triangle-free algebra]
  Working in the triangle-free flag algebra (i.e., $\mathcal{H} = \{K_3\}$),
  the $\emptyset$-flags of size $3$ are the three triangle-free graphs on three vertices:
  $\fbarKthree$, $\fbarPthree$, and $\fPthree$.
  The density of an edge in each of them is
  \[
    \den{K_2}{\fbarKthree}=0,\quad
    \den{K_2}{\fbarPthree}=\frac13,\quad
    \den{K_2}{\fPthree}=\frac23,
  \]
  so the zero-space quotient identifies
  \begin{equation}
    K_2
    \;=\;
    \frac13\,\fbarPthree
    + \frac23\,\fPthree
    \qquad\text{in } \mathcal{A}^{\emptyset}.
    \label{eq:k2-expansion}
  \end{equation}
\end{example}

\paragraph{Multiplication.}
For $\sigma$-flags $F_1 \in \mathcal{F}^\sigma_{m_1}$ and
$F_2 \in \mathcal{F}^\sigma_{m_2}$, choose an auxiliary size
$\ell \geq m_1 + m_2 - k$ and define
\begin{equation}
  \label{eq:mul}
  [F_1] \cdot [F_2]
  \;:=\;
  \sum_{G \in \mathcal{F}^\sigma_\ell} \den{F_1,F_2}{G} \cdot [G]
  \qquad\text{in } \mathcal{A}^\sigma.
\end{equation}
At first glance, this definition appears to depend on both the choice of $\ell$
and the choice of representatives for~$[F_1]$ and $[F_2]$.  The zero-space
quotient ensures neither matters: the product is independent of $\ell$ and
respects isomorphism classes of the input $\sigma$-flags.  Extending bilinearly
to all of $\mathcal{A}^\sigma$ gives an algebra over the real numbers.  Its
multiplication is commutative, meaning that $fg=gf$ for all
$f,g\in\mathcal{A}^\sigma$.
The unit of this algebra is the basis element represented by the $\sigma$-flag
$(\sigma, \operatorname{id}_{[k]})$ of size~$k$, in which the underlying graph
is $\sigma$ itself and every vertex is labeled; we write it as
$\mathbf{1}_\sigma$, and simply as $\mathbf{1}$ when
$\sigma=\emptyset$.

\begin{example}[Multiplication of $K_2^\bullet$ and $\overline{K_2}^\bullet$]
\label{ex:typed-edge-nonedge-product}
  Let $\sigma$ be the one-vertex flag type.  Taking $\ell = 3$, the density
  $\den{\fKtwobull,\,\fbarKtwobull}{G}$ equals $\frac{1}{2}$ for each size-3
  $\sigma$-flag $G$ whose labeled vertex has degree exactly $1$:
  \[
    \den{\fKtwobull,\,\fbarKtwobull}{\fbarPthreebull}
    \;=\;
    \den{\fKtwobull,\,\fbarKtwobull}{\fPthreebull}
    \;=\; \frac{1}{2},
  \]
  and vanishes on all other size-3 $\sigma$-flags:
  \[
    \den{\fKtwobull,\,\fbarKtwobull}{G} = 0
    \qquad\text{for all } G \in \mathcal{F}_3^\sigma \setminus \{\fbarPthreebull,\,\fPthreebull\}.
  \]
  Hence
  \[
    \fKtwobull \cdot \fbarKtwobull
    \;=\;
    \frac{1}{2}\,\fbarPthreebull + \frac{1}{2}\,\fPthreebull
    \qquad\text{in } \mathcal{A}^\sigma.
  \]
\end{example}

\subsection{Positive Homomorphisms and Semantic Order}
\label{sec:background-semantic-order}

Fix a flag type $\sigma$.  
The flag algebra $\mathcal{A}^\sigma$ provides algebraic syntax for expressing
densities of pattern $\sigma$-flags.  Flag algebra proofs, however, target
\emph{asymptotic} statements: rather than a single finite host $\sigma$-flag, we
consider increasing sequences~$\{G_n\}_{n\in\mathbb{N}}$ of host
$\sigma$-flags with $|V(G_n)|\to\infty$ and ask what happens to $\sigma$-flag
densities in the limit.  Positive homomorphisms can be thought of as the
limiting density profiles of convergent sequences of host $\sigma$-flags.
The semantic non-negativity cone
$\mathcal{C}^\sigma$ then consists of the flag-algebra elements
whose value under every positive homomorphism is non-negative, and establishing
membership in $\mathcal{C}^\sigma$ is precisely what flag algebra proofs aim
to do.

\paragraph{Convergent sequences.}
Call a sequence $\{G_n\}_{n\in\mathbb{N}}$ of host $\sigma$-flags \emph{increasing} if
$|V(G_1)| < |V(G_2)| < \cdots$.  Each host $\sigma$-flag $G$ defines a point
$p^G \in [0,1]^{\mathcal{F}^\sigma}$ by $p^G(F) := \den{F}{G}$.
An increasing sequence is \emph{convergent} if the sequence of points $p^{G_n}$
converges in $[0,1]^{\mathcal{F}^\sigma}$ endowed with the product topology;
equivalently, if $\lim_{n\to\infty}\den{F}{G_n}$ exists for every $\sigma$-flag $F$.

\paragraph{Positive homomorphisms.}
A \emph{positive homomorphism} on
$\mathcal{A}^\sigma$ is a map
$\phi:\mathcal{A}^\sigma \to \mathbb{R}$ that preserves addition,
multiplication, multiplication by real scalars, and the unit.  It must also
satisfy $\phi([F]) \geq 0$ for every $\sigma$-flag $F$.
Following standard flag-algebra terminology, we call such a homomorphism
\emph{positive}, although its defining inequalities are non-strict.
We write $\poshom{\sigma}$ for the set of all positive homomorphisms on
$\mathcal{A}^\sigma$.

In $\mathcal{A}^\sigma$, the following identity holds for every $\ell \geq |V(\sigma)|$:
\begin{equation}
  \sum_{F \in \mathcal{F}^\sigma_\ell} [F] \;=\; \mathbf{1}_\sigma.
  \label{eq:sum-to-one}
\end{equation}
Applying any $\phi \in \poshom{\sigma}$ to both sides gives
$\sum_{F \in \mathcal{F}^\sigma_\ell} \phi([F]) = 1$;
non-negativity then forces $\phi([F]) \in [0,1]$ for every $\sigma$-flag $F$.
Thus, each $\phi \in \poshom{\sigma}$ determines a point in $[0,1]^{\mathcal{F}^\sigma}$
via the assignment $F \mapsto \phi([F])$, so $\poshom{\sigma}$ can be identified with
a subset of $[0,1]^{\mathcal{F}^\sigma}$.  %

The following theorem makes precise the sense in which positive homomorphisms
are exactly the limiting density profiles of convergent sequences of
$\sigma$-flags~\cite{razborov2007flag}.

\begin{theorem}[Razborov]\label{thm:convergent-hom}
  \begin{enumerate}
  \item[\textup{(a)}] If $\{G_n\}_{n\in\mathbb{N}}$ is a convergent sequence of
    $\sigma$-flags, then
    \[
      \lim_{n\to\infty} p^{G_n} \in \poshom{\sigma}.
    \]
  \item[\textup{(b)}] Conversely, every $\phi \in \poshom{\sigma}$ satisfies
    $\phi = \lim_{n\to\infty} p^{G_n}$ for some convergent sequence
    $\{G_n\}_{n\in\mathbb{N}}$ of $\sigma$-flags.
  \end{enumerate}
\end{theorem}

\paragraph{Semantic non-negativity cone and order.}
For a fixed flag type $\sigma$, define the \emph{semantic non-negativity cone}
\[
  \mathcal{C}^\sigma
  :=
  \{f \in \mathcal{A}^\sigma
    \mid \phi(f) \geq 0
    \text{ for every } \phi \in \poshom{\sigma}\}.
\]
We write
\[
  f \leq_{\mathcal{H},\sigma} g
  \qquad\text{when}\qquad
  g-f \in \mathcal{C}^\sigma,
\]
where the subscript $\mathcal{H}$ records that the comparison is made in the
\(\mathcal{H}\)-free world: both $\mathcal{A}^\sigma$ and $\mathcal{C}^\sigma$ are
built from \(\mathcal{H}\)-free graphs.
Equivalently, $f\leq_{\mathcal{H},\sigma} g$ means that
$\phi(f)\leq\phi(g)$ for every $\phi \in \poshom{\sigma}$.  When $\sigma=\emptyset$, we write $f \leq_{\mathcal{H}} g$ in place of $f \leq_{\mathcal{H},\emptyset} g$, and $\mathcal{C}$ in place of $\mathcal{C}^\emptyset$.
Note that $f \in \mathcal{C}^\sigma$ is the same as saying $0 \leq_{\mathcal{H},\sigma} f$.

Two immediate examples of elements of $\mathcal{C}^\sigma$ are flag basis
elements and squares.  Every $\sigma$-flag $F$ gives an element
$[F]\in\mathcal{C}^\sigma$, because positive homomorphisms send every flag
basis element to a non-negative real by definition.
For any $f \in \mathcal{A}^\sigma$, the square $f^2$ satisfies $0 \leq_{\mathcal{H},\sigma} f^2$,
because $\phi(f^2) = \phi(f)^2 \geq 0$ for every~$\phi \in \poshom{\sigma}$.
Moreover, if $0 \leq_{\mathcal{H},\sigma} f$ and $0 \leq_{\mathcal{H},\sigma} g$, then $0 \leq_{\mathcal{H},\sigma} f + g$,
since $\phi(f+g) = \phi(f) + \phi(g) \geq 0$ for every $\phi \in \poshom{\sigma}$.
Similarly, if $0 \leq_{\mathcal{H},\sigma} f$ and $c$ is a non-negative real number, then $0 \leq_{\mathcal{H},\sigma} c \cdot f$.

\paragraph{From semantic non-negativity to Tur\'an density.}
The problem of bounding the Tur\'an density $\tdensity{F}{\mathcal{H}}$ from above
reduces to establishing membership in $\mathcal{C}^\emptyset$, as the following
lemma shows.

\begin{lemma}\label{lem:semantic-bound}
  If $[F] \leq_{\mathcal{H}} c \cdot \mathbf{1}$ for some $c \in \mathbb{R}$, then
  $\tdensity{F}{\mathcal{H}} \leq c$.
\end{lemma}
\begin{proof}
  Suppose for contradiction that $\tdensity{F}{\mathcal{H}} > c$.  Then
  there exist $\varepsilon > 0$ and an increasing sequence of
  \(\mathcal{H}\)-free graphs $\{G_k\}_{k\in\mathbb{N}}$ such that
  $\den{F}{G_k} > c + \varepsilon$ for all $k$.  Since each $p^{G_k}$
  is a point in $[0,1]^{\mathcal{F}^\emptyset}$ and
  $\mathcal{F}^\emptyset$ is countable, this product space is compact and
  metrizable.  Hence the sequence $\{p^{G_k}\}$ has a convergent
  subsequence $\{p^{G_{k_j}}\}$; by \Cref{thm:convergent-hom}(a) its limit is some
  $\phi \in \poshom{\emptyset}$.  Since $\phi$ is the pointwise limit of
  $\{p^{G_{k_j}}\}$, we have
  $\phi([F]) = \lim_j \den{F}{G_{k_j}} \geq c + \varepsilon > c$.  But
  $[F] \leq_{\mathcal{H}} c \cdot \mathbf{1}$ means $\phi([F]) \leq c$ for every~$\phi \in \poshom{\emptyset}$, a contradiction.
\end{proof}

A matching lower bound typically comes from an explicit \(\mathcal{H}\)-free construction
whose $F$-density approaches~$c$, pinning down $\tdensity{F}{\mathcal{H}}$ exactly.
Two classical results illustrate how this looks in flag algebra language.

\begin{example}[Mantel's theorem~\cite{mantel}]\label{ex:mantel}
  Mantel's theorem states $\tdensity{K_2}{K_3} = \tfrac{1}{2}$, meaning
  asymptotically at most half of all vertex pairs in a triangle-free graph
  are edges.  The upper bound translates into the flag algebra inequality
  $[K_2] \leq_{K_3} \tfrac{1}{2}\cdot\mathbf{1}$; the balanced complete bipartite
  graphs $K_{\lfloor n/2\rfloor,\lceil n/2\rceil}$ witness the lower bound.
\end{example}

\begin{example}[Erd\H{o}s Pentagon Theorem~{\cite{grzesik2012,hatami2012}}]\label{ex:pentagon}
  The Erd\H{o}s pentagon problem asks for the maximum density of $C_5$ in
  triangle-free graphs; the asymptotic answer is $\tdensity{C_5}{K_3} = \tfrac{24}{625}$.  The upper bound corresponds to the flag algebra inequality
  $[C_5] \leq_{K_3} \tfrac{24}{625}\cdot\mathbf{1}$; an explicit \(K_3\)-free construction
  (the balanced blow-up of $C_5$) witnesses the lower bound.
\end{example}

\subsection{The Downward Operator}
\label{sec:background-downward}
Labeled flag algebras (i.e., $\mathcal{A}^\sigma$ with $\sigma \neq \emptyset$) are often
richer than the unlabeled algebra $\mathcal{A}^\emptyset$: establishing membership
in $\mathcal{C}^\sigma$ can be easier than working directly in $\mathcal{C}$.
The final extremal statement, however, is always about unlabeled graphs.
The downward operator bridges this gap by mapping $\mathcal{A}^\sigma$ to
$\mathcal{A}^\emptyset$ in a way that preserves semantic non-negativity
(\Cref{thm:downward-nonneg}), so valid inequalities proved in a labeled algebra
can be transferred to the unlabeled one.

For a $\sigma$-flag $F = (M, \theta_F)$, let $F|_\emptyset$ denote the
$\emptyset$-flag obtained from $F$ by forgetting the embedding $\theta_F$,
i.e., viewing $M$ as a $\emptyset$-flag.
For a $\sigma$-flag $F$ of size $m$ with $|V(\sigma)|=k$, let
\[
  q_\sigma(F)
  \;:=\;
  \frac{
    \bigl|\{\theta' : [k]\hookrightarrow V(M)
       \mid \theta' \text{ embeds }\sigma\text{ into }M
       \text{ and } (M,\theta')\cong_\sigma F\}\bigr|}
       {m(m-1)\cdots(m-k+1)}
\]
be the probability that a uniformly random injection $[k]\hookrightarrow V(M)$
recovers the same $\sigma$-flag as $F$.
The \emph{downward operator}
$\llbracket\cdot\rrbracket_\sigma:\mathcal{A}^\sigma\to\mathcal{A}^\emptyset$
is then defined on flag basis elements by
\[
  \llbracket [F] \rrbracket_\sigma
  \;:=\;
  q_\sigma(F)\cdot [F|_\emptyset],
\]
extended linearly to all of $\mathcal{A}^\sigma$.
For $f\in\mathcal{A}^\sigma$, we call
$\llbracket f\rrbracket_\sigma$ the \emph{downward image} of $f$.

We write $\langle\sigma\rangle_0\in\mathcal{A}^\emptyset$ for the empty-type
flag-algebra basis element represented by the underlying graph of the flag type
$\sigma$, with its vertex labels forgotten.  Applying the downward operator to
the typed unit therefore gives
\[
  \llbracket \mathbf{1}_\sigma \rrbracket_\sigma
  =
  q_\sigma(\mathbf{1}_\sigma)\,\langle\sigma\rangle_0.
\]

\paragraph{Random homomorphisms.}
A positive homomorphism $\phi_0 \in \poshom{\emptyset}$ records the limiting
densities of unlabeled graphs along a convergent sequence of finite graphs.
For a nonempty flag type $\sigma$, its random extension can be understood
through any such sequence realizing $\phi_0$: in each finite graph, choose
uniformly at random an embedding of $\sigma$ and use the embedded vertices as
labels.  Razborov's random-extension theorem describes the limiting
distribution of the resulting $\sigma$-flag densities.
Whenever $\phi_0(\langle\sigma\rangle_0) > 0$, meaning that the underlying
unlabeled graph of $\sigma$ has positive limiting density, the following
holds~\cite{razborov2007flag}:
there exists a unique probability measure on $\poshom{\sigma}$,
which we denote $\Ext_\sigma(\phi_0)$,
such that for every~$f \in \mathcal{A}^\sigma$, the random positive
homomorphism $\boldsymbol{\phi}^\sigma \sim \Ext_\sigma(\phi_0)$ satisfies
\begin{equation}\label{eq:ensemble}
  \mathbb{E}\bigl[\boldsymbol{\phi}^\sigma(f)\bigr]
  \;=\;
  \frac{\phi_0\bigl(\llbracket f \rrbracket_\sigma\bigr)}
       {\phi_0\bigl(\llbracket \mathbf{1}_\sigma \rrbracket_\sigma\bigr)}.
\end{equation}
Note that the denominator is positive by the assumption
$\phi_0(\langle\sigma\rangle_0)>0$, since
$\phi_0\bigl(\llbracket \mathbf{1}_\sigma \rrbracket_\sigma\bigr)
=q_\sigma(\mathbf{1}_\sigma)\cdot
\phi_0(\langle\sigma\rangle_0)$ and $q_\sigma(\mathbf{1}_\sigma)>0$ always.

Equation~\eqref{eq:ensemble} implies that the downward operator $\llbracket \cdot \rrbracket_\sigma$ maps $\mathcal{C}^\sigma$ to $\mathcal{C}$:

\begin{theorem}[Razborov~{\cite{razborov2007flag}}]\label{thm:downward-nonneg}
  For any flag type $\sigma$ and any $f \in \mathcal{A}^\sigma$ with $0 \leq_{\mathcal{H},\sigma} f$,
  \[
    0 \;\leq_{\mathcal{H}}\; \llbracket f \rrbracket_\sigma.
  \]
\end{theorem}
\begin{proof}[Proof sketch]
  Let $\phi_0 \in \poshom{\emptyset}$; we show $\phi_0(\llbracket f \rrbracket_\sigma) \geq 0$.
  If $\phi_0(\langle\sigma\rangle_0) = 0$, then $\phi_0$ assigns density zero
  to the underlying unlabeled graph of $\sigma$.  The underlying graph of
  every $\sigma$-flag contains an induced copy of $\sigma$, so the downward
  image of every $\sigma$-flag basis element also evaluates to zero under
  $\phi_0$.  By linearity, $\phi_0(\llbracket f \rrbracket_\sigma) = 0$.
  Otherwise, let $\boldsymbol{\phi}^\sigma \sim \Ext_\sigma(\phi_0)$.
  Since $0 \leq_{\mathcal{H},\sigma} f$, every positive homomorphism on
  $\mathcal{A}^\sigma$ evaluates $f$ non-negatively, so
  $\mathbb{E}[\boldsymbol{\phi}^\sigma(f)] \geq 0$.
  Rearranging Equation~\eqref{eq:ensemble}:
  \[
    \phi_0\bigl(\llbracket f \rrbracket_\sigma\bigr)
    \;=\;
    \phi_0\bigl(\llbracket \mathbf{1}_\sigma \rrbracket_\sigma\bigr)
    \cdot
    \mathbb{E}\bigl[\boldsymbol{\phi}^\sigma(f)\bigr]
    \;\geq\; 0.
  \]
  Since $\phi_0$ was arbitrary, we have $0 \leq_{\mathcal{H}} \llbracket f \rrbracket_\sigma$.
\end{proof}

A flag algebra proof of $[F] \leq_{\mathcal{H}} c\cdot\mathbf{1}$ typically proceeds as follows:
first, one exhibits $f \in \mathcal{A}^\sigma$ with $0 \leq_{\mathcal{H},\sigma} f$ for some
nonempty flag type $\sigma$ (e.g., a square $f = g^2$) and applies
\Cref{thm:downward-nonneg} to obtain $0 \leq_{\mathcal{H}} \llbracket f \rrbracket_\sigma$;
then, one adds $\llbracket f \rrbracket_\sigma$ to $[F]$ or an expansion of $[F]$
at an appropriate size, and uses non-negativity of flag basis elements and
Equation~\eqref{eq:sum-to-one} to conclude $[F] \leq_{\mathcal{H}} c\cdot\mathbf{1}$.
The following example illustrates this pattern.

\begin{example}[Mantel's theorem via flag algebras]\label{ex:mantel-fa}
We prove $K_2 \leq_{K_3} \tfrac{1}{2}\cdot\mathbf{1}$ in the triangle-free flag algebra.
Let $\sigma$ be the one-vertex flag type.

\textit{Step 1: A key inequality.}
Since $0 \leq_{\mathcal{H},\sigma} ((\fKtwobull) - (\fbarKtwobull))^2$,
\Cref{thm:downward-nonneg} gives $0 \leq_{K_3} \llbracket((\fKtwobull) - (\fbarKtwobull))^2\rrbracket_\sigma$.
Expanding the square, computing each product at size~$3$, and
then applying the downward operator
gives the following inequality in $\mathcal{A}^{\emptyset}$:
\begin{align}
  0 \leq_{K_3} \llbracket((\fKtwobull) - (\fbarKtwobull))^2\rrbracket_\sigma\notag
  & = \llbracket(\fKtwobull)\cdot(\fKtwobull) - 2\cdot(\fKtwobull)\cdot(\fbarKtwobull) + (\fbarKtwobull)\cdot(\fbarKtwobull)\rrbracket_\sigma\notag\\
  &= \left\llbracket\fbarKthreebull + \fEthreebull
    - \fbarPthreebull - \fPthreebull
    + \fKonetwobull\right\rrbracket_\sigma\notag\\
  &= \fbarKthree - \frac{1}{3}\,\fbarPthree - \frac{1}{3}\,\fPthree.
  \label{eq:mantel-key}
\end{align}

\textit{Step 2: Algebraic derivation.}
Starting from Equation~\eqref{eq:k2-expansion}, we add two non-negative terms and simplify using Equation~\eqref{eq:sum-to-one}:
\begin{align*}
  K_2
  &= \frac{1}{3}\,\fbarPthree + \frac{2}{3}\,\fPthree
  &&\text{by Equation~\eqref{eq:k2-expansion}}\\
  &\leq_{K_3} \frac{1}{3}\,\fbarPthree + \frac{2}{3}\,\fPthree +
    \frac{1}{2}\left(\fbarKthree - \frac{1}{3}\,\fbarPthree - \frac{1}{3}\,\fPthree\right) +
    \frac{1}{3}\,\fbarPthree
  &&\text{by Equation~\eqref{eq:mantel-key} and $0 \leq_{K_3} \fbarPthree$}\\
  &= \frac{1}{2}\cdot\left(\fbarKthree + \fbarPthree + \fPthree\right)
  \;=\; \frac{1}{2}\cdot\mathbf{1}.
  &&\text{by Equation~\eqref{eq:sum-to-one}}
\end{align*}
\end{example}

\section{Formalizing Flag Algebras in Lean~4}
\label{sec:abstract}

Informal flag algebra proofs are compact partly because their notation is
deliberately overloaded.  The same symbol may denote a concrete labeled graph,
its isomorphism class, a basis vector in a free vector space, an element of the
quotient algebra, or the real number obtained by applying a positive
homomorphism to that element.  On paper, these changes of viewpoint are usually
harmless and are left implicit.  In Lean, however, they create an immediate
formalization challenge: each role must be represented by a separate type, and
each passage between roles must be justified.  In particular, constructions
on concrete graphs must be invariant under isomorphism, expansion identities
must hold in the quotient, multiplication must be well defined on the quotient
algebra, and positive homomorphisms must preserve the algebraic operations.

Our first design choice is therefore to make the Lean definitions follow the
mathematics of \Cref{sec:background} as closely as possible, even when the
resulting definitions cannot be executed directly.  Accordingly, flags are
isomorphism classes of labeled graphs, densities are finite sampling
probabilities, the flag algebra is a quotient by expansion identities, and
positive homomorphisms provide semantic interpretations of algebra elements.
The following subsections present these definitions in turn.  Together with
the maps between these representations and the theorems establishing their
correctness, they form the \emph{specification layer}.  Several of these
quotient-based definitions involve nonconstructive choices and are therefore
marked \lean{noncomputable}.  For the finite calculations needed in later
proofs, such as testing flag isomorphism and computing densities and products,
the \emph{reflection layer} of \Cref{sec:reflection} provides executable
procedures corresponding to the abstract definitions in the specification
layer.  Its adequacy theorems prove that the results returned by these
procedures agree with those abstract definitions.

Our second design choice concerns how forbidden-subgraph assumptions are
represented in the Lean development.  In the mathematical presentation of
\Cref{sec:background}, we first fix \(\mathcal{H}\) and construct flags and
flag algebras within the class of \(\mathcal{H}\)-free graphs.  Our Lean
formalization instead constructs the ambient flag algebras without a
forbidden-family parameter and imposes \(\mathcal{H}\)-freeness later by
restricting the positive homomorphisms considered by the semantic order.  Thus,
the core formalization implements the unconstrained theory.  If an arbitrary
forbidden family were carried through the definitions from the outset, even
the elementary constructions of flags, densities, expansion identities, and
multiplication would require an additional parameter and corresponding side
conditions.  By postponing the constraint, we formalize the
\(\mathcal{H}\)-independent algebraic definitions and results only once.  The
same Lean infrastructure can then be reused for different forbidden graphs,
with the appropriate constraint introduced only in the semantic comparisons
required by each extremal application.  This design is sound for the
Tur\'an-density applications formalized here: \Cref{sec:forbidden} defines the
ensemble semantic order and states the transfer theorem that converts an
empty-type ensemble inequality into the corresponding extremal-density bound.
\Cref{sec:metatheory} places this design in a broader mathematical setting by
comparing the ensemble semantic order with the order obtained by building the
constraint into the flag algebra from the outset.  It gives a criterion for
the two orders to agree, broad sufficient conditions for agreement, and an
example in which they differ.  A forthcoming paper will develop this
meta-theory and its applications in full.

For Lean-facing notation, we write the empty flag type as \(\emptyt\), rather
than the \(\emptyset\) used in \Cref{sec:background}, to distinguish it from
ordinary empty sets.

The Lean snippets below are lightly adapted for exposition.  We make otherwise
inferred parameters explicit and occasionally replace implementation-oriented
names with clearer paper-facing names, so that the relevant types and
dependencies are visible without requiring the surrounding source context.

\subsection{Flag Types and Flags}
\label{sec:flagtypes}

Our formalization builds on Mathlib, particularly its graph library.  Before
giving the Lean definitions of flags, we recall the small amount of Mathlib
notation used below.  In \Cref{sec:background}, we described a finite graph as
having a vertex set $V(G)$.  In Lean, the vertex set is represented by a type:
a term of type \lean{SimpleGraph V} is a simple graph whose vertices are the
elements of the type \lean{V}.  Finiteness is supplied separately by
assumptions such as \lean{Fintype V}.  For example,
\lean{SimpleGraph (Fin n)} is a graph on the type \lean{Fin n}, whose elements
are the natural numbers \(0,\ldots,n-1\).

Mathlib provides ${G\hookrightarrow_g H}$ for the type of graph embeddings
from $G$ to $H$.  Such an embedding is an injective vertex map that
preserves and reflects adjacency, so its image is an induced copy of $G$
in $H$.  Mathlib also provides $G\simeq_g H$ for graph isomorphisms.
With this notation in place, the concrete pair $F=(M,\theta_F)$ from
\Cref{sec:background} is represented by \lean{LabeledGraph}, while its
isomorphism class is represented by \lean{Flag}.

\begin{lstlisting}
abbrev FlagType (T : Type) := SimpleGraph T

structure LabeledGraph {T : Type} [Fintype T] ((*@$\sigma$@*) : FlagType T) (V : Type) where
  graph      : SimpleGraph V
  type_embed : (*@$\sigma$@*) (*@$\hookrightarrow$@*)g graph
\end{lstlisting}
Here \lean{FlagType T} is simply a name for \lean{SimpleGraph T}: a flag type
is represented by a simple graph on the label type \lean{T}.  The structure
\lean{LabeledGraph} packages the two pieces of a concrete $\sigma$-flag: the
field \lean{graph} stores the underlying graph $M$, and the field
\lean{type\_embed} stores the embedding
$\theta_F:\sigma\hookrightarrow_g M$.  In the signature, braces mark
implicit parameters that Lean normally infers from context, whereas square
brackets mark typeclass assumptions.  Thus, \lean{[Fintype T]} tells Lean that
the label type \lean{T} is finite; later definitions use assumptions such as
\lean{[DecidableEq V]} when equality on a vertex type must be decidable.

We define \lean{LabeledGraphIso} as a structure representing isomorphisms
between two \lean{LabeledGraph} terms, and define \lean{Flag} by quotienting
\lean{LabeledGraph} by this isomorphism relation:
\begin{lstlisting}
structure LabeledGraphIso {T V V' : Type} [Fintype T]
    {(*@$\sigma$@*) : FlagType T} (G : LabeledGraph (*@$\sigma$@*) V) (G' : LabeledGraph (*@$\sigma$@*) V') where
  graph_iso     : G.graph (*@$\simeq$@*)g G'.graph
  type_preserve : graph_iso (*@$\circ$@*) G.type_embed = G'.type_embed

infixl:50 " (*@\textcolor{orange!80!black}{$\simeq$}@*)f " => LabeledGraphIso

instance labeledGraphSetoid {T : Type} [Fintype T] ((*@$\sigma$@*) : FlagType T) (V : Type)
    : Setoid (LabeledGraph (*@$\sigma$@*) V) where
  -- r : LabeledGraph (*@\textcolor{gray}{$\sigma$}@*) V (*@\textcolor{gray}{$\to$}@*) LabeledGraph (*@\textcolor{gray}{$\sigma$}@*) V (*@\textcolor{gray}{$\to$}@*) Prop
  r     := fun G G' => Nonempty (G (*@$\simeq$@*)f G')
  iseqv := ... -- proof that r defines an equivalence relation

def Flag {T : Type} [Fintype T] ((*@$\sigma$@*) : FlagType T) (V : Type) : Type :=
  Quotient (labeledGraphSetoid (*@$\sigma$@*) V)
\end{lstlisting}
The structure \lean{LabeledGraphIso} represents an isomorphism of concrete
$\sigma$-flags.
It consists of a graph isomorphism \lean{graph\_iso} between the underlying
graphs, together with a compatibility proof \lean{type\_preserve} saying that
this isomorphism sends the labeled copy of $\sigma$ in $G$ to the labeled copy
of $\sigma$ in $G'$.  This is the Lean form of the condition
$\alpha\circ\theta_G=\theta_{G'}$.  The notation $G\simeq_f G'$ denotes this
type of isomorphism data, with $\sigma$ inferred from the two labeled graphs.

Lean's quotient type does not accept an arbitrary relation directly; it
requires a \emph{setoid}, which packages a relation with proofs that it is an
equivalence relation.  The instance \lean{labeledGraphSetoid} provides this package for
\lean{LabeledGraph}~$\sigma$~\lean{V}.  Its field \lean{r} identifies two
terms $G,G'$ when a labeled-graph isomorphism $G\simeq_f G'$ exists, and its
field \lean{iseqv} contains the proofs that this relation is reflexive,
symmetric, and transitive.  Finally,
\lean{Flag}~$\sigma$~\lean{V} is defined as the quotient of
\lean{LabeledGraph}~$\sigma$~\lean{V} by that setoid, so its elements are
precisely isomorphism classes of concrete $\sigma$-flags on the vertex type
\lean{V}.

Lean's \lean{Quotient} type comes with double-bracket notation for equivalence
classes.  In
Lean code, the unsubscripted notation \(\llbracket x\rrbracket\)
means ``the quotient class of the representative \(x\),'' with the relevant
quotient inferred from the expected type.  Thus, in this subsection, when the
expected type is \lean{Flag}~\(\sigma\)~\lean{V},
\(\llbracket G\rrbracket\) is the \(\sigma\)-flag represented by the concrete
labeled graph \(G\).  Throughout the paper, unsubscripted double brackets denote
quotient classes, whereas the subscripted brackets
\(\llbracket-\rrbracket_\sigma\) and \(\llbracket-\rrbracket_0\) denote the
downward operator.  We return to this quotient-class notation in
\Cref{sec:formal-flag-algebra}.

We can now represent $\mathcal{F}^\sigma_n$ in Lean by
\lean{Flag}~$\sigma$~\lean{(Fin}~$n$\lean{)}.  The type
\lean{FinFlag}~$\sigma$ collects $\sigma$-flags of all finite sizes and
represents $\mathcal{F}^\sigma = \bigcup_n \mathcal{F}^\sigma_n$.  It is
defined as a dependent sum type: each element is a dependent pair
$\hat F = \langle n, F\rangle$, where $n$ records the size and $F$ is the
$\sigma$-flag of that size.  In Lean, the two components are accessed as
$\hat F.1$ and $\hat F.2$.
\begin{lstlisting}
def FinFlag {k : (*@$\mathbb{N}$@*)} ((*@$\sigma$@*) : FlagType (Fin k)) : Type :=
  (*@$\Sigma$@*) (n : (*@$\mathbb{N}$@*)), Flag (*@$\sigma$@*) (Fin n)
\end{lstlisting}

Unlike the definitions in \Cref{sec:background}, none of the definitions above
carries the \(\mathcal{H}\)-free condition: the underlying graphs are
arbitrary.  This is a deliberate interface choice.  The core quotient and
density machinery is developed once for all finite graphs, while
\(\mathcal{H}\)-freeness is imposed later by restricting the positive
homomorphisms used in the ensemble semantic order
(\Cref{sec:forbidden}).

\subsection{Subflag Densities}
\label{sec:formal-subflag-densities}

Formalizing the density function is not straightforward.  The density
$\den{F_1,\ldots,F_t}{G}$ from \Cref{sec:background} should be a
function taking $t$ pattern $\sigma$-flags and a host $\sigma$-flag to a value in $[0,1]$, but
\lean{Flag} is a quotient type of \lean{LabeledGraph}.  The standard way to
define a function from a quotient type in Lean is to first define it on the
underlying type, prove that it respects the equivalence relation, and lift it
via \lean{Quotient.lift}.  That is, one must first define the density with
\lean{LabeledGraph} inputs, and then prove that replacing any input by an
isomorphic \lean{LabeledGraph} gives the same result.

We first introduce \lean{LabeledGraphList} as an abbreviation for a
(\lean{Fin t})-indexed tuple of \lean{LabeledGraph}s.
Given such a tuple of pattern \lean{LabeledGraph}s and a host \lean{LabeledGraph},
\lean{labeledGraphListDensity} computes the density as the number of tuples
$(S_1, \ldots, S_t)$ satisfying the conditions in the definition of $\den{F_1,\ldots,F_t}{G}$,
divided by the total number of candidate tuples.  Rather than adopting the
probabilistic sampling formulation of \Cref{sec:background}, this direct ratio is easier to work with in Lean.
\begin{lstlisting}
abbrev LabeledGraphList {T : Type} [Fintype T]
    ((*@$\sigma$@*) : FlagType T) (t : (*@$\mathbb{N}$@*)) (Vl : Fin t (*@$\to$@*) Type) : Type :=
  (i : Fin t) (*@$\to$@*) LabeledGraph (*@$\sigma$@*) (Vl i)

noncomputable def labeledGraphListDensity
    {T W : Type} [Fintype T] [Fintype W] [DecidableEq W]
    {t : (*@$\mathbb{N}$@*)} {Vl : Fin t (*@$\to$@*) Type} [FintypeList Vl] [DecidableEqList Vl]
    {(*@$\sigma$@*) : FlagType T} (Hl : LabeledGraphList (*@$\sigma$@*) t Vl) (G : LabeledGraph (*@$\sigma$@*) W) : (*@$\mathbb{Q}$@*) :=
  let r_list (i : Fin t) := (Hl i).size - (*@$\sigma$@*).size
  labeledGraphListCount Hl G / multinomialCoefficient r_list (G.size - (*@$\sigma$@*).size)
\end{lstlisting}
The first line defines \lean{LabeledGraphList} by spelling out what a finite
list means in this dependent setting: it is a function that assigns to each
index \lean{i : Fin t} a concrete $\sigma$-flag \lean{LabeledGraph}~$\sigma$~\lean{(Vl i)}.
The family \lean{Vl : Fin t $\to$ Type} records the vertex type of each entry, so
the $t$ pattern $\sigma$-flags are not forced to share one ambient vertex type.

The second definition gives the finite ratio corresponding to
$\den{F_1,\ldots,F_t}{G}$.  Here \lean{Hl} is the tuple of pattern labeled
graphs and \lean{G} is the concrete labeled graph in which they are counted.
The local function \lean{r\_list} records, for each pattern, how many unlabeled
vertices it contributes:
\[
  r_i = |V(F_i)| - |\sigma|.
\]
The helper \lean{labeledGraphListCount}, whose Lean definition is omitted from
the displayed excerpt, supplies the numerator.  It counts tuples of induced
labeled subgraphs of \lean{G} such that the entry at each index \lean{i} is
isomorphic, as a $\sigma$-flag, to the pattern \lean{Hl i}, and different
entries use pairwise-disjoint sets of unlabeled vertices.  The denominator
applies \lean{multinomialCoefficient} to \lean{r\_list} and
\(|V(G)| - |\sigma|\); it counts all ways to choose disjoint vertex sets of
the required sizes from the unlabeled vertices of \lean{G}.  Their ratio is
the multi-flag density $\den{F_1,\ldots,F_t}{G}$, represented as a rational
number.

The square-bracketed parameters tell Lean what finite data and decision
procedures are available.  For a type \lean{X}, \lean{[Fintype X]} supplies an
enumeration of all elements of \lean{X}, while \lean{[DecidableEq X]} supplies
a procedure for deciding whether two elements are equal.  The pattern tuple
may use a different vertex type \lean{Vl i} at each index.  Accordingly,
\lean{[FintypeList Vl]} and \lean{[DecidableEqList Vl]} package the same
information for every type \lean{Vl i}.  In Lean, such square-bracketed
parameters are called \emph{typeclass parameters}, and Lean normally supplies
them automatically through its typeclass inference mechanism.

The \lean{noncomputable} keyword signals that Lean cannot synthesize an
algorithm to evaluate this definition directly.  This poses no obstacle to
formalizing flag algebra itself, but becomes relevant when applying the algebra
to concrete extremal problems, a point we return to in \Cref{sec:reflection}.

The concrete function \lean{labeledGraphListDensity} has two inputs: the tuple
\lean{Hl} of pattern labeled graphs and the host labeled graph \lean{G}.  To
make this a function of $\sigma$-flags, we lift these two inputs separately.  The host
input is lifted from \lean{LabeledGraph} to \lean{Flag}; for the pattern tuple,
the intermediate quotient is the quotient of concrete tuples by entrywise
isomorphism.  We first lift the host input \lean{G}, then lift the pattern
tuples, and finally package the result as a function on \lean{FlagList}.  The
congruence proofs, omitted below, show that the density is unchanged when the
host $\sigma$-flag or any pattern $\sigma$-flag is replaced by an isomorphic
representative.
\begin{lstlisting}
def labeledGraphListEqv {T : Type} [Fintype T] {(*@$\sigma$@*) : FlagType T}
    {t : (*@$\mathbb{N}$@*)} {Vl : Fin t (*@$\to$@*) Type} (Gl Gl' : LabeledGraphList (*@$\sigma$@*) t Vl) : Prop :=
  (*@$\forall$@*) (i : Fin t), Nonempty (Gl i (*@$\simeq$@*)f Gl' i)

instance labeledGraphListSetoid {T : Type} [Fintype T]
    ((*@$\sigma$@*) : FlagType T) (t : (*@$\mathbb{N}$@*)) (Vl : Fin t (*@$\to$@*) Type)
    : Setoid (LabeledGraphList (*@$\sigma$@*) t Vl) where
  r     := labeledGraphListEqv
  iseqv := ... -- proof that the entrywise isomorphism is an equivalence

def QuotLabeledGraphList {T : Type} [Fintype T]
    ((*@$\sigma$@*) : FlagType T) (t : (*@$\mathbb{N}$@*)) (Vl : Fin t (*@$\to$@*) Type) : Type :=
  Quotient (labeledGraphListSetoid (*@$\sigma$@*) t Vl)

abbrev FlagList {T : Type} [Fintype T]
    ((*@$\sigma$@*) : FlagType T) (t : (*@$\mathbb{N}$@*)) (Vl : Fin t (*@$\to$@*) Type) : Type :=
  (*@$\forall$@*) (i : Fin t), Flag (*@$\sigma$@*) (Vl i)

noncomputable def labeledGraphListDensityLifted
    {T W : Type} [Fintype T] [Fintype W] [DecidableEq W]
    {t : (*@$\mathbb{N}$@*)} {Vl : Fin t (*@$\to$@*) Type} [FintypeList Vl] [DecidableEqList Vl]
    {(*@$\sigma$@*) : FlagType T} (Hl : LabeledGraphList (*@$\sigma$@*) t Vl)
    : Flag (*@$\sigma$@*) W (*@$\to$@*) (*@$\mathbb{Q}$@*) :=
  Quotient.lift
    (fun G => labeledGraphListDensity Hl G)
    (...) -- proof that host-graph equivalence preserves density

noncomputable def quotLabeledGraphListDensity
    {T W : Type} [Fintype T] [Fintype W] [DecidableEq W]
    {t : (*@$\mathbb{N}$@*)} {Vl : Fin t (*@$\to$@*) Type} [FintypeList Vl] [DecidableEqList Vl]
    {(*@$\sigma$@*) : FlagType T}
    : QuotLabeledGraphList (*@$\sigma$@*) t Vl (*@$\to$@*) Flag (*@$\sigma$@*) W (*@$\to$@*) (*@$\mathbb{Q}$@*) :=
  Quotient.lift
    labeledGraphListDensityLifted
    (...) -- proof that pattern-list equivalence preserves density

noncomputable def FlagList.coe
    {T : Type} [Fintype T] {(*@$\sigma$@*) : FlagType T}
    {t : (*@$\mathbb{N}$@*)} {Vl : Fin t (*@$\to$@*) Type}
    (Fl : FlagList (*@$\sigma$@*) t Vl)
    : QuotLabeledGraphList (*@$\sigma$@*) t Vl :=
  (*@$\llbracket$@*)fun i => (Fl i).out(*@$\rrbracket$@*)

noncomputable def flagListDensity
    {T W : Type} [Fintype T] [Fintype W] [DecidableEq W]
    {t : (*@$\mathbb{N}$@*)} {Vl : Fin t (*@$\to$@*) Type} [FintypeList Vl] [DecidableEqList Vl]
    {(*@$\sigma$@*) : FlagType T}
    : FlagList (*@$\sigma$@*) t Vl (*@$\to$@*) Flag (*@$\sigma$@*) W (*@$\to$@*) (*@$\mathbb{Q}$@*) :=
  fun Fl => quotLabeledGraphListDensity (FlagList.coe Fl)
\end{lstlisting}
Here \lean{labeledGraphListEqv} says that two concrete tuples are equivalent exactly
when their entries are isomorphic as labeled graphs: for each index \lean{i},
there must exist an isomorphism \lean{Gl i}~\(\simeq_f\)~\lean{Gl' i}.  The
instance \lean{labeledGraphListSetoid} supplies the setoid required by
\lean{Quotient}.  Its quotient is \lean{QuotLabeledGraphList}.  The
abbreviation \lean{FlagList} is the more convenient tuple of already-quotiented
$\sigma$-flags.

The three density definitions then lift the original concrete density one
argument at a time.  For a fixed concrete pattern tuple \lean{Hl},
\lean{labeledGraphListDensityLifted Hl} is a function
\lean{Flag}~$\sigma$~\lean{W}~\(\to \mathbb{Q}\), obtained by quotienting the
host argument.  The next definition, \lean{quotLabeledGraphListDensity},
lifts the pattern tuple as well, so its first input is a
\lean{QuotLabeledGraphList} and its second input is a host \lean{Flag}.  In
the final definition, the input \lean{Fl} has type
\lean{FlagList}~$\sigma$~\lean{t Vl}, that is, \lean{Fl i} is already a $\sigma$-flag
for each \lean{i : Fin t}.  The term
\lean{FlagList.coe Fl} has type
\lean{QuotLabeledGraphList}~$\sigma$~\lean{t Vl}; concretely, it chooses a
representative \lean{(Fl i).out} for each $\sigma$-flag in the tuple and forms the
quotient class of the concrete tuple \lean{fun i => (Fl i).out}.  Thus,
\lean{flagListDensity} can call
\lean{quotLabeledGraphListDensity} after converting the
input tuple of $\sigma$-flags into the quotient-of-tuples input expected by the
preceding definition.

The following \lean{flagDensity}$_1$ and \lean{flagDensity}$_2$ are specializations of \lean{flagListDensity}
to singleton and pair lists:
\begin{lstlisting}
noncomputable def flagDensity(*@$_1$@*)
    {T U W : Type} [Fintype T] [Fintype U] [DecidableEq U] [Fintype W] [DecidableEq W]
    {(*@$\sigma$@*) : FlagType T} (F : Flag (*@$\sigma$@*) U) (G : Flag (*@$\sigma$@*) W) : (*@$\mathbb{Q}$@*) :=
  flagListDensity [F](*@$^\mathtt{f}$@*) G

noncomputable def flagDensity(*@$_2$@*)
    {T U(*@$_1$@*) U(*@$_2$@*) W : Type} [Fintype T] [Fintype U(*@$_1$@*)] [DecidableEq U(*@$_1$@*)] [Fintype U(*@$_2$@*)] [DecidableEq U(*@$_2$@*)]
    [Fintype W] [DecidableEq W]
    {(*@$\sigma$@*) : FlagType T} (F(*@$_1$@*) : Flag (*@$\sigma$@*) U(*@$_1$@*)) (F(*@$_2$@*) : Flag (*@$\sigma$@*) U(*@$_2$@*)) (G : Flag (*@$\sigma$@*) W) : (*@$\mathbb{Q}$@*) :=
  flagListDensity [F(*@$_1$@*), F(*@$_2$@*)](*@$^\mathtt{f}$@*) G
\end{lstlisting}
Here the superscripted bracket notation is a tuple constructor for
\lean{FlagList}.  Thus, \lean{[F]}$^{\mathtt{f}}$ is the one-entry
\lean{FlagList} whose unique entry is \lean{F}, and
\lean{[F}$_1$\lean{, F}$_2$\lean{]}$^{\mathtt{f}}$ is the two-entry
\lean{FlagList} sending the first index of \lean{Fin~2} to \lean{F}$_1$ and
the second to \lean{F}$_2$.  The superscript \(\mathtt{f}\) distinguishes this
notation from ordinary list notation and from quotient brackets; it does not
add another quotienting operation.

Building on the definitions above, we prove the following chain rule
(\Cref{lem:chain-rule}) and several of its variants.
Although the statement is the elementary counting identity of
\Cref{sec:background}, its formal proof must account for the quotient
representation:
the counts are defined on concrete labeled graphs, and the identity between
them comes from an explicit bijection between the choices the two sampling
procedures make (\Cref{sec:proof-style}).
\begin{lstlisting}
theorem flagDensity_eq_sum_density_prods
    {(*@$\ell_0$@*) (*@$\ell_1$@*) (*@$\ell$@*) (*@$\ell'$@*) : (*@$\mathbb{N}$@*)}
    {(*@$\sigma$@*) : FlagType (Fin (*@$\ell_0$@*))} (F(*@$_1$@*) : Flag (*@$\sigma$@*) (Fin (*@$\ell_1$@*))) (G : Flag (*@$\sigma$@*) (Fin (*@$\ell$@*)))
    (h(*@$\ell_1$@*) : (*@$\ell_0$@*) (*@$\leq$@*) (*@$\ell_1$@*)) (h(*@$\ell'$@*) : (*@$\ell_1$@*) (*@$\leq$@*) (*@$\ell'$@*)) (h(*@$\ell$@*) : (*@$\ell'$@*) (*@$\leq$@*) (*@$\ell$@*))
    : flagDensity(*@$_1$@*) F(*@$_1$@*) G = (*@$\sum$@*) G' : Flag (*@$\sigma$@*) (Fin (*@$\ell'$@*)), flagDensity(*@$_1$@*) F(*@$_1$@*) G' * flagDensity(*@$_1$@*) G' G := ...
\end{lstlisting}

\subsection{The Flag Algebra}
\label{sec:formal-flag-algebra}

To define the flag algebra, we first construct the free real vector space on
$\mathcal{F}^\sigma$.  Each element of this vector space is represented as a
finitely supported function from $\mathcal{F}^\sigma$ to $\mathbb{R}$, using
Mathlib's \lean{Finsupp} type.  We call the resulting type
\lean{FlagVector}~$\sigma$; the constructor \lean{basisVector} produces the
basis vector for each $\sigma$-flag:
\begin{lstlisting}
abbrev FlagVector {k : (*@$\mathbb{N}$@*)} ((*@$\sigma$@*) : FlagType (Fin k)) : Type :=
  FinFlag (*@$\sigma$@*) (*@$\to_0$@*) (*@$\mathbb{R}$@*)

def basisVector {k : (*@$\mathbb{N}$@*)} {(*@$\sigma$@*) : FlagType (Fin k)} (F : FinFlag (*@$\sigma$@*)) : FlagVector (*@$\sigma$@*) :=
  Finsupp.single F 1
\end{lstlisting}
Here $\to_0$ denotes the type of finitely supported functions, and \lean{Finsupp.single F 1} constructs the basis vector $\mathbf e_F$ for $F \in \mathcal{F}^\sigma$: the function taking value $1$ at $F$ and $0$ elsewhere.

Next, we define the zero space $\mathcal{Z}^\sigma$.  The auxiliary function
\lean{flagExpansion F}~$\ell$ is the Lean version of
$\sum_{G \in \mathcal{F}^\sigma_\ell} \den{F}{G}\cdot \mathbf e_G$, the expansion of
$F$ at size $\ell$.  We encode the generating element of
Equation~\eqref{eq:zero-element} as \lean{zeroElement}, \lean{zeroSet} $\sigma$
collects all elements of this form, and \lean{ZeroSpace} $\sigma$ is the
linear subspace they span, represented in Lean by a \lean{Submodule}:
\begin{lstlisting}
noncomputable def flagExpansion {k : (*@$\mathbb{N}$@*)} {(*@$\sigma$@*) : FlagType (Fin k)} (F : FinFlag (*@$\sigma$@*)) ((*@$\ell$@*) : (*@$\mathbb{N}$@*))
    : FlagVector (*@$\sigma$@*) :=
  (*@$\sum$@*) F' : Flag (*@$\sigma$@*) (Fin (*@$\ell$@*)), flagDensity(*@$_1$@*) F.2 F' (*@$\leansmul$@*) basisVector (*@$\langle$@*)(*@$\ell$@*), F'(*@$\rangle$@*)

noncomputable def zeroElement {k : (*@$\mathbb{N}$@*)} {(*@$\sigma$@*) : FlagType (Fin k)} (F : FinFlag (*@$\sigma$@*)) ((*@$\ell$@*) : (*@$\mathbb{N}$@*))
    : FlagVector (*@$\sigma$@*) :=
  basisVector F - flagExpansion F (*@$\ell$@*)

noncomputable def zeroSet {k : (*@$\mathbb{N}$@*)} ((*@$\sigma$@*) : FlagType (Fin k))
    : Set (FlagVector (*@$\sigma$@*)) :=
  {z | (*@$\exists$@*) (F : FinFlag (*@$\sigma$@*)) ((*@$\ell$@*) : (*@$\mathbb{N}$@*)), F.1 (*@$\leq$@*) (*@$\ell$@*) (*@$\wedge$@*) z = zeroElement F (*@$\ell$@*)}

noncomputable def ZeroSpace {k : (*@$\mathbb{N}$@*)} ((*@$\sigma$@*) : FlagType (Fin k))
    : Submodule (*@$\mathbb{R}$@*) (FlagVector (*@$\sigma$@*)) :=
  Submodule.span (*@$\mathbb{R}$@*) (zeroSet (*@$\sigma$@*))
\end{lstlisting}
In \lean{flagExpansion}, the black dot is Lean's scalar multiplication
operation on \lean{FlagVector}s: it scales the basis vector indexed by
\(\langle\ell, F'\rangle\) by the density \lean{flagDensity$_1$ F.2 F'}.

Following the same pattern as the definition of \lean{Flag} from
\lean{LabeledGraph} in \Cref{sec:flagtypes}, we now define
the flag algebra $\mathcal{A}^\sigma$ as a quotient of \lean{FlagVector} $\sigma$.
The equivalence relation \lean{flagVectorEqv} declares two elements of \lean{FlagVector} $\sigma$ equivalent
when their difference lies in \lean{ZeroSpace} $\sigma$, and \lean{flagVectorSetoid} $\sigma$
packages this relation into a setoid so that Lean's quotient type can be applied.
\begin{lstlisting}
def flagVectorEqv {k : (*@$\mathbb{N}$@*)} {(*@$\sigma$@*) : FlagType (Fin k)} (f g : FlagVector (*@$\sigma$@*)) : Prop :=
  f - g (*@$\in$@*) ZeroSpace (*@$\sigma$@*)

instance flagVectorSetoid {k : (*@$\mathbb{N}$@*)} ((*@$\sigma$@*) : FlagType (Fin k)) : Setoid (FlagVector (*@$\sigma$@*)) where
  r     := flagVectorEqv
  iseqv := ... -- proof that flagVectorEqv is an equivalence relation

abbrev FlagAlgebra {k : (*@$\mathbb{N}$@*)} ((*@$\sigma$@*) : FlagType (Fin k)) : Type :=
  Quotient (flagVectorSetoid (*@$\sigma$@*))
\end{lstlisting}
Since \lean{FlagAlgebra}~$\sigma$ is another quotient type, if $f$ is a
\lean{FlagVector}~$\sigma$, the unsubscripted notation $\llbracket f\rrbracket$ denotes
its class in this quotient.

The quotient therefore imposes a family of expansion identities: for every
$\sigma$-flag \lean{F : FinFlag}~$\sigma$ and every size $\ell \geq F.1$, the class of
\lean{basisVector F} is identified with the class of
\lean{flagExpansion F}~$\ell$, that is, with the linear combination of all
size-$\ell$ $\sigma$-flags weighted by their densities over \(F\).

Formalizing multiplication in $\mathcal{A}^\sigma$ also requires several steps.
The term \lean{flagMulWithSize} defines the product
of two $\sigma$-flags at a chosen output size $\ell$ as in
Equation~\eqref{eq:mul}. The term \lean{flagMul} specializes this product to the minimal admissible
size $\ell = F.1 + F'.1 - k$. Then, \lean{bilinearExtension} extends
\lean{flagMul} linearly in each argument to give a product on
\lean{FlagVector} $\sigma$.  This intermediate product is only a bilinear
operation on formal sums; the quotient is where the usual algebra laws are
established.  Finally, \lean{Quotient.map$_2$} lifts the product to
\lean{FlagAlgebra} $\sigma$.
\begin{lstlisting}
noncomputable def flagMulWithSize {k : (*@$\mathbb{N}$@*)} {(*@$\sigma$@*) : FlagType (Fin k)} (F F' : FinFlag (*@$\sigma$@*)) ((*@$\ell$@*) : (*@$\mathbb{N}$@*))
    : FlagVector (*@$\sigma$@*) :=
  (*@$\sum$@*) G : Flag (*@$\sigma$@*) (Fin (*@$\ell$@*)), flagDensity(*@$_2$@*) F.2 F'.2 G (*@$\leansmul$@*) basisVector (*@$\langle$@*)(*@$\ell$@*), G(*@$\rangle$@*)

noncomputable def flagMul {k : (*@$\mathbb{N}$@*)} {(*@$\sigma$@*) : FlagType (Fin k)} (F F' : FinFlag (*@$\sigma$@*)) : FlagVector (*@$\sigma$@*) :=
  flagMulWithSize F F' (F.1 + F'.1 - k)

noncomputable instance {k : (*@$\mathbb{N}$@*)} ((*@$\sigma$@*) : FlagType (Fin k)) : Mul (FlagVector (*@$\sigma$@*)) where
  mul := bilinearExtension flagMul

noncomputable instance {k : (*@$\mathbb{N}$@*)} ((*@$\sigma$@*) : FlagType (Fin k)) : Mul (FlagAlgebra (*@$\sigma$@*)) where
  mul := by
    apply Quotient.map(*@$_2$@*) ((*@$\cdot$@*) * (*@$\cdot$@*))
    -- proof omitted: for f f' g g' : FlagVector (*@\textcolor{gray}{$\sigma$}@*),
    --   flagVectorEqv f f' and flagVectorEqv g g' imply flagVectorEqv (f * g) (f' * g')
\end{lstlisting}
Here \lean{Quotient.map$_2$} is a Lean built-in function that lifts a binary
operation to a quotient type given a proof that the operation respects the
equivalence relation.  In the displayed definition, the expression
\lean{($\cdot$ * $\cdot$)} is not a new operation: since its two arguments are
flag vectors, Lean resolves \lean{*} using the immediately preceding
\lean{Mul (FlagVector}~$\sigma$\lean{)} instance, whose multiplication is
\lean{bilinearExtension flagMul}.  The quotient instance then lifts exactly
this vector-level multiplication to \lean{FlagAlgebra} $\sigma$.  The proof
obligation amounts to showing that this multiplication respects
\lean{flagVectorEqv}, as noted in the comment above.  The load-bearing lemma
is that \lean{ZeroSpace} $\sigma$ is
closed under multiplication by arbitrary flag vectors:
\begin{lstlisting}
theorem flagVector_mul_zeroSpace {k : (*@$\mathbb{N}$@*)} {(*@$\sigma$@*) : FlagType (Fin k)}
    (f : FlagVector (*@$\sigma$@*)) {z : FlagVector (*@$\sigma$@*)} (hz_zero : z (*@$\in$@*) ZeroSpace (*@$\sigma$@*))
    : f * z (*@$\in$@*) ZeroSpace (*@$\sigma$@*) := ...
\end{lstlisting}
Together with commutativity, this says that \lean{ZeroSpace} $\sigma$ is an
ideal in the algebraic sense: multiplying a zero-space element by any flag
vector still gives a zero-space element.  This is exactly the property needed
to define multiplication on the quotient.  If two representatives differ by a
zero-space element, then multiplying both by the same vector still gives
representatives that differ by a zero-space element, so the product does not
depend on the chosen representatives.  This is the compatibility condition
that allows the Lean definition to use \lean{Quotient.map$_2$}.  A separate
theorem handles the choice of output size:
\begin{lstlisting}
theorem flagMulWithSize_indep_on_size {k : (*@$\mathbb{N}$@*)} {(*@$\sigma$@*) : FlagType (Fin k)}
    {F(*@$_1$@*) F(*@$_2$@*) : FinFlag (*@$\sigma$@*)} {(*@$\ell_1$@*) (*@$\ell_2$@*) : (*@$\mathbb{N}$@*)}
    (h(*@$\ell_1$@*) : F(*@$_1$@*).1 + F(*@$_2$@*).1 (*@$\leq$@*) (*@$\ell_1$@*) + k) (h(*@$\ell_2$@*) : F(*@$_1$@*).1 + F(*@$_2$@*).1 (*@$\leq$@*) (*@$\ell_2$@*) + k)
    : flagVectorEqv (flagMulWithSize F(*@$_1$@*) F(*@$_2$@*) (*@$\ell_1$@*)) (flagMulWithSize F(*@$_1$@*) F(*@$_2$@*) (*@$\ell_2$@*)) := ...
\end{lstlisting}
The hypotheses \lean{h$\ell_1$} and \lean{h$\ell_2$} are the admissibility
assumptions on the two output sizes: \(\ell_1\) and \(\ell_2\) must both be
large enough to contain the two input $\sigma$-flags after identifying their common flag
type \(\sigma\).  Here \(F_j.1\) is the total number of vertices of \(F_j\),
and \(k\) is the number of labeled vertices in \(\sigma\).  Equivalently, the intended output
size is at least \(F_1.1 + F_2.1 - k\): the unlabeled vertices outside the flag
type are sampled disjointly in the two $\sigma$-flags, while the \(k\) labeled vertices
are shared.  The theorem then says that any two such admissible output sizes
give equivalent flag vectors, justifying the choice of the minimal size in
\lean{flagMul}.

With these definitions in place, we prove that \lean{FlagAlgebra} $\sigma$ is
an algebra over the real numbers and that its multiplication is commutative,
completing the formalization of the algebraic structure of
$\mathcal{A}^\sigma$.

\subsection{Positive Homomorphisms and Semantic Order}

This subsection describes how we encode the semantic side of flag algebras in
Lean: sequences of finite $\sigma$-flags, their limiting density functions, positive
homomorphisms, and the order induced by testing against all positive
homomorphisms.  The convergence theorem from
Theorem~\ref{thm:convergent-hom} is the main mathematical bridge between finite
$\sigma$-flag sequences and positive homomorphisms, but the formalization first has to
make each of these surrounding notions explicit.

A sequence of $\sigma$-flags indexed by $\mathbb{N}$ is represented by
\lean{FlagSeq}~$\sigma$.  The predicate \lean{ConvergesTo} says that such a
sequence has a prescribed limiting density function.
\begin{lstlisting}
abbrev FlagSeq {k : (*@$\mathbb{N}$@*)} ((*@$\sigma$@*) : FlagType (Fin k)) :=
  (*@$\mathbb{N}$@*) (*@$\to$@*) FinFlag (*@$\sigma$@*)

def ConvergesTo {k : (*@$\mathbb{N}$@*)} {(*@$\sigma$@*) : FlagType (Fin k)} (s : FlagSeq (*@$\sigma$@*)) (a : FinFlag (*@$\sigma$@*) (*@$\to$@*) (*@$\mathbb{R}$@*)) : Prop :=
  Increases s (*@$\wedge$@*) Tendsto (flagDensitySeq s) atTop ((*@$\mathcal{N}$@*) a)
\end{lstlisting}
Here the candidate limit \(a\) is a real-valued function on finite
\(\sigma\)-flags.  The first conjunct, \lean{Increases s}, asserts that the
underlying sizes of the \(\sigma\)-flags \(s(0),s(1),\ldots\) strictly increase.  The
second conjunct uses Lean's filter notation for convergence:
\lean{atTop} is the limit \(n\to\infty\) on \(\mathbb{N}\), and
\(\mathcal{N} a\) is the neighborhood filter of the function~\(a\).  The term
\lean{flagDensitySeq s} is the sequence of density functions associated to
\(s\); its \(n\)-th value sends a finite \(\sigma\)-flag \(F\) to
\(\den{F}{s(n)}\).  Thus, the \lean{Tendsto} line is the Lean statement that,
for every finite \(\sigma\)-flag~\(F\),
\[
  \lim_{n\to\infty}\den{F}{s(n)}=a(F).
\]

Mathematically, a positive homomorphism is a map
\(\phi:\mathcal{A}^\sigma\to\mathbb{R}\) that preserves addition,
multiplication, multiplication by real scalars, and the unit, and satisfies
\(\phi([F])\ge 0\) for every finite \(\sigma\)-flag \(F\).  The Lean definition
follows this sentence almost literally:
\begin{lstlisting}
abbrev Hom {k : (*@$\mathbb{N}$@*)} ((*@$\sigma$@*) : FlagType (Fin k)) :=
  FlagAlgebra (*@$\sigma$@*) (*@$\to_a$@*)[(*@$\mathbb{R}$@*)] (*@$\mathbb{R}$@*)

def PositiveHom {k : (*@$\mathbb{N}$@*)} ((*@$\sigma$@*) : FlagType (Fin k)) : Type :=
  { (*@$\phi$@*) : Hom (*@$\sigma$@*) // (*@$\forall$@*) F : FinFlag (*@$\sigma$@*), (*@$\phi$@*) (*@$\llbracket$@*)basisVector F(*@$\rrbracket$@*) (*@$\geq$@*) 0 }
\end{lstlisting}
The arrow \lean{$\to_a$[$\mathbb{R}$]} in the first definition is Lean's
notation for the type of functions preserving addition, multiplication,
multiplication by real scalars, and the unit.  Thus, \lean{Hom}
abbreviates the type of these structure-preserving maps from
\lean{FlagAlgebra} to \(\mathbb{R}\), with the flag type \(\sigma\) supplied as
its argument.  The second definition is Lean's subtype construction.
Mathematically, we read it in the usual set-comprehension form
\[
  \{\phi\in \operatorname{Hom}(\mathcal{A}^\sigma,\mathbb{R})
    \mid \phi([F])\ge 0
    \text{ for every }F\in\mathcal{F}^{\sigma}\}.
\]
In other words, the only extra condition imposed by \lean{PositiveHom} is
non-negativity on flag basis elements.  In Lean syntax, the expression after
the subtype separator \lean{//},
\[
  \lean{$\forall$ F : FinFlag $\sigma$, $\phi$ $\llbracket$basisVector F$\rrbracket$ $\geq$ 0},
\]
is the predicate that imposes this condition.
Here \lean{basisVector F} is the formal basis vector representing \(F\), and
the unsubscripted quotient brackets send that flag vector to its class in the flag
algebra before \(\phi\) evaluates it.

Using these definitions, we formalize Theorem~\ref{thm:convergent-hom} as two
separate theorems.\footnote{For readability, the displayed statements give the
limit profile the ambient function type
\lean{FinFlag $\sigma$ $\to$ $\mathbb{R}$}.  The Lean implementation instead
packages such profiles in the auxiliary subtype
\lean{FlagDensitySpace}~\(\sigma\), which records that every value lies in
\([0,1]\), and packages the profile induced by a positive homomorphism as
\lean{$\phi$.coe}.  Accordingly, the implemented conclusions are written more
compactly as \lean{$\phi$.coe = a} and
\lean{ConvergesTo s $\phi$.coe}.  The formulations are equivalent: every flag
density lies in \([0,1]\), so the pointwise limit of a convergent flag sequence
does as well.}
\begin{lstlisting}
(*@\textcolor{gray}{\itshape -- Theorem~\ref{thm:convergent-hom} (a)}@*)
theorem flagSeq_limit_mem_positiveHom {k : (*@$\mathbb{N}$@*)} {(*@$\sigma$@*) : FlagType (Fin k)}
    (s : FlagSeq (*@$\sigma$@*)) {a : FinFlag (*@$\sigma$@*) (*@$\to$@*) (*@$\mathbb{R}$@*)} (hs_conv : ConvergesTo s a)
    : (*@$\exists$@*) ((*@$\phi$@*) : PositiveHom (*@$\sigma$@*)), (*@$\forall$@*) (F : FinFlag (*@$\sigma$@*)), (*@$\phi$@*) (*@$\llbracket$@*)basisVector F(*@$\rrbracket$@*) = a F := ...

(*@\textcolor{gray}{\itshape -- Theorem~\ref{thm:convergent-hom} (b)}@*)
theorem positiveHom_as_flagSeq_limit {k : (*@$\mathbb{N}$@*)} {(*@$\sigma$@*) : FlagType (Fin k)}
    ((*@$\phi$@*) : PositiveHom (*@$\sigma$@*))
    : (*@$\exists$@*) (s : FlagSeq (*@$\sigma$@*)), ConvergesTo s (fun F => (*@$\phi$@*) (*@$\llbracket$@*)basisVector F(*@$\rrbracket$@*)) := ...
\end{lstlisting}
Here \(a\) is the same real-valued function on finite \(\sigma\)-flags that
appears in the definition of \lean{ConvergesTo}.  The first theorem assumes
\lean{ConvergesTo s a} and produces a positive homomorphism \(\phi\) satisfying
\[
  \phi\bigl(\llbracket\lean{basisVector F}\rrbracket\bigr)=a(F)
  \qquad\text{for every finite \(\sigma\)-flag \(F\)}.
\]
Thus, the limiting density profile of every convergent \(\sigma\)-flag
sequence is induced by a positive homomorphism.  Conversely, the second
theorem starts with a positive homomorphism \(\phi\) and produces a
\(\sigma\)-flag sequence \(s\) satisfying
\[
  \lean{ConvergesTo}\;s\;
    \bigl(F\mapsto\phi(\llbracket\lean{basisVector F}\rrbracket)\bigr).
\]
In other words, the density profile induced by every positive homomorphism
arises as the limit of a \(\sigma\)-flag sequence.
Both directions require nontrivial formalization work.  In part (a), one must
show that the limit \(a\) of a \(\sigma\)-flag sequence \(s\) satisfies all algebraic
properties of a positive homomorphism, including details that informal proofs
leave implicit.  Part~(b) goes in the opposite direction.  Starting from a
positive homomorphism~\(\phi\), the proof uses the values~\(\phi([G])\) to put a
probability distribution on the finite \(\sigma\)-flags of each sufficiently
large size: a \(\sigma\)-flag \(G\) is sampled with weight \(\phi([G])\).  These weights are
indeed probabilities because \(\phi\) is non-negative on flag basis elements and satisfies the normalization
identity
\[
  \sum_{G\in\mathcal{F}^{\sigma}_{\ell}} \phi([G])
  =
  \phi\left(\sum_{G\in\mathcal{F}^{\sigma}_{\ell}} [G]\right)
  =
  \phi(\mathbf{1}_\sigma)
  =
  1.
\]
One then samples an entire sequence
\[
  X=(X_0,X_1,\ldots),
\]
where \(X_n\) is chosen from the distribution on \(\sigma\)-flags of size \(n^2+k\).
The key estimate says that, for each fixed test \(\sigma\)-flag \(F\), the density
\(\den{F}{X_n}\) is very likely to be close to \(\phi([F])\) once \(n\) is large.
Thus, almost every sampled sequence has the limiting density profile prescribed
by \(\phi\).  Selecting one such realization and calling it
\(G_0,G_1,\ldots\) gives the deterministic \(\sigma\)-flag sequence required by
Theorem~\ref{thm:convergent-hom}(b).  The Lean proof formalizes this argument
by constructing the product measure on flag sequences, proving that the
required coordinatewise convergence event has probability one, and selecting
a realization from that event.

We also define the semantic non-negativity cone $\mathcal{C}^\sigma$
and the order it induces on $\mathcal{A}^\sigma$.
Here, 
\lean{semanticCone}~$\sigma$ is the set of elements in
\lean{FlagAlgebra}~$\sigma$ that every positive homomorphism
evaluates non-negatively.  The induced order is installed as a Lean
\lean{LE} instance, with $f \leq g$ defined to mean
$g - f \in \mathcal{C}^\sigma$.
\begin{lstlisting}
def semanticCone {k : (*@$\mathbb{N}$@*)} ((*@$\sigma$@*) : FlagType (Fin k)) : Set (FlagAlgebra (*@$\sigma$@*)) :=
  { f | (*@$\forall$@*) (*@$\phi$@*) : PositiveHom (*@$\sigma$@*), (*@$\phi$@*) f (*@$\geq$@*) 0 }

instance {k : (*@$\mathbb{N}$@*)} ((*@$\sigma$@*) : FlagType (Fin k)) : LE (FlagAlgebra (*@$\sigma$@*)) where
  le := fun f g => g - f (*@$\in$@*) semanticCone (*@$\sigma$@*)
\end{lstlisting}

\subsection{The Downward Operator}
\label{sec:formal-downward}

The downward operator from \Cref{sec:background-downward} forgets the
labels of a typed flag-algebra expression while keeping the correct averaging
factor.  On a single $\sigma$-flag $F$, the intended formula is
\[
  \llbracket [F] \rrbracket_\sigma = q_\sigma(F)\cdot [F|_\emptyset].
\]
The displayed formula specifies the downward operator on flag basis elements.
To implement it, we first define the normalizing factor and the operation of
forgetting labels on concrete $\sigma$-typed labeled graphs, and prove that
both are independent of the chosen representative.  This gives the downward
image of an individual $\sigma$-flag.  We then extend this operation linearly
to flag vectors and prove that two flag vectors representing the same element
of $\mathcal{A}^\sigma$ have the same image in $\mathcal{A}^\emptyset$.  The
result is a well-defined operator
\[
  \mathcal{A}^\sigma\longrightarrow\mathcal{A}^\emptyset.
\]

Fix a flag type $\sigma$ with $k$ vertices. We regard
these vertices of $\sigma$ as labels.  A concrete $\sigma$-typed labeled
graph~$G$ on~$n$ vertices includes an embedding of $\sigma$ into $G$, which
places these labels on $k$ vertices of $G$.  Its normalizing factor is
computed by counting such label placements.  The numerator,
\lean{isomorphismCount G}, counts the injections of the $k$ labels into the
underlying graph that produce a typed labeled graph isomorphic to~$G$.  Its
denominator
\[
  \frac{n!}{(n-k)!}
\]
is the total number of injections of those labels into the $n$ vertices.
Their ratio is $q_\sigma(G)$.

The other ingredient, $F|_\emptyset$, simply forgets the type embedding while
retaining the underlying graph.  The following listing shows the concrete
definitions of the two ingredients, their lifts to isomorphism classes, and
the map \lean{downwardFlag} that combines them.  The subsequent listing carries
out the linear extension and the passage to the flag algebra described above.
\begin{lstlisting}
noncomputable def downwardNormalizingFactor_labeledGraph {k n : (*@$\mathbb{N}$@*)} {(*@$\sigma$@*) : FlagType (Fin k)}
    (G : LabeledGraph (*@$\sigma$@*) (Fin n)) : (*@$\mathbb{Q}$@*) :=
  let num_of_all_injections := n.factorial / (n - k).factorial
  isomorphismCount G / num_of_all_injections

noncomputable def downwardNormalizingFactor {k n : (*@$\mathbb{N}$@*)} {(*@$\sigma$@*) : FlagType (Fin k)}
    : Flag (*@$\sigma$@*) (Fin n) (*@$\to$@*) (*@$\mathbb{Q}$@*) :=
  Quotient.lift (fun G : LabeledGraph (*@$\sigma$@*) (Fin n) => downwardNormalizingFactor_labeledGraph G) (...)

def unlabel_labeledGraph {k : (*@$\mathbb{N}$@*)} {V : Type} {(*@$\sigma$@*) : FlagType (Fin k)} (G : LabeledGraph (*@$\sigma$@*) V)
    : LabeledGraph (*@$\emptyt$@*) V := ...

noncomputable def unlabel {k : (*@$\mathbb{N}$@*)} {V : Type} {(*@$\sigma$@*) : FlagType (Fin k)}
    : Flag (*@$\sigma$@*) V (*@$\to$@*) Flag (*@$\emptyt$@*) V :=
  Quotient.lift (fun G : LabeledGraph (*@$\sigma$@*) V => (*@$\llbracket$@*)unlabel_labeledGraph G(*@$\rrbracket$@*)) (...)

noncomputable def downwardFlag {k n : (*@$\mathbb{N}$@*)} {(*@$\sigma$@*) : FlagType (Fin k)} (F : Flag (*@$\sigma$@*) (Fin n))
    : FlagVector (*@$\emptyt$@*) :=
  downwardNormalizingFactor F (*@$\leansmul$@*) basisVector (*@$\langle$@*)n, unlabel F(*@$\rangle$@*)
\end{lstlisting}

The definitions appear in the order of this construction.  The counting ratio
on concrete $\sigma$-typed labeled graphs is implemented by
\lean{downwardNormalizingFactor\_labeledGraph}, and its lift to isomorphism
classes is \lean{downwardNormalizingFactor}.
Similarly, \lean{unlabel\_labeledGraph} forgets the type embedding at the concrete
level, while \lean{unlabel} lifts this operation to flags.  Finally,
\lean{downwardFlag} forms the empty-type basis vector indexed by the pair
$\langle n,\lean{unlabel F}\rangle$ and scales it by
\lean{downwardNormalizingFactor F}.

We next lift the single-flag construction to the algebra.  The map
\lean{downwardFlagVector} first extends \lean{downwardFlag} linearly from basis
\(\sigma\)-flags to arbitrary $\sigma$-typed flag vectors.  Its output is an empty-type
flag vector; taking its zero-space quotient gives
\lean{downwardFlagVectorQuot}, whose values lie in
\lean{FlagAlgebra}~$\emptyt$.  At this stage the input is still a flag vector
rather than an element of \lean{FlagAlgebra}~$\sigma$.

To quotient the input as well, we must show that the construction is
independent of the chosen representative.  The key respect lemma says that if
two $\sigma$-typed flag vectors differ by an element of the zero space, then
their downward images represent the same element of the empty-type algebra.
This is exactly the compatibility condition needed to define \lean{downward}
by a final quotient lift.
\begin{lstlisting}
noncomputable def downwardFlagVector {k : (*@$\mathbb{N}$@*)} {(*@$\sigma$@*) : FlagType (Fin k)}
    : FlagVector (*@$\sigma$@*) (*@$\to$@*) FlagVector (*@$\emptyt$@*) :=
  linearExtension (fun F : FinFlag (*@$\sigma$@*) => downwardFlag F.2)

noncomputable def downwardFlagVectorQuot {k : (*@$\mathbb{N}$@*)} {(*@$\sigma$@*) : FlagType (Fin k)}
    : FlagVector (*@$\sigma$@*) (*@$\to$@*) FlagAlgebra (*@$\emptyt$@*) :=
  fun f : FlagVector (*@$\sigma$@*) => (*@$\llbracket$@*)downwardFlagVector f(*@$\rrbracket$@*)

lemma downwardFlagVectorQuot_respect_eqv {k : (*@$\mathbb{N}$@*)} {(*@$\sigma$@*) : FlagType (Fin k)}
    {f f' : FlagVector (*@$\sigma$@*)} (h : f - f' (*@$\in$@*) ZeroSpace (*@$\sigma$@*))
    : downwardFlagVectorQuot f = downwardFlagVectorQuot f' := ...

noncomputable def downward {k : (*@$\mathbb{N}$@*)} {(*@$\sigma$@*) : FlagType (Fin k)}
    : FlagAlgebra (*@$\sigma$@*) (*@$\to$@*) FlagAlgebra (*@$\emptyt$@*) :=
  Quotient.lift (fun f : FlagVector (*@$\sigma$@*) => downwardFlagVectorQuot f) (...)

notation "(*@\textcolor{orange!80!black}{$\llbracket$}@*)" f "(*@\textcolor{orange!80!black}{$\rrbracket_0$}@*)" => (downward f)
\end{lstlisting}
The last definition is now a genuine map
\(\lean{FlagAlgebra}~\sigma\to\lean{FlagAlgebra}~\emptyt\), independent of all
representative choices.  The final line introduces the Lean notation
\(\llbracket f\rrbracket_0\) for \lean{downward f}; the subscript $0$
emphasizes that the result lies in the empty-type algebra.  In mathematical
prose we use \(\llbracket f\rrbracket_\sigma\), following
\Cref{sec:background-downward}, where the subscript instead records the type
being forgotten.  Altogether, the implementation consists of a quotient lift
from concrete labeled graphs to flags, a linear extension to flag vectors,
and a second quotient lift to the flag algebra.

Having constructed the algebraic map, we turn to its semantic meaning.  The
key result is the formal random-extension identity from
Equation~\eqref{eq:ensemble}; preservation of semantic non-negativity will
follow from it.

For the measure-theoretic statement, the implementation represents a random
positive homomorphism by its density profile.  The sample space
\lean{PositiveHomSpace}~$\sigma$ is the measurable subspace of
\([0,1]^{\mathcal{F}^\sigma}\) consisting of profiles induced by positive
homomorphisms.  Each point $\phi$ of this space determines a positive
homomorphism, and \lean{PositiveHomSpace.toPosHom $\phi$ f} evaluates that
homomorphism on $f$.  Thus, a probability measure on
\lean{PositiveHomSpace}~$\sigma$ can be read simply as a distribution of
positive homomorphisms.  With this representation, the random-extension
identity is stated as follows.
\begin{lstlisting}
theorem (*@\texttt{exists\_probMeasure\_extend\_emptyType\_positiveHom}@*) {k : (*@$\mathbb{N}$@*)} {(*@$\sigma$@*) : FlagType (Fin k)}
    {(*@$\phi_0$@*) : PositiveHom (*@$\emptyt$@*)} (h(*@$\sigma$@*) : (*@$\phi_0$@*) ((*@$\langle\sigma\rangle_0$@*)) > 0)
    : (*@$\exists$@*) (*@$\mathbb{P}$@*) : ProbabilityMeasure (PositiveHomSpace (*@$\sigma$@*)), (*@$\forall$@*) f : FlagAlgebra (*@$\sigma$@*),
        (*@$\int$@*) ((*@$\phi$@*) : PositiveHomSpace (*@$\sigma$@*)), PositiveHomSpace.toPosHom (*@$\phi$@*) f (*@$\partial$@*)(*@$\mathbb{P}$@*)
        = (*@$\phi_0$@*) (*@$\llbracket$@*)f(*@$\rrbracket_0$@*) / (*@$\phi_0$@*) (*@$\llbracket$@*)(1 : FlagAlgebra (*@$\sigma$@*))(*@$\rrbracket_0$@*)
  := ...
\end{lstlisting}
Mathlib writes \lean{$\int$ x, g x $\partial\mu$} for the integral of
$x\mapsto g(x)$ with respect to $\mu$; here $\partial\mu$ specifies the
measure and is not a derivative.  In the displayed theorem, $\phi$ ranges over
\lean{PositiveHomSpace}~$\sigma$, the integrand evaluates the corresponding
positive homomorphism on $f$, and the measure is $\mathbb{P}$.  Because
$\mathbb{P}$ is a probability measure, the left-hand side is the expectation
\[
  \mathbb{E}_{\phi\sim \mathbb{P}}[\phi(f)].
\]
The theorem therefore says exactly that $\mathbb{P}$ satisfies
Equation~\eqref{eq:ensemble}, expressed using Lean's downward notation:
\[
  \mathbb{E}_{\phi\sim \mathbb{P}}[\phi(f)]
  =
  \frac{\phi_0(\llbracket f\rrbracket_0)}
       {\phi_0(\llbracket \mathbf{1}_\sigma\rrbracket_0)} .
\]
Thus, $\mathbb{P}$ realizes the random-extension measure
$\Ext_\sigma(\phi_0)$ from \Cref{sec:background-downward}.  Equivalently,
\[
  \phi_0(\llbracket f\rrbracket_0)
  =
  \phi_0(\llbracket \mathbf{1}_\sigma\rrbracket_0)
  \mathbb{E}_{\phi\sim\mathbb{P}}[\phi(f)].
\]
This rearranged form explains the probabilistic construction behind the
formal proof.  Take a convergent sequence of finite graphs representing
$\phi_0$; in each graph, choose an embedding of $\sigma$ uniformly at random,
use the embedded vertices as labels, and evaluate $f$.  Averaging and passing
to the limit produces the expectation above, while the positive factor
$\phi_0(\llbracket \mathbf{1}_\sigma\rrbracket_0)$ accounts for the
normalization of the random embedding.

The random-extension identity now yields the desired transfer property.  If
$f$ is non-negative under every positive homomorphism of the $\sigma$-typed
algebra, then every value $\phi(f)$ in the expectation is non-negative.  The
rearranged identity therefore shows that the downward image of $f$ is
non-negative under every empty-type positive homomorphism.  This is the Lean
counterpart of \Cref{thm:downward-nonneg}:
\begin{lstlisting}
theorem downward_preserve_semanticCone {k : (*@$\mathbb{N}$@*)} {(*@$\sigma$@*) : FlagType (Fin k)}
    (f : FlagAlgebra (*@$\sigma$@*)) (hf : f (*@$\in$@*) semanticCone (*@$\sigma$@*))
    : (*@$\llbracket$@*)f(*@$\rrbracket_0$@*) (*@$\in$@*) semanticCone (*@$\emptyt$@*) := ...
\end{lstlisting}
The proof considers an arbitrary empty-type positive homomorphism $\phi_0$.
When $\phi_0(\langle\sigma\rangle_0)>0$, it applies the random-extension
measure above.  When $\phi_0(\langle\sigma\rangle_0)=0$, a separate lemma shows
that every downward image evaluates to zero, so the conclusion is immediate.
The density bounds developed later use this theorem to transfer inequalities
from a labeled algebra to the unlabeled one.

\subsection{Forbidden Graphs and the Ensemble Semantic Order}
\label{sec:forbidden}

In \Cref{sec:background}, a finite forbidden family \(\mathcal{H}\) is fixed
first, and the flag types, flags, algebras, and positive homomorphisms are all
built inside the \(\mathcal{H}\)-free world.  In that setup, an inequality
\(f\leq_{\mathcal{H},\sigma} g\) is already an inequality in the constrained
\(\sigma\)-typed algebra: the forbidden configurations have been removed
before the semantic order is tested.

Our Lean formalization deliberately separates building the algebra from
imposing the constraint.  The core development constructs the ambient flag
algebra \(\mathcal{A}^\sigma\) of all finite simple graphs, with no forbidden
family as a parameter.  A forbidden graph enters only when we compare two
ambient algebra elements \(f,g\in\mathcal{A}^\sigma\).  We restrict attention to
empty-type positive homomorphisms~\(\phi_0\) that assign density zero to every
graph containing the forbidden graph as a (not necessarily induced) subgraph.
For each such \(\phi_0\) for which \(\sigma\) has positive density, take a
sequence of finite graphs whose unlabeled density profiles converge to
\(\phi_0\).  In each graph, choose an embedding of \(\sigma\) uniformly at random
and use the embedded vertices as labels.  The limiting distribution of the
resulting \(\sigma\)-typed density profiles is \(\Ext_\sigma(\phi_0)\), and we
require \(\phi(f)\leq\phi(g)\) with probability one when \(\phi\) is sampled from
this distribution.

Thus, the forbidden-graph condition does not define a separate flag-algebra
universe in the implementation.  Instead, it restricts the empty-type positive
homomorphisms~\(\phi_0\) from which the random extensions used in the comparison
are formed.  The elements \(f\) and \(g\), as well as the definitions of flags,
densities, products, the downward operator, and positive homomorphisms, remain
those of the ambient algebra and can therefore be reused across extremal
problems.

The order used for this a posteriori condition is the
\emph{ensemble semantic order}.  We state it for a single forbidden graph
\(H\); a forbidden family would impose the zero-density condition below for
every member.

\Needspace*{12\baselineskip}
\begin{definition}[Ensemble semantic order]\label{def:ensemble-semantic-order}
Fix a forbidden graph \(H\), a flag type \(\sigma\), and elements
\(f,g\in\mathcal{A}^\sigma\) in the ambient flag algebra.  We write
\(f\leq^{\mathrm{ens}}_{H,\sigma} g\) if, for every empty-type positive
homomorphism \(\phi_0\) such that, viewing graphs as
empty-type flag-algebra elements,
\[
  \phi_0([F])=0
  \quad\text{for every graph \(F\) containing \(H\) as a subgraph}
  \qquad\text{and}\qquad
  \phi_0(\langle\sigma\rangle_0)>0,
\]
the random extension \(\phi\sim\Ext_\sigma(\phi_0)\) satisfies
\[
  \mathbb{P}\bigl[\phi(f)\leq\phi(g)\bigr]=1.
\]
\end{definition}

The first condition restricts \(\phi_0\) to the \(H\)-free unlabeled limits: it
assigns density zero to every graph containing \(H\).  The second condition,
\(\phi_0(\langle\sigma\rangle_0)>0\), says that the underlying unlabeled graph of
\(\sigma\) has positive limiting density along any sequence of finite
graphs converging to \(\phi_0\).  This is precisely the hypothesis under which the random-extension
measure \(\Ext_\sigma(\phi_0)\) is defined.  The final display says that, when a
\(\sigma\)-typed positive homomorphism \(\phi\) is sampled from this measure, the
inequality \(\phi(f)\leq\phi(g)\) holds with probability one.

\begin{remark}[Relation to built-in \(H\)-free semantics]\label{rem:builtin-hfree}
The background approach builds the \(H\)-free condition into the flag algebra
from the outset, whereas the ensemble approach retains the ambient algebra and
imposes the condition only when comparing its elements.  At the empty type,
where no labels are chosen, the two approaches agree, so this distinction does
not affect the final unlabeled Tur\'an-density bounds.  For nonempty types,
however, the two approaches can differ: the built-in approach may fix the
labels at an exceptional vertex configuration, whereas the ensemble approach
obtains its labels by uniform random sampling.  An exceptional choice of labels
may therefore affect the built-in comparison while remaining invisible to the
ensemble comparison.
\Cref{sec:metatheory} gives a precise criterion for when the two approaches
agree and an example showing that they need not.
\end{remark}

The Lean predicate \lean{forbidLE} is the formalization of
\Cref{def:ensemble-semantic-order}:
\begin{lstlisting}
def forbidLEWith {k : (*@$\mathbb{N}$@*)} {(*@$\sigma$@*) : FlagType (Fin k)}
    (C : ForbidCondition) (f g : FlagAlgebra (*@$\sigma$@*)) : Prop :=
  (*@$\forall\;$@*)((*@$\phi_0$@*) : PositiveHom (*@$\emptyt$@*)),
    (*@$\phi_0$@*) (*@$\langle\sigma\rangle_0$@*) > 0
    -> (*@$\,$@*) C (*@$\phi_0$@*)
    -> (*@$\,$@*) (*@$\mathbb{P}$@*)[(*@$\phi_0$@*)] {(*@$\phi$@*) : PositiveHomSpace (*@$\sigma$@*) | (*@$\phi$@*) f (*@$\leq\;$@*)(*@$\phi$@*) g} = 1

def forbidLE {k m : (*@$\mathbb{N}$@*)} {(*@$\sigma$@*) : FlagType (Fin k)}
    (H : SimpleGraph (Fin m)) (f g : FlagAlgebra (*@$\sigma$@*)) : Prop :=
  forbidLEWith (forbiddenCondition H) f g
\end{lstlisting}
Read \lean{forbidLE} first.  Its explicit argument
\lean{H} is the forbidden graph, and \lean{f} and \lean{g} are the two
ambient \(\sigma\)-typed flag-algebra elements being compared.  In particular,
\lean{f} and \lean{g} do not live in a separately built \(H\)-free algebra.
The definition delegates to \lean{forbidLEWith}, the condition-generic version
that takes a \lean{ForbidCondition} as a parameter.  A \lean{ForbidCondition}
is a predicate on empty-type positive homomorphisms,
specifying which empty-type positive homomorphisms are allowed.  The particular condition
\lean{forbiddenCondition H} selects the \(H\)-free ones: it considers only
positive homomorphisms \(\phi_0\) such that
\[
  \phi_0([F])=0
  \qquad
  \text{for every \(\emptyset\)-flag } F
  \text{ whose underlying graph contains } H \text{ as a subgraph}.
\]

The body of \lean{forbidLEWith} ranges over empty-type positive homomorphisms
\lean{$\phi_0$}.  The hypothesis \lean{C $\phi_0$} imposes the chosen condition
on \lean{$\phi_0$}; for \lean{forbidLE H}, this is the \(H\)-free condition
described above.  Recall that \lean{PositiveHomSpace $\sigma$} is the measurable
space of \(\sigma\)-typed density profiles, each of which determines a positive
homomorphism on \(\mathcal{A}^\sigma\).  The measure
\lean{$\mathbb{P}$[$\phi_0$]} is the random-extension measure
\(\Ext_\sigma(\phi_0)\) on this space.  In the final line of the definition, the
set
\[
  \{\phi\mid \phi(f)\leq\phi(g)\}
\]
is therefore the event that the real number assigned to \(f\) is at most the
real number assigned to \(g\).  The equality to \(1\) says that this inequality
holds with probability one when a density profile \(\phi\) is sampled from
\(\Ext_\sigma(\phi_0)\).

What remains is to connect
the semantic predicate above to the corresponding asymptotic bound on induced densities.
This is the Lean counterpart of \Cref{lem:semantic-bound}, with the hypothesis
stated in the ensemble order of \Cref{def:ensemble-semantic-order} rather than
in the built-in order of \Cref{sec:background-semantic-order}.
In the theorem below, \(H\) is
the forbidden graph and \(F\) is the target graph.  The expression \lean{F.toFlagAlgebra} converts
the simple graph \(F\) into the corresponding element of the empty-type flag
algebra \lean{FlagAlgebra $\emptyt$}, whose evaluation is the induced
\(F\)-density.

The formal bridge packages this passage from semantic order to the finite
asymptotic statement:
\begin{lstlisting}
theorem generalizedTuranDensity_le_of_forbidLE
    {n m : (*@$\mathbb{N}$@*)} {H : SimpleGraph (Fin n)}
    {F : SimpleGraph (Fin m)} {c : (*@$\mathbb{R}$@*)} (hc : 0 (*@$\leq$@*) c)
    (h : forbidLE H F.toFlagAlgebra (c (*@$\leansmul$@*) (1 : FlagAlgebra (*@$\emptyt$@*))))
    : generalizedTuranDensity H F (*@$\leq$@*) c := ...
\end{lstlisting}
The hypothesis \lean{h} is an empty-typed \lean{forbidLE} inequality: the target
graph \(F\), viewed as an element of the empty-type flag algebra, is at most
\(c\cdot \mathbf{1}\) under the ensemble semantic order for the forbidden
graph \(H\).  The conclusion is the corresponding
Tur\'an-density bound, written in Lean as
\lean{generalizedTuranDensity~H~F~$\leq$~c}; mathematically, it is the assertion
\[
  \tdensity{F}{H}
  =
  \lim_{N\to\infty}
  \max_{\substack{|V(G)|=N\\ G\text{ is }H\text{-free}}}
  \den{F}{G}
  \leq c .
\]
Consequently, for each later application, it is enough to prove the empty-type
ensemble inequality
\[
  [F]\leq^{\mathrm{ens}}_{H,\emptyset}c\cdot\mathbf{1}.
\]
The theorem above then yields the corresponding Tur\'an-density bound
\(\tdensity{F}{H}\leq c\).  The same theorem performs this final conversion in
every application.

The Tur\'an density \(\tdensity{F}{H}\) is represented in Lean by
\lean{generalizedTuranDensity H F}, whose definition uses Mathlib's
\lean{limUnder}.  This operation returns a real number even before the
normalized extremal numbers have been shown to converge; if they did not
converge, the returned value would have no asymptotic meaning.  We therefore
prove \lean{tendsto\_\allowbreak generalizedTuranDensity}, which shows that the
normalized extremal numbers converge to \(\tdensity{F}{H}\), namely to
\lean{generalizedTuranDensity H F}.  The proof is standard: once
\(n\geq |V(F)|\), the normalized extremal numbers are non-increasing, and they
are bounded below by zero, so they converge.

\section{The Reflection Layer}
\label{sec:reflection}

The finite calculations a formal proof depends on can be discharged in two
ways.  One is to proceed as a
mathematician usually would: when a finite claim is needed, write a proof of
that claim directly.  In the flag-algebra setting, this would mean proving, one
by one, that two particular flags are isomorphic, that a certain density is a
given rational number, or that a product of flag-algebra basis elements expands into a specified
linear combination.  This is possible in principle, but it is the wrong scale:
a single nontrivial application contains many such finite checks, and hand-proving each
one would bury the mathematical argument under routine case analysis.

The second approach is computation-oriented.  Instead of proving each finite
calculation from scratch, we prove once that an algorithm correctly performs the
kind of calculation we need, and then we let Lean run that algorithm inside later
proofs.  This is the idea of \emph{reflection}.  The reflection layer of our
Lean formalization is the computable counterpart of the specification layer
described in \Cref{sec:abstract}.  For each finite operation that would otherwise require a hand proof,
such as testing whether two labeled graphs represent the same flag, computing a
density, or expanding a product,
it provides a program that carries out the calculation and a theorem saying that
the program's output agrees with the abstract definition.  Later proofs can then
run the program and invoke the correctness theorem, rather than reconstructing
the finite verification case by case.

Flag algebras are especially well suited to this style.  In applications, the
candidate bounds and their supporting numerical data are often found by
computer search or semidefinite programming, and validating the resulting
proof already amounts to checking finite data: same-flag tests, subflag
densities, multiplication tables, and rational arithmetic.  A formal
proof should neither trust those external tables nor replay each entry as a
hand proof.  Reflection gives the middle path.

\Cref{sec:abstract} gave the specification layer: definitions close to the
mathematics, where flags are identified up to isomorphism and quotient classes
hide arbitrary choices of representatives.  Those definitions are the right ones
for stating theorems, but many are marked \lean{noncomputable} because they are
specifications rather than efficient programs.  This section develops the
executable mirror of that layer: concrete graph and flag representations,
procedures for isomorphism testing, for computing densities, and for computing
downward normalizing factors, and adequacy theorems
proving that those procedures agree with the quotient-based definitions.

\subsection{Concrete Graph Representations and Finite Search}
\label{sec:reflection-concrete}

The finite coefficients in a flag-algebra proof come from concrete counting
problems.  For example, suppose \(F\) and \(G\) are \(\emptyset\)-flags.  To
compute the density \(\den{F}{G}\), one can conceptually enumerate all
\(|V(F)|\)-vertex subsets of \(G\), regard their induced graphs as
\(\emptyset\)-flags, test which of them are isomorphic to \(F\), count the
successful ones, and divide by
\(\binom{|V(G)|}{|V(F)|}\).  For a \(\sigma\)-flag density, the same experiment
keeps the labeled vertices fixed and chooses only the additional unlabeled
vertices.  The downward normalizing factor \(q_\sigma(F)\) is another finite
count of the same flavor: it counts label placements rather than subflags.

Thus, an executable flag-algebra development needs two basic ingredients.  It
must be able to enumerate the finite objects being counted, such as candidate
induced subflags, and it must be able to decide yes/no questions, such
as whether two flags are isomorphic.  In Lean, these ingredients are usually
provided by typeclass instances.  A \lean{Fintype} instance gives Lean an
explicit finite collection containing every element of a type; a
\lean{Decidable} instance gives a procedure for deciding a proposition.  Once
such instances are available, a definition that says ``count all objects
satisfying this predicate'' can be evaluated by iterating over the finite
collection and filtering by the decision procedure.

The specification layer of \Cref{sec:abstract}, however, is not designed
around producing such executable \lean{Fintype} and \lean{Decidable} instances.
Its starting point is Mathlib's graph type \lean{SimpleGraph V}, where
\lean{V} is the type whose elements are the vertices of the graph and the
adjacency relation \lean{Adj} takes two vertices and returns a proposition:
\begin{lstlisting}
structure SimpleGraph (V : Type*) where
  Adj      : V -> (*@$\,$@*) V -> (*@$\,$@*) Prop
  symm     : Symmetric Adj
  loopless : Irreflexive Adj
\end{lstlisting}
This works well for the mathematical specification, but it is too general to be
the data structure of the evaluator.  The vertex type \lean{V} need not come
with an algorithm for enumerating its vertices, and the relation~\lean{Adj} need
not come with a Boolean decision procedure.  Since \lean{LabeledGraph}s and
\lean{Flag}s are built on top of this graph type, they inherit the same
generality.  Even if the specification layer knows that, for each size, there are
only finitely many flags up to isomorphism, this does not by itself provide the
computational choices needed by the evaluator: a concrete enumeration of
representatives and a Boolean procedure for comparing them up to isomorphism.

The reflection layer therefore does not use \lean{SimpleGraph} directly as the
evaluator's representation.  Instead, it introduces a family of more rigid
edge-set representations: an \(n\)-vertex graph has vertex type \lean{Fin n},
and its edge set is stored as \lean{Finset (Sym2 (Fin n))}, a finite set of
unordered pairs of vertices.  This representation is packaged as the structure
\lean{Sym2Graph}.  We then reuse the same graph structure to define a
concrete flag type, \lean{Sym2FlagType}, and a concrete labeled graph,
\lean{Sym2LabeledGraph}.
\begin{lstlisting}
structure Sym2Graph (n : (*@$\mathbb{N}$@*)) where
  edges       : Finset (Sym2 (Fin n))
  edges_valid : forall e in edges, (*@$\neg$@*) e.IsDiag

abbrev Sym2FlagType (k : (*@$\mathbb{N}$@*)) := Sym2Graph k

structure Sym2LabeledGraph {k : (*@$\mathbb{N}$@*)} ((*@$\sigma$@*) : Sym2FlagType k) (n : (*@$\mathbb{N}$@*)) where
  edges       : Finset (Sym2 (Fin n))
  edges_valid : forall e in edges, (*@$\neg$@*) e.IsDiag
  type_embed  : fromEdgeSet (SetLike.coe (*@$\sigma$@*).edges) (*@$\hookrightarrow$@*)g fromEdgeSet (SetLike.coe edges)
\end{lstlisting}
The expression \lean{SetLike.coe edges} is the coercion that views the
finite edge set \lean{edges : Finset (Sym2 (Fin n))} as an ordinary set of
unordered pairs.  This coercion is needed because Mathlib's constructor
\lean{fromEdgeSet} builds an abstract \lean{SimpleGraph} from a set of
unordered vertex pairs, not from a \lean{Finset}.  Thus,
\lean{fromEdgeSet (SetLike.coe edges)} is the abstract graph decoded from the
concrete edge data.  The field \lean{type\_embed} then says that the graph
decoded from \(\sigma\)'s edge set embeds into the graph decoded from the field
\lean{edges}.  In other words, the concrete
\lean{Finset}-based representation still carries the usual flag-labeling map
used by the specification layer.

In this concrete definition, unlike in a general \lean{SimpleGraph}, Lean can
enumerate all vertices of the type \lean{Fin n}, namely \(0,1,\ldots,n-1\).
It can also check whether two vertices \(u\) and \(v\) are adjacent by testing
whether the unordered pair \(\{u,v\}\) belongs to the explicit finite edge set
\lean{edges}.

These basic choices then scale to the finite searches used by the evaluator.
For example, to compute a density \(\den{F}{G}\),
the evaluator must range over possible placements of \(F\)
inside the host \(G\).  Such a placement first chooses a finite set of vertices
of \(G\) that contains all labeled vertices; this set then determines an
induced subflag of \(G\), which can be tested for isomorphism with \(F\).
The following \lean{Fintype} instance provides the enumeration used to range
over these possible placements:
\begin{lstlisting}
structure Sym2InducedLabeledSubgraph {k : (*@$\mathbb{N}$@*)} {(*@$\sigma$@*) : Sym2FlagType k} {n : (*@$\mathbb{N}$@*)}
    (G : Sym2LabeledGraph (*@$\sigma$@*) n) where
  verts : Finset (Fin n)
  -- G.type_verts is the image of G.type_embed: the labeled vertices of G.
  verts_subset : G.type_verts (*@$\subseteq$@*) verts

def Sym2InducedLabeledSubgraph.edges {k : (*@$\mathbb{N}$@*)} {(*@$\sigma$@*) : Sym2FlagType k} {n : (*@$\mathbb{N}$@*)}
    {G : Sym2LabeledGraph (*@$\sigma$@*) n} (H : Sym2InducedLabeledSubgraph G)
    : Finset (Sym2 (Fin n)) :=
  G.edges.filter (fun (e : Sym2 (Fin n)) =>
    forall (v : Fin n), v (*@$\in$@*) e -> (*@$\,$@*) v (*@$\in$@*) H.verts)

instance {k : (*@$\mathbb{N}$@*)} {(*@$\sigma$@*) : Sym2FlagType k} {n : (*@$\mathbb{N}$@*)} (G : Sym2LabeledGraph (*@$\sigma$@*) n)
    : Fintype (Sym2InducedLabeledSubgraph G) where
  -- This Finset.univ enumerates all vertex subsets V : Finset (Fin n).
  elems := (Finset.univ : Finset (Finset (Fin n))).filterMap
    (fun (V : Finset (Fin n)) =>
      if hV : G.type_verts (*@$\subseteq$@*) V
      then some ({ verts := V, verts_subset := hV } : Sym2InducedLabeledSubgraph G)
      else none)
    (by ...)
  -- complete proves that every H : Sym2InducedLabeledSubgraph G appears in elems.
  complete := fun (H : Sym2InducedLabeledSubgraph G) => ...
\end{lstlisting}
The type \lean{Sym2InducedLabeledSubgraph G} represents induced subflags
of \(G\): it stores the chosen vertex set together with a proof that this set
includes the labeled vertices.  Its edge set is then computed by filtering
\(G\)'s explicit edge set to the edges whose endpoints both lie in the
chosen vertex set.  The \lean{Fintype (Sym2InducedLabeledSubgraph G)} instance
provides an enumeration of all induced subflags of \(G\).  Concretely,
the instance contains an \lean{elems} field listing all elements of the type \lean{Sym2InducedLabeledSubgraph~G},
together with a \lean{complete} proof that this enumeration misses none of
them.  Here the
enumeration is deliberately simple: since the vertices of \(G\) are
\lean{Fin n}, Lean enumerates all subsets of \(\{0,\ldots,n-1\}\), keeps
exactly those containing \lean{G.type\_verts}, and stores the retained subset
with the proof of containment.  Each retained subset is then viewed as an
induced subflag by using \(G\) to compute the corresponding induced edge set.
Thus, the subset enumeration becomes a finite list of candidates
\(H : \lean{Sym2InducedLabeledSubgraph G}\), which is the search space for the
density calculation.  To compute the numerator of \(\den{F}{G}\), we still need
to recognize which of these candidates are isomorphic to the pattern \(F\).
This is the second basic ingredient: a decidable predicate on candidates which,
given \(H\), decides whether the labeled graph associated to \(H\) is
isomorphic to~\(F\), with the type labels respected.

The schematic predicate \lean{candidateMatches F H} below isolates this test.
It is a one-pattern presentation of the isomorphism check that appears inside
the implementation's tuple-level predicate \lean{predIsoSym2LabeledHl}.  It
turns the induced subgraph \(H\) into its associated labeled graph and asks
whether the result is isomorphic to \(F\), with the type labels respected.  As
stated, \lean{candidateMatches F H} is a proposition, not a program.  The
reflection step is to show Lean how to decide it from the concrete
data stored in \(F\), \(G\), and~\(H\).  The following small instance displays
the mechanism used in the implementation: expose the finite vertex set of the
candidate, read adjacency from explicit edge membership, reject candidates of
the wrong size, and then invoke the finite isomorphism search.
\begin{lstlisting}
/-- Schematic one-pattern helper: H is isomorphic to the concrete pattern F. -/
def candidateMatches {k m n : (*@$\mathbb{N}$@*)} {(*@$\sigma$@*) : Sym2FlagType k}
    (F : Sym2LabeledGraph (*@$\sigma$@*) m) {G : Sym2LabeledGraph (*@$\sigma$@*) n} (H : Sym2InducedLabeledSubgraph G)
    : Prop :=
  Nonempty (H.toLabeledSubgraph.coe (*@$\simeq$@*)f F.toLabeledGraph)

/-- The finite edge-set representation turns the predicate into executable code. -/
instance candidateMatchesDecidable {k m n : (*@$\mathbb{N}$@*)} {(*@$\sigma$@*) : Sym2FlagType k}
    (F : Sym2LabeledGraph (*@$\sigma$@*) m) (G : Sym2LabeledGraph (*@$\sigma$@*) n)
    : DecidablePred (fun H : Sym2InducedLabeledSubgraph G => candidateMatches F H) :=
  fun H => by
    -- the vertex set of H is finite
    have : Fintype H.toLabeledSubgraph.subgraph.verts := by ...
    -- adjacency in H is decidable
    have : DecidableRel H.toLabeledSubgraph.coe.graph.Adj := by ...
    -- adjacency in F is decidable
    have : DecidableRel F.toLabeledGraph.graph.Adj := by ...
    exact if hsize : H.verts.card = m then
      (by
        dsimp [candidateMatches] -- unfold the definition of candidateMatches
        infer_instance)
    else
      isFalse (fun hIso => hsize (verts_card_of_coe_iso H F hIso))

/-- A one-pattern illustrative version of the reflection-layer density computation. -/
def sym2LabeledGraphDensity(*@$_1$@*) {k m n : (*@$\mathbb{N}$@*)} {(*@$\sigma$@*) : Sym2FlagType k}
    (F : Sym2LabeledGraph (*@$\sigma$@*) m) (G : Sym2LabeledGraph (*@$\sigma$@*) n)
    [DecidablePred (fun H : Sym2InducedLabeledSubgraph G => candidateMatches F H)]
    : (*@$\mathbb{Q}$@*) :=
  ((Finset.univ : Finset (Sym2InducedLabeledSubgraph G)).filter
    (fun H => candidateMatches F H)).card / Nat.choose (n - k) (m - k)
\end{lstlisting}
The body of \lean{candidateMatches} decodes both sides before comparing them.
The pattern is decoded by \lean{F.toLabeledGraph}, which reads \(F\)'s edge
\lean{Finset} through \lean{fromEdgeSet} and keeps the type embedding, giving
the abstract \lean{LabeledGraph} of \Cref{sec:flagtypes}.  The candidate is
decoded in two steps.  First, \lean{H.toLabeledSubgraph} is the labeled
subgraph of \lean{G.toLabeledGraph} whose vertices are \lean{H.verts} and whose
edges are the induced ones computed above.  Second, \lean{.coe} views that
subgraph as a labeled graph in its own right, on the vertex type
\lean{H.verts}, forgetting that it sits inside \(G\).  The two decoded labeled
graphs are then compared with the isomorphism type \(\simeq_f\) of
\Cref{sec:flagtypes}, and \lean{Nonempty} turns it into the proposition that
some such isomorphism exists.

The proof of the instance is the computational heart of the example.  After the
local \lean{Fintype} and \lean{DecidableRel} instances have been assembled, the
\lean{hsize} branch reaches the \lean{infer\_instance} line; this is the
generic inference step.  It asks typeclass search to decide the
\lean{Nonempty} proposition in the definition of \lean{candidateMatches}, which
says that such an isomorphism exists.  The generic decidability instance does
this by enumerating finitely many graph isomorphisms and checking the
type-preservation condition.
The schematic predicate \lean{candidateMatches~F~H} has therefore been reduced
to a finite search.  The \lean{hsize} branch is not conceptually necessary for the
definition, but it matches the implementation's computation-oriented proof: a
candidate whose vertex set has the wrong cardinality cannot be isomorphic to
\(F\), so it is rejected before the isomorphism search is run.

Once the matching predicate is decidable, the density computation is ordinary
finite evaluation.  The illustrative definition
\lean{sym2LabeledGraphDensity}$_1$
enumerates all candidates \(H : \lean{Sym2InducedLabeledSubgraph G}\), filters
them by \lean{candidateMatches F H}, counts the survivors, and divides by the
number of ways to choose the non-type vertices.  The same pattern is used for
the multi-argument density \(\den{F_1,\ldots,F_t}{G}\) formalized in
\Cref{sec:formal-subflag-densities}.  There the numerator counts tuples of
induced subgraphs that realize \(F_1,\ldots,F_t\) inside \(G\), with the
chosen non-type vertices pairwise disjoint.  The actual predicate
\lean{predIsoSym2LabeledHl} in our Lean formalization is the computable version of this
condition for tuples.  Given the tuple of patterns \lean{Hl}
and a tuple of candidate induced subgraphs \lean{Gl}, it checks, for every
\lean{i : Fin t}, that \lean{Gl i} matches the corresponding pattern
\lean{Hl i}, using the same decidable matching test as above, and that the
candidates are disjoint away from the shared type vertices.  Our
computable implementation of \(\den{F_1,\ldots,F_t}{G}\) uses this
predicate to filter the finite set of candidate tuples, and then counts the
surviving tuples.  In particular, the computable counterparts
of the specification-level densities \(\mathtt{flagDensity}_1\) and
\(\mathtt{flagDensity}_2\) are:
\begin{lstlisting}
def sym2FlagDensity(*@$_1$@*) {k : (*@$\mathbb{N}$@*)} {(*@$\sigma$@*) : Sym2FlagType k} {m n : (*@$\mathbb{N}$@*)}
    (F : Sym2Flag (*@$\sigma$@*) m) (G : Sym2Flag (*@$\sigma$@*) n) : (*@$\mathbb{Q}$@*) := ...

def sym2FlagDensity(*@$_2$@*) {k : (*@$\mathbb{N}$@*)} {(*@$\sigma$@*) : Sym2FlagType k} {m(*@$_0$@*) m(*@$_1$@*) n : (*@$\mathbb{N}$@*)}
    (F(*@$_0$@*) : Sym2Flag (*@$\sigma$@*) m(*@$_0$@*)) (F(*@$_1$@*) : Sym2Flag (*@$\sigma$@*) m(*@$_1$@*)) (G : Sym2Flag (*@$\sigma$@*) n) : (*@$\mathbb{Q}$@*) := ...
\end{lstlisting}
The arguments here are of type \lean{Sym2Flag}, the reflection-layer counterpart
of \lean{Flag} obtained by quotienting \lean{Sym2LabeledGraph} by concrete
isomorphism; \Cref{sec:reflection-adequacy} constructs it.
Note that unlike the specification-level density functions, these definitions do not
carry the \lean{noncomputable} keyword: Lean has an executable algorithm for
evaluating them.  But computability alone does not say that they implement the
density functions from the specification layer.  The correspondence between
this implementation and our noncomputable specification of
\(\den{F_1,\ldots,F_t}{G}\) is established by the adequacy theorems explained
in the next subsection.

\subsection{Adequacy Theorems for Reflection}
\label{sec:reflection-adequacy}

\Cref{sec:flagtypes} defined flags by quotienting labeled graphs by
isomorphism.
The reflection layer repeats the same construction for the concrete graph
representation from \Cref{sec:reflection-concrete}:
\begin{lstlisting}
def sym2LabeledGraphEqv {k : (*@$\mathbb{N}$@*)} {(*@$\sigma$@*) : Sym2FlagType k} {n : (*@$\mathbb{N}$@*)} (G G' : Sym2LabeledGraph (*@$\sigma$@*) n) 
    : Prop := ...

instance sym2LabeledGraphSetoid {k : (*@$\mathbb{N}$@*)} ((*@$\sigma$@*) : Sym2FlagType k) (n : (*@$\mathbb{N}$@*)) 
    : Setoid (Sym2LabeledGraph (*@$\sigma$@*) n) where
  r     := sym2LabeledGraphEqv
  iseqv := ... -- reflexivity, symmetry, and transitivity of r

def Sym2Flag {k : (*@$\mathbb{N}$@*)} ((*@$\sigma$@*) : Sym2FlagType k) (n : (*@$\mathbb{N}$@*)) : Type :=
  Quotient (sym2LabeledGraphSetoid (*@$\sigma$@*) n)
\end{lstlisting}
The omitted definition of \lean{sym2LabeledGraphEqv} is standard, and says
that two concrete labeled graphs are equivalent if there
exists a label-preserving graph isomorphism between them.
The displayed setoid instance packages this relation with its equivalence
proofs, and \lean{Sym2Flag} takes the resulting quotient.

We now have two corresponding representations of flags.  The
specification-oriented development of \Cref{sec:flagtypes} uses
\lean{Flag}, while the computation-oriented development
of \Cref{sec:reflection-concrete} uses
\lean{Sym2Flag}.  To use a computation from the reflection layer in a theorem
stated in the
specification layer, we need an explicit map from the latter representation to
the former.  We call this map \lean{Sym2Flag.toFlag}:
\begin{lstlisting}
def Sym2Flag.toFlag {k : (*@$\mathbb{N}$@*)} {(*@$\sigma$@*) : Sym2FlagType k} {n : (*@$\mathbb{N}$@*)} (G : Sym2Flag (*@$\sigma$@*) n) 
    : Flag (fromEdgeSet (SetLike.coe (*@$\sigma$@*).edges)) (Fin n) :=
  Quotient.lift
    (fun G : Sym2LabeledGraph (*@$\sigma$@*) n => (*@$\llbracket$@*)G.toLabeledGraph(*@$\rrbracket$@*))
    ... -- proof that the lift respects the equivalence relation of the quotient
    G
\end{lstlisting}

With \lean{Sym2Flag.toFlag} in hand, we can carry out the reflection described
in \Cref{sec:reflection}.  It asks for a program and a theorem saying that the
program's output agrees with the abstract definition in the specification
layer.
\Cref{sec:reflection-concrete} supplied the programs, and the theorems are the
ones below.  Our
\emph{adequacy theorems} state that
\lean{sym2FlagDensity$_1$} and \lean{sym2FlagDensity$_2$}
agree
with the abstract definitions
of \lean{flagDensity$_1$} and \lean{flagDensity$_2$}
after we decode \lean{Sym2Flag}s back to ordinary
\lean{Flag}s using \lean{Sym2Flag.toFlag}:
\begin{lstlisting}
theorem flagDensity(*@$_1$@*)_eq_sym2FlagDensity(*@$_1$@*) {k : (*@$\mathbb{N}$@*)} {(*@$\sigma$@*) : Sym2FlagType k} {m n : (*@$\mathbb{N}$@*)}
    (F : Sym2Flag (*@$\sigma$@*) m) (G : Sym2Flag (*@$\sigma$@*) n)
    : flagDensity(*@$_1$@*) F.toFlag G.toFlag = sym2FlagDensity(*@$_1$@*) F G := ...

theorem flagDensity(*@$_2$@*)_eq_sym2FlagDensity(*@$_2$@*) {k : (*@$\mathbb{N}$@*)} {(*@$\sigma$@*) : Sym2FlagType k} {m(*@$_0$@*) m(*@$_1$@*) n : (*@$\mathbb{N}$@*)}
    (F(*@$_0$@*) : Sym2Flag (*@$\sigma$@*) m(*@$_0$@*)) (F(*@$_1$@*) : Sym2Flag (*@$\sigma$@*) m(*@$_1$@*)) (G : Sym2Flag (*@$\sigma$@*) n)
    : flagDensity(*@$_2$@*) F(*@$_0$@*).toFlag F(*@$_1$@*).toFlag G.toFlag = sym2FlagDensity(*@$_2$@*) F(*@$_0$@*) F(*@$_1$@*) G := ...
\end{lstlisting}
A computable definition together with its adequacy theorem is what we call a
\emph{reflected} computation.
With these adequacy theorems in place, a proof about a density in the
specification layer can be reduced to a proof about the corresponding
computable density in the reflection layer.  Since this density is computable,
Lean can then run the executable algorithm certified by the adequacy theorem to
check the claimed value.  For example, the density
\[
  \den{\fKtwobull,\,\fbarKtwobull}{\fbarPthreebull}
  = \frac{1}{2}
\]
from \Cref{ex:typed-edge-nonedge-product} is proved as follows:
\begin{lstlisting}
example : flagDensity(*@$_2$@*)
             Sym2Flag_2_1_0_0.toFlag
             Sym2Flag_2_1_0_1.toFlag
             Sym2Flag_3_1_0_1.toFlag = 1 / 2
  := by
  rw [flagDensity(*@$_2$@*)_eq_sym2FlagDensity(*@$_2$@*)]
  decide +kernel
\end{lstlisting}
In the code, the three displayed constants are concrete \lean{Sym2Flag}s in
the reflection layer: \lean{Sym2Flag\_2\_1\_0\_0} corresponds to
\(\fbarKtwobull\), \lean{Sym2Flag\_2\_1\_0\_1} corresponds to \(\fKtwobull\),
and \lean{Sym2Flag\_3\_1\_0\_1} corresponds to \(\fbarPthreebull\).  The
statement applies \lean{toFlag} to each of these concrete flags, so it is a
specification-level statement about
\lean{flagDensity}$_2$.  The command \lean{rw} uses the adequacy theorem
\lean{flagDensity}$_2$\lean{\_eq\_sym2FlagDensity}$_2$ and changes the goal to
the closed reflection-layer computation
\begin{lstlisting}
  sym2FlagDensity(*@$_2$@*)  Sym2Flag_2_1_0_0  Sym2Flag_2_1_0_1  Sym2Flag_3_1_0_1  =  1 / 2
\end{lstlisting}
The tactic \lean{decide +kernel} closes a decidable goal by evaluating it
inside the kernel.  Here it runs the reflection-layer density computation above
and closes the goal, so nothing outside the kernel is trusted.

The downward operator from \Cref{sec:formal-downward} follows the same
reflection pattern.  The specification-level definition uses the abstract
downward normalizing factor \(q_\sigma(F)\), implemented as
\lean{downwardNormalizingFactor}.  Its executable counterpart is
\lean{downwardNormalizingFactor\_Sym2Flag}.  On a concrete
\lean{Sym2LabeledGraph} over a flag type with \(k\) labels, on \(n\) vertices,
Lean enumerates the finite set of type embeddings that
produce an equivalent flag, takes its cardinality, and divides by the total
number \(n!/(n-k)!\) of injections of the \(k\) labels into the \(n\) vertices.
The resulting computation is invariant under flag equivalence, so it can then
be lifted through the \lean{Sym2Flag} quotient.
\begin{lstlisting}
def isoEmbeddingCount_sym2LabeledGraph {k : (*@$\mathbb{N}$@*)} {(*@$\sigma$@*) : Sym2FlagType k} {n : (*@$\mathbb{N}$@*)}
    (G : Sym2LabeledGraph (*@$\sigma$@*) n) : (*@$\mathbb{N}$@*) :=
  (isoSym2TypeEmbeddingSetWithSameGraph G).card

def downwardNormalizingFactor_sym2LabeledGraph {k : (*@$\mathbb{N}$@*)} {(*@$\sigma$@*) : Sym2FlagType k} {n : (*@$\mathbb{N}$@*)}
    (G : Sym2LabeledGraph (*@$\sigma$@*) n) : (*@$\mathbb{Q}$@*) :=
  let num_of_all_injections := n.factorial / (n - k).factorial
  isoEmbeddingCount_sym2LabeledGraph G / num_of_all_injections

def downwardNormalizingFactor_Sym2Flag {k : (*@$\mathbb{N}$@*)} {(*@$\sigma$@*) : Sym2FlagType k} {n : (*@$\mathbb{N}$@*)}
    (F : Sym2Flag (*@$\sigma$@*) n) : (*@$\mathbb{Q}$@*) :=
  Quotient.lift
    (fun G => downwardNormalizingFactor_sym2LabeledGraph G)
    ... -- proof that downwardNormalizingFactor_sym2LabeledGraph respects the equivalence relation
    F
\end{lstlisting}
These are executable \lean{def}s, rather than specification-level
\lean{noncomputable def}s.  The first definition computes a natural number by
taking the cardinality of the finite set
\lean{isoSym2TypeEmbeddingSetWithSameGraph G}.  Its elements are the type
embeddings of \(\sigma\) into \(G\)'s graph, in the sense of the
\lean{type\_embed} field of \lean{Sym2LabeledGraph}, that give a labeled graph
isomorphic to \(G\).  The
second definition divides that count by the number of injections, so its value
is exactly \(q_\sigma(G)\).  The final \lean{Quotient.lift} applies this same
representative-level computation to a \lean{Sym2Flag}; its omitted
representative-invariance argument proves that equivalent representatives give
the same value.

We then prove that the specification-level and executable implementations of
the downward normalizing factor agree.
\begin{lstlisting}
theorem downwardNormalizingFactor_eq {k : (*@$\mathbb{N}$@*)} {(*@$\sigma$@*) : Sym2FlagType k} {n : (*@$\mathbb{N}$@*)} (F : Sym2Flag (*@$\sigma$@*) n)
    : downwardNormalizingFactor F.toFlag = downwardNormalizingFactor_Sym2Flag F := ...
\end{lstlisting}
After decoding a concrete \lean{Sym2Flag} back to an abstract \lean{Flag}, this
adequacy theorem rewrites a goal about the specification-level coefficient as
a goal about the executable function
\lean{downwardNormalizingFactor\_Sym2Flag}.  For instance, for the concrete
$\sigma$-flag corresponding to \(\fbarPthreebull\), where \(\sigma\) is the one-vertex
flag type, Lean proves \(q_\sigma(\fbarPthreebull)=2/3\) by the same
rewrite-and-evaluate pattern:
\begin{lstlisting}
example : downwardNormalizingFactor Sym2Flag_3_1_0_1.toFlag = 2 / 3 := by
  rw [downwardNormalizingFactor_eq]
  decide +kernel
\end{lstlisting}
Here the command \lean{rw} changes the goal to the closed computable equality
\begin{lstlisting}
  downwardNormalizingFactor_Sym2Flag  Sym2Flag_3_1_0_1  =  2 / 3
\end{lstlisting}
and \lean{decide +kernel} evaluates that finite computation.

The reflected normalizing factor is what makes the downward operator executable
on concrete $\sigma$-flags.  The
reflection layer computes the underlying $\emptyset$-flag and its normalizing
factor, and the adequacy theorem certifies their product as the abstract
downward image.  By linearity, a finite linear combination of flag basis
elements is handled term by term, so the downward identities needed in later
flag-algebra proofs can be established by computation rather than by separate
counting arguments.

\subsection{Elaboration-Time Generation of Flags and Theorems}
\label{sec:reflection-generation}

The two examples in the preceding subsection show the basic unit of reflected
computation.  To prove the concrete density
\(\den{\fKtwobull,\,\fbarKtwobull}{\fbarPthreebull}=1/2\), Lean rewrites the
specification-level goal with the density adequacy theorem and then evaluates
the resulting closed \lean{sym2FlagDensity}$_2$ computation.  To prove
\(q_\sigma(\fbarPthreebull)=2/3\), it performs the same two steps with
\lean{downwardNormalizingFactor\_eq}.  The adequacy theorems themselves are
uniform: they prove the reflected evaluators correct for every possible input.
What repeats in a concrete development is not the correctness argument, but
its specialization to each concrete $\sigma$-flag and to each concrete density
instance used by the later algebraic proof.

The generated flags are named by a fixed scheme, already visible in the
constants of \Cref{sec:reflection-adequacy}.  A name such as
\lean{Flag\_n\_k\_m\_i} or \lean{FlagAlgebra\_n\_k\_m\_i} carries four
indices: \(n\) is the size of the flag, \(k\) is the size of its flag type,
\(m\) is the position of the flag type's graph in the enumeration of the
\(k\)-vertex graphs, and \(i\) is the position of the flag itself in the
enumeration of the \(n\)-vertex flags of that flag type.  Both enumerations are
in a fixed canonical order.  Graphs are ordered first by their number of edges
and then by their canonical edge lists.  The canonical edge list of a graph is
the lexicographically smallest sorted edge list obtained by relabeling its
vertices in all possible ways.  The flags over a graph are ordered by the
placement of their labels.  The
reflection-layer constants \lean{Sym2Flag\_n\_k\_m\_i} follow the same
scheme, and the empty flag type is written \(k=m=0\).  Thus,
\lean{Sym2Flag\_3\_1\_0\_1} from \Cref{sec:reflection-adequacy} is flag
number \(1\) among the three-vertex flags of the unique one-vertex type,
namely \(\fbarPthreebull\), and \lean{FlagAlgebra\_3\_1\_0\_1} is the
corresponding basis element of the flag algebra.

We use Lean's metaprogramming facilities to automate this repetition.  They
specialize the uniform adequacy theorems to concrete flags and to the required
single-pattern and pattern-pair density instances.  More precisely, the
generation commands displayed below are custom command
elaborators written using Lean's metaprogramming
framework~\cite{ebner2017structured}.  During
elaboration, the phase in which Lean processes source commands and produces
checked definitions and theorems, these commands run the finite enumerators
and generate proofs of the required specialized facts.  Once the resulting
theorem families have been accepted by Lean, subsequent flag-algebra proofs
use them as rewrite rules instead of evaluating the finite enumerators again.
For the
sizes and flag type used in the two examples of
\Cref{sec:reflection-adequacy}, the relevant commands have
the following form:
\begin{lstlisting}[emph={generate_empty_typed_flags,generate_flags,generate_flag_pair_density_theorems_no_forbid,generate_mul_theorems},emphstyle=\color{blue!70!black}\bfseries]
-- Generate empty-typed flags at pattern size 2 and host size 3.
generate_empty_typed_flags 2
generate_empty_typed_flags 3

-- Use the unique flag type on one labeled vertex (k = 1, m = 0),
-- and generate its flags at the sizes 2 and 3.
generate_flags 2 1 0
generate_flags 3 1 0

-- Prove all pair densities, then all product expansions.
generate_flag_pair_density_theorems_no_forbid 2 3 1 0
generate_mul_theorems 2 3 1 0
\end{lstlisting}

The numerical arguments are the \(n\), \(k\), \(m\) of the naming scheme.  The
command \lean{generate\allowbreak\_empty\allowbreak\_typed\allowbreak\_flags n} generates the empty-typed flags
on \(n\) vertices, and \lean{generate\_flags n k m} generates the \(n\)-vertex
flags of the flag type selected by \(k\) and \(m\); here \(k=1,m=0\) is the unique
flag type on one labeled vertex.  The density and multiplication commands take a
pattern size \(p\) and a host size \(h\), followed by the same \(k,m\)
describing the flag type, and the suffix \lean{no\_forbid} says that no
forbidden-graph filter is being applied.
Here, a completeness lemma states that the generated names exhaust the
relevant finite set: every \(n\)-vertex \(\sigma\)-flag, up to
label-preserving isomorphism, appears among the generated flags.  For the empty
flag type, this means that every \(n\)-vertex graph is represented up to
isomorphism; for a forbid-free enumeration, the same statement is restricted
to the \(H\)-free flags.
\Cref{tab:generation-commands} collects the whole command family, including
the forbid-free variants described below.

\begin{table}
  \centering
  \caption{The generation commands.  \(n\) is a flag size and \(k\), \(m\)
  select the flag type, as in the naming scheme; the density and
  multiplication commands take a pattern size \(p\) and a host size \(h\).
  The forbid-free variants take a further argument \lean{H}, a
  \lean{Sym2Graph} term for the forbidden graph.}
  \label{tab:generation-commands}
  \small
  \begin{tabular}{@{}lp{0.34\textwidth}@{}}
    \toprule
    Command & Generates \\
    \midrule
    \lean{generate\_empty\_typed\_flags n} &
      the empty-typed flags on \(n\) vertices, with the completeness lemma \\
    \lean{generate\_flags n k m} &
      the \(n\)-vertex flags of flag type \((k,m)\), with unlabelings, downward
      coefficients, and the completeness lemma \\
    \lean{generate\_flag\_pair\_density\_theorems\_no\_forbid p h k m} &
      the pair-density theorems for size-\(p\) patterns in size-\(h\) hosts \\
    \lean{generate\_mul\_theorems p h k m} &
      for every pair \(F_1,F_2\in\mathcal F^\sigma_p\), the product expansion
      \([F_1]\cdot[F_2]=\sum_{G\in\mathcal F^\sigma_h} c_G\,[G]\) \\
    \midrule
    \lean{generate\_forbid\_free\_empty\_typed\_flags n H} &
      the \(H\)-free empty-typed flags, with the filtered completeness lemma \\
    \lean{generate\_forbid\_free\_flags n k m H} &
      the \(H\)-free typed flags, with unlabelings, downward coefficients,
      and the filtered completeness lemma \\
    \lean{generate\_forbid\_free\_flag\_pair\_density\_theorems p h k m H} &
      the pair-density theorems over the \(H\)-free flags \\
    \lean{generate\_forbid\_free\_mul\_theorems p h k m H} &
      the product expansions, valid modulo flags containing \(H\) \\
    \lean{generate\_forbid\_free\_flag\_density\_theorems p i h H} &
      the single-flag densities of \lean{Flag\_p\_0\_0\_i} in every
      \(H\)-free size-\(h\) host \\
    \bottomrule
  \end{tabular}
\end{table}

\Needspace*{4\baselineskip}
The important point is what each command contributes to subsequent proofs.  The
first four commands create the named two- and three-vertex flags
used in the examples of \Cref{sec:reflection-adequacy}, and each of the four
also proves the corresponding completeness lemma.  Generating the typed
flags also computes and proves their downward identities, including the
coefficient \(q_\sigma(\fbarPthreebull)=2/3\).  The pair-density command then
generates a theorem \(\lean{flagDensity}_2\,F_1\,F_2\,G=c\) for every unordered
pair of two-vertex patterns~\((F_1,F_2)\) and every three-vertex host \(G\);
this family includes the identity
\(\den{\fKtwobull,\,\fbarKtwobull}{\fbarPthreebull}=1/2\) displayed in
\Cref{sec:reflection-adequacy}.
Finally, the multiplication
command combines those coefficients with the completeness lemma for the hosts
to prove the product expansion
\([F_1] \cdot [F_2]=\sum_G c_G\,[G]\) for every pattern pair; completeness is
what justifies summing over exactly the named hosts.  Thus, elaboration produces the
flags and their corresponding algebra basis elements, their unlabelings and downward coefficients, the density table,
and the multiplication table, together with proofs of all
the entries of these tables.

Extremal applications work under a forbidden-graph assumption, and each of
the commands above has a forbid variant for this situation.  The
\lean{generate\_forbid\_free\_*} commands take one additional argument, a
concrete \lean{Sym2Graph} term for the forbidden graph \(H\); the Mantel
development, for example, defines \lean{K3} as the complete graph on three
vertices and runs \lean{generate\_forbid\_free\_flags 3 1 0 K3}.  These
variants emit constants and theorems only for flags whose underlying graphs
avoid \(H\) as an ordinary, not necessarily induced, subgraph; their
completeness lemmas state that the named flags exhaust precisely this set,
using the same notion of \(H\)-freeness as the forbidden-graph machinery of
\Cref{sec:forbidden}.
One member of the family has no unforbidden counterpart:
\lean{generate\_forbid\_free\_flag\_density\_theorems p i h H} evaluates the
single-flag densities \(\lean{flagDensity}_1\) of the empty-typed flag
\lean{Flag\_p\_0\_0\_i} in every \(H\)-free host of size \(h\).  The
certificate compiler of \Cref{sec:flagmatic} uses this table when the target is
the \(p\)-vertex empty-typed flag indexed by \(i\) and \(p<h\): the table
provides the coefficients needed to express its density in terms of the
\(H\)-free \(h\)-vertex flags.

\Needspace*{4\baselineskip}
\medskip
\noindent\textbf{Optimizations.}
The pipeline employs the following optimizations to reduce the computational
cost of elaboration.  These optimizations do not eliminate the inherent
combinatorial growth of flag enumeration and density computation, but they
make the nontrivial examples considered in this paper feasible.

\begin{itemize}
\item \textbf{Size-inductive enumeration.}
The enumeration of flags is inductive in the size: the underlying
\(n\)-vertex graph representatives are obtained by extending the
\((n-1)\)-vertex representatives by one vertex in all possible ways, and the
typed flags on \(n\) vertices are obtained by placing the flag type's labels on
those representatives in every valid way.  The working lists therefore stay
close to one representative per isomorphism class, rather than covering all
\(2^{\binom{n}{2}}\) graphs on a fixed vertex set.

\item \textbf{Deduplication by cheap invariants.}
Both steps of that induction produce isomorphic duplicates: different vertex
extensions, and different label placements, can yield the same flag.  The
enumeration therefore deduplicates as it goes, keeping a candidate only when
it is isomorphic to no kept representative.  These isomorphism tests are the
expensive part, so each candidate carries a cheap isomorphism-invariant key,
derived from its edge count and degree sequence, and the expensive test runs
only when two keys collide.

\item \textbf{Completeness from a verified enumerator.}
One way to certify a particular generated list would be to enumerate the full
quotient type independently and compare it with that list.  Such a check would
establish completeness only for that particular run and would repeat the
brute-force search that the size-inductive enumeration is designed to avoid.
Instead, we prove once and for all, by induction on the size, that the
enumeration algorithm of our Lean implementation produces a representative of every isomorphism class.
A concrete generation run then needs only to check that the emitted constants
coincide with the algorithm's computed output.  The completeness of that
output follows by specializing the general theorem about the algorithm, rather
than by independently searching the full space again.

\item \textbf{Pruning under a forbidden graph.}
For a complete forbidden graph, the size induction itself prunes: only
\(H\)-free representatives are extended, only \(H\)-free extensions are kept,
and a cheap clique test decides containment, so flags containing \(H\) are
never generated.  For other forbidden graphs, the representative enumeration
is filtered before any typed structure is built, so the labeled-graph work is
spent only on the surviving graphs.  The pruned generation nevertheless
proves the same filtered completeness lemma.  Its enumeration misses no
\(H\)-free isomorphism class because \(H\)-freeness is invariant under
isomorphism and preserved by deleting a vertex; the completeness statement
therefore follows by induction on size.  The pruning algorithm
decides \(H\)-freeness using ordinary subgraph containment.  The filtered
completeness lemma expresses the same condition in flag-algebra terms, by
requiring every induced pattern whose underlying graph contains \(H\) to have
density zero.  Lean proves that these two tests retain exactly the same flags,
so the enumeration can use the executable containment test while later proofs
use the density-based completeness statement.

\end{itemize}

Metaprogramming automates theorem construction, not verification: Lean still
checks every generated proof, and a wrong coefficient makes that proof fail.
Once accepted, these theorems let later arguments simplify downward images,
densities, and products by rewriting rather than reopening the finite searches.
Thus, the story of this section has three steps: concrete representations make
the finite operations executable, adequacy theorems connect those operations
to the abstract definitions, and elaboration-time generation turns their
individual results into tables of proved facts.  The next section
uses those tables to automate complete flag-algebra arguments.

\section{Automating Flag-Algebra Proofs}
\label{sec:flagmatic}

Fix a forbidden family \(\mathcal H\), a target graph \(F\), and a proposed
bound \(c\).  The task addressed in this section is to turn finite certificate
data into a Lean proof of
\[
  [F]\leq_{\mathcal H}c\cdot\mathbf{1},
\]
which expresses the desired asymptotic upper bound on the induced density of
\(F\) in \(\mathcal H\)-free graphs.  A \emph{flag-algebra certificate} is a
finite witness for this inequality: it specifies flags, rational matrices, and
coefficients that purport to express the gap
\(c\cdot\mathbf{1}-[F]\) as a sum of terms that are non-negative under the
\(\mathcal H\)-free condition.  Flagmatic~\cite{vaughan2013flagmatic} uses
semidefinite programming to search for such a witness.

This section concerns certificate verification rather than certificate
search.  Our compiler accepts a certificate produced by Flagmatic as untrusted
input, translates its flags, matrices, and target data into Lean, regenerates
the required flags, recomputes their density and multiplication identities,
verifies the rational matrix factorizations, and assembles the final
inequality.  Automating this translation is important because nontrivial
certificates contain extensive finite bookkeeping that would be tedious and
error-prone to reproduce manually.  The resulting theorem is accepted only
after Lean has checked every mathematical obligation.

We first explain how positive semidefinite matrices justify the inequalities
encoded by a certificate.  We then identify the four obligations common to our
examples and show how the compiler discharges them.  Finally, we state the
trust model and evaluate the compiler on seven certificates.

\subsection{How a Certificate Proves a Bound}
\label{sec:flagmatic-background}

At a high level, a flag-algebra certificate proves
$[F] \leq_{\mathcal H} c\cdot\mathbf{1}$ by constructing two kinds of
non-negative correction term.  One kind consists of terms $Q_t$, each obtained
by applying the downward operator to a positive-semidefinite quadratic
expression.  The other is a remainder $r$, a linear combination of
$\mathcal H$-free flags with non-negative coefficients.  Under the
$\mathcal H$-free condition, these data yield the chain
\[
  [F]
  \leq_{\mathcal H} [F]+\sum_t Q_t
  = c\cdot\mathbf{1}-r
  \leq_{\mathcal H} c\cdot\mathbf{1}.
\]
The first inequality uses the non-negativity of the $Q_t$, the equality is the
finite calculation recorded by the certificate, and the final inequality uses
the non-negativity of $r$.  Equivalently, the certificate decomposes the gap as
$c\cdot\mathbf{1}-[F]=\sum_t Q_t+r$.  Flagmatic~\cite{vaughan2013flagmatic}
searches for this data using an SDP solver and outputs it as a certificate.  Our
compiler consumes the certificate as untrusted input and asks Lean to verify
the non-negativity of the $Q_t$ and $r$ and the intervening equality.

Mantel's theorem gives a small certificate with exactly the structure above.
Here the target is $F=K_2$, whose density is the edge density, and we
abbreviate its basis element $[K_2]$ by
$K_2$.
We work with the one-vertex type $\sigma$, choose a non-negative element
$f\in\mathcal A^\sigma$, and use \Cref{thm:downward-nonneg} to make its
downward image $Q:=\llbracket f\rrbracket_\sigma$ a non-negative correction
term.  We then show that
\[
  K_2 + Q = \tfrac12\cdot\mathbf{1} - \tfrac13\cdot\fbarPthree.
\]
Write $r:=\tfrac13\cdot\fbarPthree$ for the displayed remainder.  The
non-negativity of $Q$ and $r$ then proves the desired bound.

To construct $f$, take the two size-$2$ $\sigma$-flags $\fKtwobull$ and
$\fbarKtwobull$.  We use these pictures directly in the algebraic calculation
below.\footnote{Strictly,
each picture denotes a finite $\sigma$-flag, whereas the algebraic operations
below apply to its quotient class in $\mathcal A^\sigma$.  We follow the usual
convention of omitting brackets around pictorial basis elements.}
Set
\[
  f := \tfrac12((\fKtwobull)-(\fbarKtwobull))^2,
  \qquad
  Q := \llbracket f \rrbracket_\sigma.
\]
For every positive homomorphism $\phi\in\poshom{\sigma}$,
\begin{align*}
  \phi(f)
  = \tfrac12\,\phi\bigl(((\fKtwobull)-(\fbarKtwobull))^2\bigr)
  = \tfrac12\,\bigl(\phi((\fKtwobull)-(\fbarKtwobull))\bigr)^2
  \;\geq\; 0.
\end{align*}
The first equality unfolds the definition of $f$ and uses preservation of
scalar multiplication; the second uses preservation of multiplication, which
rewrites the flag-algebra square as the square of the real number
\(\phi((\fKtwobull)-(\fbarKtwobull))\).  This makes the source of the
non-negativity explicit: the final expression for \(\phi(f)\),
\[
  \tfrac12\bigl(\phi((\fKtwobull)-(\fbarKtwobull))\bigr)^2,
\]
is non-negative because it is a positive scalar multiple of the square of a
real number.  Thus, $f$ is non-negative in $\mathcal A^\sigma$, and
\Cref{thm:downward-nonneg} gives $0\leq_{K_3}Q$.

To calculate $Q$ in the common size-three basis, we use the triangle-free
flags $\fbarKthree$, $\fbarPthree$, and $\fPthree$.  The size-three expansion
of $K_2$ from Equation~\eqref{eq:k2-expansion} is
\[
  K_2 = \tfrac13\cdot\fbarPthree + \tfrac23\cdot\fPthree.
\]
Applying Equation~\eqref{eq:sum-to-one} to the triangle-free empty-type flags of size
three gives
\[
  \mathbf{1} = \fbarKthree + \fbarPthree + \fPthree.
\]
A size-three product expansion followed by the downward operator gives
\begin{align*}
  Q
  &=
  \bigl\llbracket
    \tfrac12\bigl((\fKtwobull)-(\fbarKtwobull)\bigr)^2
  \bigr\rrbracket_\sigma\\
  &=
  \bigl\llbracket
    \tfrac12\left(
       \fbarKthreebull+\fEthreebull
       -\fbarPthreebull-\fPthreebull
       +\fKonetwobull
     \right)
  \bigr\rrbracket_\sigma\\
  &= \tfrac12\left(
       \fbarKthree
       + \tfrac13\cdot\fbarPthree
       - \tfrac23\cdot\fbarPthree
       - \tfrac23\cdot\fPthree
       + \tfrac13\cdot\fPthree
     \right)\\
  &= \tfrac12\cdot\fbarKthree
     -\tfrac16\cdot\fbarPthree
     -\tfrac16\cdot\fPthree.
\end{align*}
The fractions produced by the downward operator are its normalizing factors:
they record whether one, two, or all three choices of the labeled vertex
produce the displayed $\sigma$-flag.

Combining these identities now gives a derivation of the desired bound:
\begin{align*}
  K_2
  &\leq_{K_3} K_2+Q \\
  &= \left(\tfrac13\cdot\fbarPthree + \tfrac23\cdot\fPthree\right)
     + \left(\tfrac12\cdot\fbarKthree
     -\tfrac16\cdot\fbarPthree
     -\tfrac16\cdot\fPthree\right)\\
  &= \tfrac12\cdot\fbarKthree
     + \tfrac16\cdot\fbarPthree
     + \tfrac12\cdot\fPthree\\
  &= \tfrac12\cdot
     \left(\fbarKthree+\fbarPthree+\fPthree\right)
     - \tfrac13\cdot\fbarPthree\\
  &= \tfrac12\cdot\mathbf{1} - \tfrac13\cdot\fbarPthree\\
  &\leq_{K_3} \tfrac12\cdot\mathbf{1}.
\end{align*}
The first inequality adds the non-negative downward square $Q$.  The
successive equalities substitute the expansions of $K_2$ and $Q$, collect their
coefficients, and regroup the result as $\tfrac12\cdot\mathbf{1}$ minus the
remainder using the unit expansion
$\mathbf{1} = \fbarKthree+\fbarPthree+\fPthree$.  The final inequality follows
because the remainder $r=\tfrac13\cdot\fbarPthree$ is non-negative.

The essential certificate data for this example are precisely the nontrivial
terms in this derivation: the two $\sigma$-flags $\fKtwobull$ and
$\fbarKtwobull$, the downward image $Q$ of their halved squared difference, and
the non-negative remainder $r=\tfrac13\cdot\fbarPthree$.  We chose this small
certificate by hand; in larger examples, Flagmatic's SDP search supplies the
analogous flags, positive semidefinite matrices, and remainder coefficients.
Checking even this certificate requires expanding a product, computing its
downward image, expanding the target and the unit to a common host size,
and collecting coefficients.  These calculations remain suitable for
automation but become impractical to perform and transcribe by hand for larger
certificates; they are exactly what our compiler automates.

The same idea scales by replacing one square with weighted sums of many
squares.  Choose $\sigma$-flags $E_1,\ldots,E_r$, let
$e_i:=[E_i]\in\mathcal A^\sigma$ be their basis elements, and collect these
algebra elements into a vector $v=(e_1,\ldots,e_r)^\top$.  A symmetric matrix $M$ is \emph{positive semidefinite}
(PSD) if $x^\top Mx\geq 0$ for every real vector~$x$.  Consequently, when $M=(M_{ij})$
is PSD, the flag-algebra quadratic form
$\langle v,Mv\rangle=\sum_{i,j}M_{ij} e_i e_j$ is non-negative under every
positive homomorphism: evaluating the algebra elements turns it into the ordinary real
quadratic form $x^\top Mx$.  The downward operator then transfers this
non-negativity from the labeled algebra to the unlabeled one.

Flagmatic finds the matrices by turning this decomposition problem into a
\emph{semidefinite program} (SDP).  It first fixes a host size $N$ and chooses
the flag types and flags that may occur.  It expands the target, the unit, and
the downward quadratic forms in the common basis of allowed $N$-vertex graphs,
namely those that avoid $\mathcal H$.  At this point, the entries of the matrices
are unknowns.  Requiring the coefficient left over for every host graph to be
non-negative gives linear constraints on those entries, and requiring each
matrix to be PSD gives the semidefinite constraints.  An SDP solver finds
numerical matrices satisfying these conditions, after which Flagmatic converts
them into exact rational matrices.

A full certificate uses several flag types at once, one PSD matrix per flag
type, so the single matrix $M$ above becomes a family indexed by the flag type.
The index $t$ ranges over the chosen flag types $\sigma_t$, each contributing
its own flag vector $v_t$ and matrix $M_t$.  We call the data for one such
index, the flag type $\sigma_t$ together with $v_t$ and $M_t$, an \emph{SDP block};
the name reflects that the matrices $M_t$ are the diagonal blocks of the
semidefinite program's matrix variable.  The downward quadratic form
block~$t$ produces is
\[
  Q_t := \bigl\llbracket
    \langle v_t, M_t\, v_t\rangle
  \bigr\rrbracket_{\sigma_t}.
\]
The resulting certificate specifies an identity, modulo the $\mathcal H$-free
condition, of the form
\[
  [F] + \sum_t Q_t
  =
  c\cdot\mathbf{1} - r,
\]
where every $M_t$ is PSD and $r$ is a linear combination of $\mathcal H$-free
flags with non-negative coefficients.  Equivalently,
the certificate justifies the derivation
\begin{align*}
  [F]
  \leq_{\mathcal H}
  \bigl([F]+\sum_t Q_t\bigr)
  =
  \bigl(c\cdot\mathbf{1}-r\bigr)
  \leq_{\mathcal H} c\cdot\mathbf{1}.
\end{align*}
The first inequality adds the non-negative downward quadratic forms.  The
equality is the certificate identity checked after all terms have been
expanded in the common host-size basis, and the final inequality follows from
the non-negativity of $r$.  The certificate records the chosen
flags, the rational matrices $M_t$, and the target-density data needed to state
this calculation.  Our compiler does not trust the certificate's derived
density or multiplication tables: it recomputes them in Lean, proves the
matrices are PSD, and checks the calculation.  The rest of this section explains
the four recurring parts of that verification and how the compiler generates
them.

\subsection{The Four-Part Compiler}
\label{sec:flagmatic-shape}
\label{sec:flagmatic-pipeline}
\label{sec:flagmatic-gen}\label{sec:flagmatic-psd}\label{sec:flagmatic-expand}%
\label{sec:flagmatic-close}\label{sec:flagmatic-script}%

We organize the compiler around a four-part proof skeleton, with one part
for each obligation in turning a certificate into a verified proof.  The four
parts are:

\begin{enumerate}
  \item \textbf{Generating the flags and density data.}  The proof must name
    the flags and corresponding flag-algebra basis elements it uses,
    restricted to those that are free of the chosen forbidden graph, and must
    prove the identities for single-flag densities, flag-pair densities, and
    products of flags sharing a common flag type.
    These identities support the algebraic rewrites used later in the proof.
  \item \textbf{Certifying the matrices.}  For each SDP block $t$, the
    certificate supplies a rational matrix $M_t$ and a list of
    $\sigma_t$-flags $E_{t,i}$.  The proof must define the vector~$v_t$ of
    their basis elements $e_{t,i}:=[E_{t,i}]$ and must prove that $M_t$ is
    positive semidefinite.
  \item \textbf{Lifting the target (when needed).}  If the target graph $F$
    has fewer than $N$ vertices, where $N$ is the host size chosen up front to
    bound the certificate search in \Cref{sec:flagmatic-background}, the proof
    must first lift $[F]$ to the
    host-size basis, dropping the contributions that contain a copy of the
    forbidden graph.
  \item \textbf{Assembling the main theorem.}  With
    $Q_t:=\bigl\llbracket\langle v_t,M_t v_t\rangle\bigr\rrbracket_{\sigma_t}$
    and the certificate's non-negative remainder~$r$, the proof must assemble
    the bound $[F] \leq_{\mathcal{H}} c\cdot\mathbf{1}$ through the same chain as
    in \Cref{sec:flagmatic-background}:
    \[
      [F]
      \;\leq_{\mathcal H}\;
      [F] + \sum_t Q_t
      \;=\;
      c\cdot\mathbf{1}-r
      \;\leq_{\mathcal H}\;
      c\cdot\mathbf{1}.
    \]
    The first inequality adds the non-negative downward quadratic forms arising
    from the matrices certified in Part~2.  The middle equality is checked using
    the identities from Part~1 and, when needed, the target-lifting identity
    from Part~3.  The final inequality follows because $r$ has non-negative
    coefficients.
\end{enumerate}

The compiler handles each part with a small set of reusable proof-producing
components.  Consequently, a new theorem primarily requires new certificate
data rather than a new proof architecture.

For the certificate shapes supported by the current implementation, this
four-part process is exposed through the \lean{gen-skeleton} subcommand of
\lean{flagmatic\_to\_lean.py}:
\begin{lstlisting}[language={}]
python flagmatic_to_lean.py gen-skeleton <cert>.json <target>.lean
\end{lstlisting}
Here \lean{<cert>.json} is a placeholder for a Flagmatic certificate file, which
identifies the target and forbidden graphs and contains the proposed bound,
the flag types and flags used by each SDP block, and the rational matrix data.  The
argument \lean{<target>.lean} is the path at which the script should write the
generated Lean source file.  The output file contains the flag and
density generation commands of Part~1, the matrices and basis-element vectors
of Part~2, the target-lifting lemma when Part~3 is needed, and the complete
proof of the bound in Part~4.

For example, applying the compiler to the Mantel certificate with
the following command
\begin{lstlisting}[language={}]
python flagmatic_to_lean.py gen-skeleton <mantel-cert>.json Mantel.lean --namespace Mantel --force
\end{lstlisting}
generates \lean{Mantel.lean} from the Mantel certificate.  The optional
\lean{--namespace Mantel} argument selects the namespace of the generated
declarations, and \lean{--force} permits the script to replace an existing
target file.  Before generating the file, one can run the separate diagnostic
command
\begin{lstlisting}[language={}]
python flagmatic_to_lean.py inspect <mantel-cert>.json
\end{lstlisting}
to check that the certificate can be parsed and that its flags can be matched
to Lean's enumeration.  This command does not generate a proof file; it prints
the resulting identifier mapping and the Part~1 generation commands.  We now
explain the generated Lean file one part at a time before showing an abridged
assembled output.

\paragraph{Generating the flags and density data (Part~1).}
This is where the reflection layer of \Cref{sec:reflection} enters the
compiler, making the finite bookkeeping involving flags and densities
executable and checkable.  Part~1 defines the certificate's forbidden graph as a
\lean{Sym2Graph} term and invokes the \lean{generate\_forbid\_free\_*}
commands of \Cref{sec:reflection-generation} on it.  The commands choose
between two enumeration procedures according to the forbidden graph: they use
size-inductive clique pruning when the forbidden graph is complete, and the
general ordinary-subgraph filtering procedure otherwise.  During elaboration,
they enumerate the flags that avoid the forbidden graph, with isomorphic
flags appearing only once.  For each enumerated flag, the commands emit a concrete
reflection-layer constant, either a \lean{Sym2EmptyTypedFlag} at the empty flag type
or a \lean{Sym2Flag} at a nonempty flag type, together with its decoded
specification-level \lean{Flag} and \lean{FlagAlgebra} basis element.  As in
\Cref{sec:reflection-generation}, the commands also prove the completeness
lemma stating that the emitted flags exhaust the flags avoiding the forbidden
graph.

After emitting these flags, Part 1 generates
the rewrite theorems, which are used later
in Parts 3 and 4 of the proof.  Concretely, for each emitted $\sigma$-flag
$F$, Part~1 evaluates the downward normalizing
factor~$q_\sigma(F)$ from \Cref{sec:formal-downward}.
For each relevant pair of $\sigma$-flags $F_1,F_2$ and host $\sigma$-flag $G$, it also
evaluates the pair density
\(\lean{flagDensity}_2\,F_1\,F_2\,G\), which supplies the coefficient of $[G]$
in the product expansion of $[F_1]\cdot[F_2]$.  When the target must be
lifted in Part~3, the generated file similarly contains proved evaluations of
the required single-flag densities.  The adequacy theorems from
\Cref{sec:reflection} turn all these reflection-layer evaluations into
specification-level statements about the decoded flags.  The multiplication
generator then combines the pair-density identities with the general
multiplication theorem and the completeness lemma for the generated flag set
to derive the required product expansions.

\paragraph{Certifying the matrices (Part~2).}
With the flags and their rewrite theorems in place, Part~2 discharges the
positive-semidefiniteness obligations.  For each SDP block, the certificate
stores rational matrices $Q'_t$ and $R_t$ with
$M_t=R_tQ'_tR_t^\top$, where $R_t$ may be rectangular and $Q'_t$ is a smaller
symmetric matrix. Flagmatic's user guide~\cite{vaughan2012flagmaticguide}
explains the format: $Q'_t$ is
chosen to be positive definite, which implies the positive semidefiniteness of $M_t$.

Our verification takes these data through a different, equally elementary
witness.  We state one Lean lemma for every PSD obligation, based on an exact
rational $LDL^\top$ factorization:
\begin{lstlisting}
theorem posSemidef_real_of_LDLt {n : (*@$\mathbb{N}$@*)}
    {M L : Matrix (Fin n) (Fin n) (*@$\mathbb{Q}$@*)} {d : Fin n -> (*@$\,\mathbb{Q}$@*)}
    (hd : forall (*@$\,$@*)i, 0 <= (*@$\,$@*)d i) (hM : M = L * Matrix.diagonal d * L(*@$^\top$@*))
    : (ratMatrixToReal M).PosSemidef := ...
\end{lstlisting}
The lemma says that if a rational matrix $M$ factors as
$L\operatorname{diag}(d)L^\top$ with every entry of $d$ non-negative, then the
entrywise real cast of $M$, written \lean{ratMatrixToReal M}, is positive
semidefinite; the $\mathbb{Q}$-to-$\mathbb{R}$ casting happens inside the
lemma.  To use it, the compiler script reconstructs $M_t$ from $Q'_t$ and
$R_t$ in exact rational arithmetic, computes a factorization
$M_t=L_t\operatorname{diag}(d_t)L_t^\top$ with $L_t$ square and unit lower
triangular, and emits $M_t$, $L_t$, and $d_t$ in the generated Lean file.
Lean treats the emitted factorization only as a candidate witness: the
\lean{psd\_real\_ldlt} tactic must prove the lemma's two hypotheses.  The
non-negativity of every entry of $d_t$ is checked with \lean{fin\_cases} and
\lean{norm\_num}.  For the factorization equality, both sides are concrete
finite matrices over $\mathbb{Q}$, so \lean{decide +kernel} evaluates their
decidable equality and has the Lean kernel certify that their rational
entries agree.  No eigenvalue or floating-point calculation enters this
check.  Lean does not separately compare $M_t$ with the original $Q'_t$ and
$R_t$ fields of the certificate; that reconstruction is performed by the
unverified compiler script.  An incorrect $LDL^\top$ witness is therefore
rejected, while fidelity to the source certificate is a separate concern
discussed in \Cref{sec:flagmatic-trust}.  In either case, the emitted matrix
can contribute to the final theorem only if both its PSD proof and the
subsequent algebraic calculation are accepted by Lean.

\paragraph{Lifting the target (Part~3).}
The first two parts provide the algebraic data used by the certificate.  Part~3
aligns the target with that data when necessary.  The certificate identities
are expressed at the host size $N$, so when the target graph $F$ itself has
$N$ vertices, its basis element $[F]$ already lives in that basis and nothing
needs to be done.  When $F$ is smaller, the compiler proves an auxiliary
identity that expands $[F]$ over the $N$-vertex $\mathcal{H}$-free
$\emptyset$-flags,
\[
  [F]
  \;=\;
  \sum_{\substack{G \in \mathcal{F}^{\emptyset}_{N} \\ G\ \mathcal{H}\text{-free}}}
  \den{F}{G}\,[G]
  \qquad\text{modulo the \(\mathcal{H}\)-free condition,}
\]
with coefficients $\den{F}{G}$ given by the single-flag densities
generated in Part~1.

All certificates considered here have a singleton forbidden family
$\mathcal{H}=\{H\}$.  The generated proof invokes the
\lean{flag\_expand\_hfree N K} tactic, where $N$ is the host size and $K$ is the
reflection-layer \lean{Sym2Graph} encoding of $H$.  The qualifier on the
identity above matters, because our development works in the ambient algebra.  In
the built-in \(H\)-free algebra (\Cref{rem:builtin-hfree}), the restricted sum
would simply be the expansion identity~\eqref{eq:expansion-identity}, all of
whose flags are \(H\)-free there.  In the ambient algebra, however, the flags containing
\(H\) are nonzero basis elements, so dropping their terms cannot give a
literal equality; the auxiliary identity instead says that every empty-type
positive homomorphism satisfying the \(H\)-freeness condition of
\Cref{def:ensemble-semantic-order} gives both sides the same value.
Part~1 selected the $H$-free flags by a concrete
computation on $K$.  Bridging the two requires proving that the underlying
graph of every flag that computation discarded contains $H$ as a subgraph; the
$H$-free condition then forces those flags' terms to vanish, which justifies
expanding over exactly the flags Part~1 emitted.  From $N$ and $K$ alone, the
tactic establishes this bridge itself.

\paragraph{Assembling the main theorem (Part~4).}
Part~4 assembles the outputs of the preceding parts and realizes the
chain of inequalities from \Cref{sec:flagmatic-background}.  The certificate identity
$[F]+\sum_tQ_t=c\cdot\mathbf{1}-r$ supplies its middle equality, while the
non-negativity of the $Q_t$ and of $r$ supplies its two inequalities.  In the
Mantel example, there is one $Q_t$, namely
$Q=\bigl\llbracket\tfrac12((\fKtwobull)-(\fbarKtwobull))^2\bigr\rrbracket_\sigma$,
and $r=\tfrac13\cdot\fbarPthree$.  The generated Lean proof does not
define $r$ explicitly; it appears as the coefficientwise difference left after
the two sides are expanded in a common basis.

The first step is to prove
$[F]\leq_H[F]+\sum_tQ_t$.  It uses the following general lemma:
\begin{lstlisting}
theorem forbidLEWith_add_QuadraticForm {k n : (*@$\mathbb{N}$@*)} {(*@$\sigma$@*) : FlagType (Fin k)}
    {C : ForbidCondition} {f g : FlagAlgebra (*@$\emptyt$@*)}
    (M : Matrix (Fin n) (Fin n) (*@$\mathbb{R}$@*)) (hM : M.PosSemidef) (v : FlagAlgebraVec (*@$\sigma$@*) n)
    : forbidLEWith C f g -> (*@$\,$@*) forbidLEWith C f (g + (*@$\llbracket$@*)flagQuadraticForm M v(*@$\rrbracket_0$@*)) := ...
\end{lstlisting}
Here $C$ is a predicate encoding the forbidden-subgraph set $\mathcal{H}$; for a problem
forbidding a single graph $H$, it is \lean{forbiddenCondition H}, the
$H$-freeness condition of \Cref{def:ensemble-semantic-order} used throughout
this section.  Thus,
\lean{forbidLEWith C f g} is the implementation-level form of $f\leq_H g$.
The lemma says that if $f\leq_H g$, then the inequality remains true after the
non-negative downward quadratic form determined by the PSD matrix $M$ and flag
vector $v$ is added to the right-hand side; \lean{hM} is the proof of positive
semidefiniteness generated in Part~2.
The proof starts from the reflexive inequality $[F]\leq_H[F]$ and applies this
lemma once for each SDP block.  After all blocks have been added, its conclusion is precisely
$[F]\leq_H[F]+\sum_tQ_t$.

It remains to prove $[F]+\sum_tQ_t\leq_H c\cdot\mathbf{1}$.  The compiler checks
this comparison by rewriting both sides in the basis of $H$-free flags at the
host size $N$.  When necessary, the Part~3 lemma first expands $[F]$ to that
basis.
Expanding each quadratic form produces downward images of products of basis
elements; the \lean{reduce\_downward\_flagmul} tactic rewrites them using the multiplication
and downward identities generated in Part~1.  The
\lean{expand\_one\_hfree\_at} tactic similarly expands $\mathbf{1}$ in the same
basis, dispatching on the forbidden graph exactly as
\lean{flag\_expand\_hfree} does in Part~3.  Finally,
the \lean{flagsum\_ac\_sort\_rhs\_pipeline} tactic collects the coefficients of matching
basis elements.

After these rewrites, the difference between the right- and left-hand sides is
exactly the remainder $r$ from the certificate calculation.  The tactic
\lean{flag\_nonneg} checks that its coefficients are non-negative and then uses
the non-negativity of every basis element to conclude $0\leq_H r$.  In this way,
the normalization and the final non-negativity check together verify the
conceptual equality $[F]+\sum_tQ_t=c\cdot\mathbf{1}-r$ and the last inequality
$c\cdot\mathbf{1}-r\leq_H c\cdot\mathbf{1}$.

\paragraph{The compiler script.}
The compiler script performs the certificate-specific assembly.  It parses
the target, bound, forbidden graph, host size, and SDP blocks from the
Flagmatic file.  For each graph or flag represented by a compact string, it
finds the isomorphic entry in Lean's canonical enumeration; the resulting
index selects the corresponding \lean{FlagAlgebra\_n\_k\_m\_i} constant.  The
script also reconstructs the rational matrices, computes candidate exact
$LDL^\top$ witnesses, and determines whether the target requires lifting.

From these data, the script emits one Lean source file organized according to
the four parts above: the generation commands of Part~1, the matrix and flag
vector declarations and PSD proofs of Part~2, the optional target expansion
of Part~3, and the final theorem of Part~4.  Certificates change the generated
data and the number of SDP blocks, but not this proof structure.

The generated file is a proof proposal that Lean checks independently,
including all its intermediate propositions and its final bound.  The source
certificate itself is not an input to Lean, however, so this check establishes
the validity of the generated theorem rather than the fidelity of the
translation.  The next subsection makes this boundary precise.

\subsection{Trust Model}
\label{sec:flagmatic-trust}

There are two distinct questions of trust: is the theorem that Lean accepts
valid, and is it the statement the certificate was meant to yield?

For the first question, neither the SDP solver nor the compiler script is
trusted for the logical validity of a theorem that Lean accepts.  The
complete generated files used here introduce neither axiom declarations nor
\lean{sorry}, so every claim is carried by a proof term Lean checks: the
emitted matrices must be proved positive semidefinite, the certificate
calculation must normalize to non-negative coefficients, and the final
theorem must typecheck.  Subject to Lean's trusted base, erroneous data
cannot make an invalid proof term valid.

For the second question, the compiler script is trusted for \emph{translation
fidelity}: it parses the
certificate, chooses the corresponding entry in Lean's canonical enumeration
for each certificate flag, reconstructs $M_t=R_tQ'_tR_t^\top$, and selects the
target, bound, and forbidden graph.
These choices are visible in the generated declarations and theorem statement,
but Lean does not prove that they faithfully reproduce the source certificate.
A translation error could therefore produce a valid theorem different from
the one intended.  Guarding against this does not require trusting the
script; it requires reading the final theorem statement.  In this sense the
script marks an auditability boundary, rather than a soundness gap in the
generated Lean theorem.

The proof obligations also use two evaluation mechanisms.  Five of the seven
generated files evaluated in \Cref{sec:results}, including the Erd\H{o}s
pentagon, discharge every finite
computation, from the exact matrix equalities of the PSD proofs to the
enumerations and density calculations, by \lean{decide +kernel}; the kernel
checks all of it, and no compiled-evaluation axiom appears.  In the $K_5$-free
and $C_5$-free edge-density certificates, kernel reduction of the generated
computations was too slow to be practical, so those two files use
\lean{native\_decide} and additionally trust Lean's native compiler and
runtime.  Finally, the custom
tactics are unverified metaprograms whose dependence on elaborated expression
shapes affects robustness and portability.  They construct proof terms that
Lean checks, so a tactic failure leaves the proof incomplete rather than
establishing an invalid theorem.

\subsection{Evaluation}
\label{sec:results}\label{sec:flagmatic-scenarios}

We evaluated the compiler on seven certificates spanning host sizes
$N\in\{3,4,5\}$, targets at and below the host size, between one and four SDP
blocks, and four
forbidden graphs ($K_3$, $K_4$, $K_5$, and $C_5$).  For each certificate,
the compiler produces a complete Lean source file.  All seven files compile
and contain neither \lean{sorry} nor generated axiom declarations.  No
certificate coefficient is transcribed into Lean by hand.  All but the
$K_5$- and $C_5$-free edge-density files are checked entirely by
\lean{decide +kernel}; those two use
\lean{native\_decide}.
Table~\ref{tab:flagmatic-scenarios} summarizes the seven cases;
\Cref{sec:compile-times} reports their compilation times and peak memory
under both evaluation modes.

\begin{table}
  \centering
\caption{Certificates carried end-to-end through the compiler.  \(H\) is the
  forbidden graph, \(n_F\) the number of vertices of the target graph \(F\),
  \(N\) the host size, and \(T\) the number of SDP blocks.}
  \label{tab:flagmatic-scenarios}
  \begin{tabular}{lccccc}
    \toprule
    Case & $H$ & $n_F$ & $N$ & $T$ & Bound \\
    \midrule
    Mantel (edge density)               & $K_3$ & 2 & 3 & 1 & $1/2$    \\
    $P_3$ density                       & $K_3$ & 3 & 3 & 1 & $3/4$    \\
    $C_4$ density                       & $K_3$ & 4 & 4 & 2 & $3/8$    \\
    Edge density                        & $K_4$ & 2 & 4 & 2 & $2/3$    \\
    Erd\H{o}s pentagon ($C_5$ density)  & $K_3$ & 5 & 5 & 3 & $24/625$ \\
    Edge density                        & $K_5$ & 2 & 5 & 4 & $3/4$    \\
    Edge density                        & $C_5$ & 2 & 5 & 4 & $1/2$    \\
    \bottomrule
  \end{tabular}
\end{table}

The cases play complementary roles.  Mantel is a hand-checkable baseline: its
single $2\times2$ SDP block is the explicit square from
\Cref{sec:flagmatic-background}.  Here the dimension of a block is the
dimension of its certificate matrix, equivalently the length of its associated
flag vector.  The $C_4$-density certificate has two blocks, with matrices of
dimensions $4\times4$ and $3\times3$, while the Erd\H{o}s pentagon certificate
has three blocks, with matrices of dimensions $8\times8$, $6\times6$, and
$5\times5$.  The $K_5$- and $C_5$-free edge-density cases each use four blocks,
all with~$8\times8$ matrices.  These certificates are Flagmatic's output as
produced, with no attempt to minimize the number of blocks or the matrix
dimensions; more compact certificates for the same bounds may exist, but the
unoptimized output already compiles end to end.  Together, these examples
exercise targets both
at and below the host size, and increasingly large multi-block normalizations.

The $C_5$-free edge-density case shows that the forbidden graph need not be
complete; the generated file proves that graphs containing no $C_5$ subgraph
have edge density at most $1/2$.

The generated certificate argument establishes only the upper-bound direction.
For Mantel and the Erd\H{o}s pentagon,
the matching lower bounds are proved directly in Lean using the underlying
formalization (\Cref{sec:lowerbounds}), completing the
Tur\'an densities $\tfrac12$ and $\tfrac{24}{625}$ of Examples~\ref{ex:mantel}
and~\ref{ex:pentagon}.

\paragraph{Limitations.}
\label{sec:flagmatic-open}
We point out three limitations of the current compiler.
First, this evaluation establishes coverage of the seven cases above, not
completeness of the compiler.  Even a certificate within the compiler's
intended scope may expose a normalization shape not handled by the custom
tactics.  Such failures are safe: an unsupported or malformed input is
rejected or yields an explicitly incomplete skeleton containing \lean{sorry},
and a tactic failure leaves the proof incomplete, so neither case produces an
unchecked theorem.
Second, fully kernel-checked compilation does not yet scale.  For the
$K_5$-free and $C_5$-free edge-density certificates, we attempted to check the
generated files entirely by \lean{decide +kernel}; the runs did not finish in
an acceptable time and were abandoned, which is why those two files use
\lean{native\_decide} (\Cref{sec:flagmatic-trust}).
Third, as explained in the trust model, we do not verify translation fidelity:
Lean does not check that the generated declarations faithfully transcribe the
source certificate.  For each of the seven evaluated certificates, however,
the compiler fills the fixed proof skeleton without manual editing, and Lean
checks every proof obligation in the generated file under the trust model of
\Cref{sec:flagmatic-trust}.  Thus, the fidelity of the generated theorem to the
source certificate remains a limitation, while the logical validity of the
generated theorem itself is checked by Lean.

\section{Using the Formalization Beyond Certificate Compilation}
\label{sec:lowerbounds}

The certificate compiler of \Cref{sec:flagmatic} is one application of our
development, not its full scope.  Underlying the compiler is a reusable
formalization of flag-algebra expressions, positive homomorphisms and their
semantics, and the passage from finite graphs to limiting densities.  These
components can also be used directly, without first translating an argument
into a certificate.

This section illustrates three such uses.  The Mantel development gives a
direct flag-algebra proof and shows what it is like to work with the
formalization without the compiler.  The matching lower bounds for Mantel's
theorem and the Erd\H{o}s pentagon theorem combine explicit finite
constructions with the convergence of normalized extremal numbers, completing
the corresponding exact Tur\'an-density results.  Finally, two inequalities of
Goodman show that the same framework supports flag-algebra arguments outside
the certificate problems considered in \Cref{sec:flagmatic}.

\paragraph{Mantel's theorem.}
\Cref{ex:mantel} states $\tdensity{K_2}{K_3}=\tfrac12$.  The Lean development
records it in four statements:
\begin{lstlisting}
theorem Mantel_flag_bound : K2 (*@$\leq$@*) (1 / 2 : (*@$\mathbb{R}$@*)) (*@$\leansmul$@*) 1 + K3 := ...

theorem Mantel_flag_bound' : (*@$\forall$@*) ((*@$\phi$@*) : PositiveHom (*@$\emptyt$@*)),
    (*@$\phi$@*) K3 = 0 (*@$\to$@*) (*@$\phi$@*) K2 (*@$\leq$@*) 1 / 2 := ...

lemma extremal_density_K3_ge (n : (*@$\mathbb{N}$@*)) (hn2 : n (*@$\geq$@*) 2)
    : (extremalNumber n (completeGraph (Fin 3)) / n.choose 2 : (*@$\mathbb{R}$@*)) (*@$\geq$@*) 1 / 2 := ...

theorem Mantel_Turan : turanDensity (completeGraph (Fin 3)) = 1 / 2 := ...
\end{lstlisting}
Here \lean{K2} and \lean{K3} denote the basis elements $[K_2]$ and $[K_3]$.
The compiler of \Cref{sec:flagmatic} already produces the upper bound in
Mantel's theorem.  We include this separate development for a different
reason: it shows what a proof looks like when one works directly with the
flag-algebra formalization instead of supplying a certificate to the compiler.
Accordingly, this proof uses neither the compiler nor the
\lean{forbidLE} interface of \Cref{sec:forbidden}.

The direct upper-bound argument takes place in the ambient algebra and keeps
the triangle contribution explicitly:
\[
  K_2\leq\tfrac12\cdot\mathbf{1}+K_3.
\]
This inequality is formalized as \lean{Mantel\_flag\_bound}.  If a positive
homomorphism $\phi$ represents a triangle-free limit, then
$\phi([K_3])=0$, and evaluating the displayed inequality gives
$\phi([K_2])\leq\tfrac12$.  This pointwise statement is recorded as
\lean{Mantel\_flag\_bound'}, and the remaining argument proceeds from it.

The matching lower bound comes from the balanced complete bipartite graph
$K_{\lfloor n/2\rfloor,\lceil n/2\rceil}$ and is formalized as
\lean{extremal\_density\_K3\_ge}.  This graph is triangle-free:
among any three vertices, two lie in the same part and are therefore
non-adjacent.  Its $\lfloor n^2/4\rfloor$ edges make up at least half of the
$\binom{n}{2}$ vertex pairs.  Here \lean{extremalNumber} denotes the maximum
number of edges in a $K_3$-free graph on $n$ vertices.  It agrees with
$\mathrm{ex}(n,\,K_2;\,K_3)$ because the induced copies of $K_2$ are precisely
the edges.  When $n=2k$, Lean uses
\lean{completeBipartiteGraph (Fin k) (Fin k)} as the witness and compares its
$k^2$ edges with $\tfrac12\binom{2k}{2}$.  The case of odd $n$ follows by
comparison with the next even graph size, using the fact that the normalized
extremal numbers are non-increasing.

Finally, the construction and the upper bound are combined in
\lean{Mantel\_Turan}, using the compactness argument of
\Cref{lem:semantic-bound} directly rather than the machinery of
\Cref{sec:forbidden}.  If the normalized extremal
numbers exceeded $\tfrac12+\varepsilon$ for infinitely many values of $n$, we
could choose extremal triangle-free graphs of those sizes and pass, by
\Cref{thm:convergent-hom}, to a convergent subsequence.  Its limit would be a
positive homomorphism $\phi$ with $\phi([K_3])=0$ and
$\phi([K_2])\geq\tfrac12+\varepsilon$, contradicting
\lean{Mantel\_flag\_bound'}.  The construction gives the lower bound
$\tfrac12$, so the normalized extremal numbers converge to $\tfrac12$, which
is the statement of \lean{Mantel\_Turan}.

\paragraph{The Erd\H{o}s pentagon theorem.}
This example shows how an explicit construction and the passage from finite
graphs to limiting densities complement a compiler-generated upper bound.
The full result of \Cref{ex:pentagon},
\begin{lstlisting}
theorem ErdosPentagon_Turan : generalizedTuranDensity K3 C5 = 24/625 := ...
\end{lstlisting}
is the conjunction of the upper bound of \Cref{sec:results} and the lower
bound
\begin{lstlisting}
theorem ErdosPentagon_Turan_lowerBound : generalizedTuranDensity K3 C5 (*@$\geq$@*) 24/625 := ...
\end{lstlisting}
proved by an explicit construction in Lean.  We first describe the
construction and the analytic steps mathematically, and then match them with
the formal development.

Since, by the convergence theorem
\lean{tendsto\_\allowbreak generalizedTuranDensity} of \Cref{sec:forbidden},
the normalized extremal numbers $\mathrm{ex}(n,\,C_5;\,K_3)/\binom{n}{5}$
converge to $\tdensity{C_5}{K_3}$, it suffices to exhibit, for infinitely many
graph sizes, one triangle-free graph whose normalized count of induced copies of
$C_5$ is at least $\tfrac{24}{625}$.  The
witness is the \emph{$n$-fold blow-up} of $C_5$: replace each vertex of
$C_5$ by a class of $n$ vertices, and join two vertices of distinct
classes exactly when the corresponding vertices of $C_5$ are adjacent;
vertices
of a common class remain non-adjacent.  The blow-up is triangle-free:
adjacent vertices lie in distinct classes, so the classes of a hypothetical
triangle would form a triangle in $C_5$ itself, which has none.  It also
contains many copies of $C_5$: picking one vertex from each of the five
classes
yields five vertices that inherit exactly the adjacencies of $C_5$, so each
of the $n^5$ choice functions induces a copy of $C_5$.  These copies are
induced, as the density
convention of \Cref{sec:intro} requires.  (In a triangle-free graph, the
distinction disappears anyway: every $C_5$ subgraph is induced, since any
chord, an edge joining two non-consecutive vertices of the cycle, would
create a triangle.)  The blow-up is therefore a
triangle-free graph on $5n$ vertices with at least $n^5$ induced copies
of~$C_5$,
so $\mathrm{ex}(5n,\,C_5;\,K_3)\geq n^5$, and normalizing gives, for every
$n\geq1$,
\[
  \frac{\mathrm{ex}(5n,\,C_5;\,K_3)}{\binom{5n}{5}}
  \;\geq\;
  \frac{n^5}{\binom{5n}{5}}
  \;\geq\;
  \frac{n^5}{(5n)^5 / 5!}
  \;=\;
  \frac{24}{625}.
\]
The bound thus holds along the subsequence of graph sizes $5n$, and taking the
limit yields $\tdensity{C_5}{K_3}\geq\tfrac{24}{625}$.

The Lean development mirrors each step of this outline:
\begin{lstlisting}
def blowUp {V : Type} [Fintype V] (G : SimpleGraph V) (n : (*@$\mathbb{N}$@*))
    : SimpleGraph (V (*@$\times$@*) Fin n) :=
  { Adj v w := G.Adj v.1 w.1
    symm v w := by apply G.symm }

theorem blowUp_K3_free {m : (*@$\mathbb{N}$@*)} {G : SimpleGraph (Fin m)}
    (n : (*@$\mathbb{N}$@*)) (hfree : K3.Free G) : K3.Free (blowUp G n) := ...

lemma subgraphCount_blowUp_C5_ge (n : (*@$\mathbb{N}$@*))
    : subgraphCount C5 (blowUp C5 n) (*@$\geq$@*) n ^ 5 := ...

theorem generalizedExtremalNumber_K3_C5_ge (n : (*@$\mathbb{N}$@*))
    : generalizedExtremalNumber (5 * n) K3 C5 (*@$\geq$@*) n ^ 5 := ...
\end{lstlisting}
The \lean{blowUp G n} is the $n$-fold blow-up of an arbitrary base graph $G$: its
vertex type is the product type \lean{V $\times$ Fin n}, whose elements are
the pairs $(v,i)$ of a vertex \lean{v : V} and an index \lean{i : Fin n}, so
the class of a base vertex $v$ consists of the pairs with first coordinate
$v$, and the \lean{Adj} field compares first coordinates through $G$.
Vertices of a common class are non-adjacent because $G$ has no loops, and two
classes over adjacent base vertices are completely joined.
The \lean{blowUp\_K3\_free} theorem states the triangle-freeness property.  Its proof
transports a hypothetical triangle back to the base graph: composing a copy
of $K_3$ in the blow-up with the first projection gives a copy of $K_3$ in
$G$, where injectivity survives the projection because vertices of a common
class are non-adjacent.  Instantiating $G$ with $C_5$, whose own
triangle-freeness is a finite check, shows that the blow-up of $C_5$
is triangle-free.  The
\lean{subgraphCount\_blowUp\_C5\_ge} lemma
expresses the counting property; here
\lean{subgraphCount C5 G} is the number of induced subgraphs of \lean{G}
isomorphic to $C_5$, matching the induced-copy convention of
\Cref{sec:intro}.  The proof converts the one-vertex-per-class recipe into an
injection from \lean{(Fin 5 $\to$ Fin n)} into the induced copies of $C_5$
in the blow-up: a choice function $g$ is sent to the subgraph induced on the five
vertices $(i, g(i))$ for \lean{i : Fin 5}, and distinct choice functions give
distinct subgraphs.

The preceding two results give a triangle-free graph on $5n$ vertices
containing at least $n^5$ induced copies of $C_5$.  The theorem
\lean{generalizedExtremalNumber\_K3\_C5\_ge} records the resulting finite bound
\[
  \mathrm{ex}(5n,\,C_5;\,K_3)\geq n^5.
\]
There is one additional representation step in Lean.  The natural vertex set
of \lean{blowUp C5 n} is the product \lean{Fin 5 $\times$ Fin n}, while
\lean{generalizedExtremalNumber} takes its maximum over graphs on
\lean{Fin (5 * n)}.  The graph \lean{blowUp\_fin} simply reindexes the blow-up
on this latter vertex set; it is proved isomorphic to \lean{blowUp} and is used
as the witness.  Normalizing the finite bound above and passing to the limit
along the graph sizes $5n$ then gives
\lean{ErdosPentagon\_Turan\_lowerBound}.

\paragraph{Direct Goodman inequalities.}
The final examples use the core flag-algebra infrastructure directly for two
inequalities outside the Tur\'an-density certificate problems considered in
\Cref{sec:flagmatic}.  Both are due to Goodman~\cite{goodman1959}.
The first is the algebraic lower bound
\[
  K_3 \;\geq\; K_2\cdot(2K_2-\mathbf{1}),
\]
which bounds the triangle density from below in terms of the edge density:
under a positive homomorphism with edge density $e$, the triangle density is
at least $e(2e-1)$, a quantity that becomes positive as soon as $e$ exceeds
the Mantel threshold $\tfrac12$.  The proof combines the size-three expansion
of $K_2$ with the Cauchy--Schwarz inequality for the downward operator,
\[
  \llbracket f\,g\rrbracket_\sigma^{\,2}
  \;\leq\;
  \llbracket f^2\rrbracket_\sigma\cdot\llbracket g^2\rrbracket_\sigma,
\]
which is a separate theorem in Razborov's paper~\cite{razborov2007flag} and
is proved in the development in this general form as
\lean{Cauchy\_\allowbreak Schwarz\_\allowbreak inequality}; here it is
applied to the edge flag labeled at one endpoint.  The resulting inequality is
formalized as
\lean{Goodman\_\allowbreak bound\_\allowbreak on\_\allowbreak triangle\_\allowbreak density}.

The second result is
\[
  \overline{K_3} + K_3 \;\geq\; \tfrac14\cdot\mathbf{1},
\]
where $\overline{K_3}$ is the empty graph on three vertices: asymptotically,
at least a quarter of all
vertex triples induce a triangle or an independent set.  Reading a graph and
its complement as the two color classes of a $2$-edge-colored complete graph
turns this into Goodman's classical bound on the number of monochromatic
triangles; the proof reuses the downward square certificate of Mantel's
theorem.  It is formalized as
\lean{Goodman\_theorem\_on\_Ramsey\_multiplicity}.

Together, the examples in this section show that certificate compilation is
only one use of the formalization.  The same framework supports direct
flag-algebra derivations, explicit finite constructions and their limiting
arguments, and inequalities beyond maximizing the induced density of one graph
while forbidding another as a subgraph.

\section{Engineering Obstacles}
\label{sec:obstacles}

The specification layer of \Cref{sec:abstract} makes explicit several proof
obligations that ordinary combinatorial notation hides.  A construction
defined on a concrete labeled graph must be proved invariant under
isomorphism before it can be lifted to flags.  Because a flag's size is part
of its Lean type, identities relating flags of different sizes, such as the
expansion identity, require additional type-level bookkeeping.  Reusing a
theorem about a tuple of flags can create a similar problem after the tuple
is reordered: mathematically identical tuples need not have types that Lean
identifies.  Finally, finite counting depends on a chosen enumeration, and
the same choice must be used consistently throughout a proof.  These
obligations do not change the underlying mathematics, but they recur often
enough to impose a substantial proof-engineering cost.

Although these difficulties often appear together, their causes differ.
Invariance under isomorphism is the well-definedness condition required to
define an operation on a quotient.  The size and tuple-reordering issues
reflect a limitation of Lean's underlying type theory: mathematically equal
types need not be definitionally equal, and therefore may require explicit
casts.  The enumeration issue, by contrast, is not primarily a
consequence of the underlying type theory.  It arises from elaboration and
typeclass inference, which may synthesize different enumerations at different
occurrences.  The following subsections treat these obstacles separately.

Counting identities present one further challenge: two density expressions
may count the same configurations using different finite representations.
Our most reliable proof method is to identify the two finite sets, construct
an explicit bijection between them, and deduce equality of their
cardinalities.  The final subsection explains this recurring pattern.

\subsection{Computing with Graphs up to Isomorphism}

As described in \Cref{sec:flagtypes}, a flag is an isomorphism class of
concrete labeled graphs.  Thus, two isomorphic labeled graphs $G$ and $G'$ may
use different vertex names and be represented by different Lean values, but
they denote the same flag.  Representing flags by isomorphism classes is
mathematically natural because subsequent statements on flags should not depend on
arbitrary vertex names.  At the same time, this quotient representation
creates a recurring well-definedness obligation whenever a construction is
first described using a concrete representative.

Suppose such a construction assigns a value $A(G)$ to a labeled graph $G$.
To define the corresponding value of the flag represented by $G$, one must
prove
\[
  G \simeq_f G' \quad\Longrightarrow\quad A(G)=A(G').
\]
Only then is the rule $[G]\mapsto A(G)$ independent of the representative.
Lean packages this step as \lean{Quotient.lift}.  The use of
\lean{Quotient.lift} is itself short; the substantive obligation is proving
that the representative-level construction is invariant under isomorphism.

This pattern occurs when lifting densities and graph operators to flags and
also at the reflection boundary.  For example, decoding an executable
edge-set representation must yield the same abstract flag for every
isomorphic concrete encoding.  Rather than exposing the two representations in
every later theorem, we prove this compatibility once and let subsequent
proofs operate on the resulting flag.  The same issue appears one level higher
when flag vectors are quotiented by the zero-space relations to form
\lean{FlagAlgebra}: a construction on vectors must respect those relations
before it can be lifted to the algebra.

One could avoid these compatibility proofs by choosing a canonical
representative for every isomorphism class.  That choice would make the
representation easier to compute with, but it would introduce
canonicalization details into the statements and proofs of algebraic lemmas.
We instead keep the quotient specification close to Razborov's mathematics
and concentrate the representation-specific work at the boundaries where a
construction is lifted to equivalence classes.

\subsection{Comparing Flags of Different Sizes}

In Lean, a flag's size is part of its type: the $\sigma$-flags of size $\ell$
form the type \lean{Flag $\sigma$ (Fin $\ell$)}, which varies with $\ell$.  A
statement comparing flags of different sizes must therefore cross between
different types, a crossing that paper mathematics performs silently.  The
situation arises throughout the development: the expansion
identity~\eqref{eq:expansion-identity} equates a size-$m$ flag with a
linear combination of size-$n$ flags, and the product formula of
Equation~\eqref{eq:mul} must be proved independent of its auxiliary size.
This subsection explains how such cross-size statements are made and proved.

Consider the expansion identity~\eqref{eq:expansion-identity}.  Its
left-hand side involves a flag of type \lean{Flag $\sigma$ (Fin $m$)},
whereas the flags on the right have type \lean{Flag $\sigma$ (Fin $n$)}.  A
term of the first type cannot be used where Lean expects a term of the
second, and an ordinary equality can only compare terms having the same
type.  Thus, the identity is not merely hard to prove; as written, it is not
a well-typed statement.

The first formalization task is therefore to construct explicitly
the common ambient space that paper mathematics leaves implicit.
Mathematically, we regard every fixed-size family
$\mathcal{F}^{\sigma}_{\ell}$ as part of the all-size family
$\mathcal{F}^{\sigma}$, and hence regard combinations of flags of different
sizes as vectors in the same space.  In Lean, however, the types
\lean{Flag $\sigma$ (Fin $\ell$)} vary with $\ell$ and do not come with such a
common ambient type.  We therefore use the dependent sum
\[
  \text{\lean{FinFlag}}~\sigma
  \;=\;
  \mathop{\Sigma}_{\ell : \mathbb{N}}
    \text{\lean{Flag}}~\sigma~(\text{\lean{Fin}}~\ell).
\]
A size-$m$ $\sigma$-flag $F$ is therefore stored as the tagged pair
$\langle m,F\rangle$, while a size-$n$ $\sigma$-flag $G$ is stored as
$\langle n,G\rangle$.  The second components have different types, but
both tagged pairs belong to the single type \lean{FinFlag}~$\sigma$.
Recall that \lean{FlagVector}~$\sigma$ is the vector space of finitely
supported real-valued functions on \lean{FinFlag}~$\sigma$, or equivalently
of finite formal real linear combinations of its elements.  Thus, the two
tagged $\sigma$-flags index basis vectors in the same vector space, and
equations mixing flags of different sizes can be stated in that common space.

The packaging provides a common home, not an identification: when their size
tags differ, $\langle m,F\rangle$ and $\langle n,G\rangle$ remain distinct
elements of \lean{FinFlag}~$\sigma$ and index distinct basis vectors.
Identifying them is the quotient's job, and it needs no further argument:
the difference between the two sides of the expansion identity~\eqref{eq:expansion-identity} 
is one of the generators of the zero space, so the identity holds
in the flag algebra by definition.  Later cross-size comparisons, including
the independence of the
product formula from its auxiliary size, reason modulo this relation.

\subsection{Comparing Reordered Flag Tuples}
\label{sec:reordered-flag-tuples}

The multi-flag density is invariant under reordering its pattern flags.  One
instance is the symmetry
\(
  \den{F_0,F_1}{G}=\den{F_1,F_0}{G}
\),
which is what makes flag-algebra multiplication commutative, since the
product formula~\eqref{eq:mul} expands a product of two flags into a sum
weighted by exactly these pair densities.  On paper, the invariance is proved
once, in full generality: for any number of pattern flags, reordering the
tuple does not change its density in a host.  The two-flag symmetry is then
read off by taking the permutation that swaps the two positions, and this
instantiation is not counted as a proof step.  The formalization takes the
same route, and the friction appears exactly in the instantiation, a step
with no counterpart on paper.

Recall from \Cref{sec:formal-subflag-densities} that a tuple of pattern
flags is a dependent function: an element of
\lean{FlagList}~$\sigma$~\lean{t Vl} assigns 
a $\sigma$-flag on the vertex type \lean{Vl i} to each index \lean{i : Fin t}, and \lean{flagListDensity}
computes the multi-flag density of such a tuple in a host flag.  Both the
entries and their types vary with the index, and the vertex-type family
\lean{Vl} is itself an argument of the tuple's type.  Reordering is defined
by precomposition:
\begin{lstlisting}
abbrev Perm (t : (*@$\mathbb{N}$@*)) := Fin t (*@$\simeq$@*) Fin t

def listTypePermute {t : (*@$\mathbb{N}$@*)} (Vl : Fin t (*@$\to$@*) Type) ((*@$\pi$@*) : Perm t)
    : Fin t (*@$\to$@*) Type :=
  fun i => Vl ((*@$\pi$@*) i)

def FlagList.permute {T : Type} {(*@$\sigma$@*) : FlagType T} {t : (*@$\mathbb{N}$@*)} {Vl : Fin t (*@$\to$@*) Type}
    (Fl : FlagList (*@$\sigma$@*) t Vl) ((*@$\pi$@*) : Perm t)
    : FlagList (*@$\sigma$@*) t (listTypePermute Vl (*@$\pi$@*)) :=
  fun i => Fl ((*@$\pi$@*) i)
\end{lstlisting}
A permutation \lean{$\pi$ : Perm t} is a bijection of the index set
\lean{Fin t}.  The permuted tuple \lean{Fl.permute $\pi$} places
\lean{Fl ($\pi$ i)} at position \lean{i}, so the entry types are reordered
along with the entries: the vertex-type family of the result is
\lean{listTypePermute Vl $\pi$}, the reordered family
\lean{fun i => Vl ($\pi$ i)}, and it appears in the type of the result.  The
invariance is then proved once, for every length \lean{t} and every
permutation; the symmetry to be derived from it is the pair case:
\begin{lstlisting}
theorem flagDensity_permute {T W : Type} {(*@$\sigma$@*) : FlagType T} {t : (*@$\mathbb{N}$@*)} {Vl : Fin t (*@$\to$@*) Type}
    (Fl : FlagList (*@$\sigma$@*) t Vl) (G : Flag (*@$\sigma$@*) W) ((*@$\pi$@*) : Perm t)
    : flagListDensity Fl G = flagListDensity (Fl.permute (*@$\pi$@*)) G := ...

theorem flagPairDensity_comm {T U(*@$_0$@*) U(*@$_1$@*) W : Type} {(*@$\sigma$@*) : FlagType T} 
    (F(*@$_0$@*) : Flag (*@$\sigma$@*) U(*@$_0$@*)) (F(*@$_1$@*) : Flag (*@$\sigma$@*) U(*@$_1$@*)) (G : Flag (*@$\sigma$@*) W)
    : flagDensity(*@$_2$@*) F(*@$_0$@*) F(*@$_1$@*) G = flagDensity(*@$_2$@*) F(*@$_1$@*) F(*@$_0$@*) G := ...
\end{lstlisting}
Here \lean{flagDensity}$_2$ is the pair specialization of
\lean{flagListDensity}, so the second theorem states precisely the two-flag
symmetry introduced at the beginning of this subsection.

To derive the second statement from the first, suppose the $\sigma$-flags
$F_0$ and $F_1$ have vertex types $U_0$ and $U_1$, and let \lean{Vl} be the
vertex-type family with
\lean{Vl 0 = U}$_0$ and \lean{Vl 1 = U}$_1$, so that the two-entry flag tuple
\lean{[F}$_0$\lean{, F}$_1$\lean{]}$^{\mathtt{f}}$, using the
\lean{FlagList} notation introduced in
\Cref{sec:formal-subflag-densities}, is an element of
\lean{FlagList}~$\sigma$~\lean{2 Vl}.  Take $\pi$ to be the permutation of
\lean{Fin 2} with $\pi(0)=1$ and $\pi(1)=0$.  Then
\lean{flagDensity\_permute} equates the density of
\lean{[F}$_0$\lean{, F}$_1$\lean{]}$^{\mathtt{f}}$ with that of its permuted
tuple; the permuted tuple places $F_1$ at position $0$ and $F_0$ at
position $1$, and its vertex-type family is
\lean{listTypePermute Vl $\pi$}.  The theorem \lean{flagPairDensity\_comm} speaks instead
of the directly constructed pair
\lean{[F}$_1$\lean{, F}$_0$\lean{]}$^{\mathtt{f}}$, an element of
\lean{FlagList}~$\sigma$~\lean{2 Vl'} for the vertex-type family \lean{Vl'}
with \lean{Vl' 0 = U}$_1$ and \lean{Vl' 1 = U}$_0$.  The families 
\lean{listTypePermute Vl $\pi$} and \lean{Vl'} are pointwise equal, since both send $0$
to $U_1$ and $1$ to $U_0$, but Lean need not reduce the two expressions to
the same term.  The family is an argument of the tuple's type, so the
permuted tuple and \lean{[F}$_1$\lean{, F}$_0$\lean{]}$^{\mathtt{f}}$
inhabit types, \lean{(FlagList}~$\sigma$~\lean{2 (listTypePermute Vl $\pi$))}
and \lean{FlagList}~$\sigma$~\lean{2 Vl'}, that are not definitionally the
same.  The permuted tuple and the directly constructed pair therefore cannot
yet be compared by ordinary equality.  The formal proof must first establish
equality of their vertex-type families and then use that equality to compare
the tuples.  This type-alignment step has no counterpart on paper, where the
two expressions simply denote the same reordered pair.

The bridging lemma \lean{flagListDensity\_HEq\_eq} packages these
requirements: if the two vertex-type families are equal and the two tuples
agree after their types are aligned, then the two densities coincide.
\begin{lstlisting}
theorem flagListDensity_HEq_eq {T W : Type} {(*@$\sigma$@*) : FlagType T} {t : (*@$\mathbb{N}$@*)} {Vl(*@$_0$@*) Vl(*@$_1$@*) : Fin t (*@$\to$@*) Type}
    {Fl(*@$_0$@*) : FlagList (*@$\sigma$@*) t Vl(*@$_0$@*)} {Fl(*@$_1$@*) : FlagList (*@$\sigma$@*) t Vl(*@$_1$@*)}
    (h_Vl_eq : Vl(*@$_0$@*) = Vl(*@$_1$@*)) (h_HEq : HEq Fl(*@$_0$@*) Fl(*@$_1$@*)) (G : Flag (*@$\sigma$@*) W)
    : flagListDensity Fl(*@$_0$@*) G = flagListDensity Fl(*@$_1$@*) G := ...
\end{lstlisting}
In this application, \lean{Vl}$_0$ is
\lean{listTypePermute Vl $\pi$} and \lean{Vl}$_1$ is \lean{Vl'}, so
\lean{h\_Vl\_eq} states
\[
  \lean{listTypePermute Vl $\pi$ = Vl'},
\]
proved by checking the two indices of \lean{Fin 2}.  The hypothesis
\lean{h\_HEq} expresses that the permuted tuple and the directly constructed
tuple \lean{[F}$_1$\lean{, F}$_0$\lean{]}$^{\mathtt{f}}$ are equal even
though Lean initially assigns them different types.  Concretely, both tuples
place $F_1$ at position $0$ and $F_0$ at position $1$.

The role of \lean{h\_Vl\_eq} is to align these two tuple types.  Once their
vertex-type families have been identified, \lean{h\_HEq} can be used as an
ordinary equality, allowing one tuple to be substituted for the other.  The
lemma \lean{flagListDensity\_HEq\_eq} then concludes that their densities are
equal.  Thus, the formal argument recovers the mathematically immediate fact
that permuting \lean{[F}$_0$\lean{, F}$_1$\lean{]}$^{\mathtt{f}}$ produces
the same pair as directly constructing
\lean{[F}$_1$\lean{, F}$_0$\lean{]}$^{\mathtt{f}}$.

\subsection{Controlling Finite Enumerations}
\label{sec:fintype}

Lean distinguishes a proof that a type is finite from a mechanism for
enumerating its elements.  An instance of \lean{Finite T} provides only the
former: it asserts that $T$ has finitely many elements but gives no way to
visit them.  An instance of \lean{Fintype T} provides the latter:
\begin{lstlisting}
class Fintype (T : Type) where
  elems : Finset T
  complete : (*@$\forall$@*) x : T, x (*@$\in$@*) elems
\end{lstlisting}
The field \lean{elems} is a \lean{Finset T}, a concrete finite set of
elements of \lean{T}, and \lean{complete} proves that every element of
\lean{T} belongs to it.  The set \lean{elems} carries no mathematically
significant order, but it allows generic operations to count the elements,
inspect every candidate, and form finite sums and products.  Thus,
\lean{Fintype} is the interface needed by counting algorithms,
whereas \lean{Finite} records only the underlying finiteness fact.  Every
\lean{Fintype} instance implies \lean{Finite}.  In the other direction,
\lean{Fintype.ofFinite} produces a \lean{Fintype} instance from a proof of
finiteness by invoking Lean's choice axiom \lean{Classical.choice}.  The
resulting instance is noncomputable: its \lean{elems} exists as a term that
proofs can mention, and every lemma about \lean{Fintype} applies to it, but
the definition carries no executable code, so Lean cannot evaluate
\lean{elems} to list the elements of $T$.  This costs nothing when a proof relates cardinalities to one another by
lemmas, which is how the counting below proceeds.  What a noncomputable
instance rules out is evaluation.  Proving a concrete numeric value of
\lean{Fintype.card T} by computation, or enumerating the elements
executably, requires a computable \lean{Fintype} instance; the reflection
layer (\Cref{sec:reflection}) builds its enumerations in that form.

This distinction becomes especially important when formalizing cardinalities of finite sets
and finite counting arguments.  A set $S=\{x:T\mid P(x)\}$ is represented as a
predicate of type \lean{Set T}; when
we count its elements, the relevant type is the subtype of terms $x:T$
satisfying $P(x)$.  Lean can often infer that this subtype is \lean{Finite}
without having a particular \lean{Fintype} enumeration available, but the
cardinality function \lean{Fintype.card} and the lemmas about it, such as
\lean{Finset.card\_univ} and \lean{Fintype.card\_congr}, require a
\lean{Fintype} instance.  The elaborator does synthesize \lean{Fintype}
instances automatically; the difficulty is not obtaining an instance but
obtaining the same one every time.  In a development with multiple
overlapping constructions (quotient types, embedded subgraphs,
partial-function spaces), the same type can receive different \lean{Fintype}
instances from different elaboration paths.  As sets, their \lean{elems}
fields cannot disagree, since \lean{complete} forces each to contain every
element of the type; what can differ is the term denoting that set.  When
the two \lean{elems} terms do not reduce to a common term definitionally,
the kernel rejects the goal even though the instances are provably equal.

The formalization responds in two ways.  The first is preventive: the
moment a finite set is named, the intended instance is chosen and bound
with \lean{let}, so that every later occurrence elaborates against the same
binding rather than a fresh search:
\begin{lstlisting}
-- S0 and S1 are the two finite sets being counted.
-- Choose and name the Fintype enumeration used for each one.
let hS0 : Fintype S0 := Fintype.ofFinite S0
let hS1 : Fintype S1 := Fintype.ofFinite S1
-- h_iso_S0_S1 is a bijection between the two sets.
-- The @ notation exposes the Fintype arguments that Lean normally inserts.
have card_eq : @Fintype.card S0 hS0 = @Fintype.card S1 hS1 :=
  @Fintype.card_congr S0 S1 hS0 hS1 h_iso_S0_S1
-- Rewrite the toFinset counts in the goal to Fintype cardinalities.
have count_eq : S0.toFinset.card = S1.toFinset.card := by
  -- These two toFinset terms also use hS0 and hS1, respectively.
  rw [Set.toFinset_card, Set.toFinset_card]
  exact card_eq
\end{lstlisting}
The displayed \lean{@} notation makes explicit where the two named instances
enter the proof: \lean{hS0} supplies the enumeration used by
\lean{Fintype.card S0}, and \lean{hS1} supplies the one used by
\lean{Fintype.card S1}; both are also arguments of
\lean{Fintype.card\_congr}.  The subsequent \lean{S0.toFinset} and
\lean{S1.toFinset} expressions use the same instances implicitly, as do the
two applications of \lean{Set.toFinset\_card}.  Thus,
\lean{Set.toFinset\_card} rewrites the two concrete counts to exactly the
cardinality terms appearing in \lean{card\_eq}.  Without the local bindings,
Lean could synthesize those hidden arguments independently at the different
occurrences, producing terms that are propositionally equal but do not match
for rewriting.

This pattern appears in \lean{flagDensity\allowbreak\_permute} from
\Cref{sec:reordered-flag-tuples}.  There, \lean{S0} and \lean{S1} are the
sets of tuples of pairwise-disjoint induced copies that realize \lean{Fl} and
\lean{Fl.permute $\pi$}, respectively.  The mathematical step is the
reindexing bijection \lean{h\_iso\_S0\_S1}; the named \lean{Fintype}
instances ensure that the cardinality equality obtained from this bijection
rewrites the counts occurring in the density definition.  The flag-operator
proofs use the same pattern for the isomorphism set associated with a fixed
underlying graph.

The second response repairs mismatches that do occur.  Once the
mathematical steps of such a proof are done, what remains is a goal of a
characteristic shape: an equality whose two sides print identically and
differ only in the inferred instances inside them.  The
tactic \lean{congr!} closes such goals; it decomposes the equality into
congruence subgoals and discharges the instance mismatches, which is sound
because any two \lean{Fintype} instances of the same type are provably
equal (\lean{Fintype T} is a subsingleton).  The bridging lemma
\lean{flagListDensity\allowbreak\_HEq\allowbreak\_eq} of
\Cref{sec:reordered-flag-tuples} is finished exactly this way: it carries
a \lean{FintypeList} instance for each vertex-type family, supplying a
\lean{Fintype} instance for every component type.  Once the families and the
tuples have been identified, the two sides of the goal differ only in the
\lean{Fintype} instances supplied for the component types; \lean{congr!}
resolves these remaining differences.

These explicit \lean{Fintype} choices are needed only where finite sets are
converted to \lean{Finset}s or their cardinalities are compared.  Once the
required counting identity has been proved, the subsequent algebraic argument
uses that identity without referring to the chosen instances.

\subsection{Proving Counting Identities by Bijections}
\label{sec:proof-style}

Explicit bijections became an unexpectedly prevalent proof pattern in the
specification layer.  We did not set out to make bijective arguments the
default: we
expected many of these equalities to follow from existing counting lemmas or
symbolic simplification, with explicit equivalences needed only occasionally.
In practice, however, unfolding and simplifying the definitions repeatedly
exposed the same structure: the two sides are cardinalities of finite sets of
configurations presented in different ways.  Our most reliable proof method
is to identify the two sets, construct an explicit bijection between them, and
use the bijection to prove equality of their cardinalities.  The bijection
carries the main mathematical content of the proof.

Our experience proving the density chain rules (\Cref{lem:chain-rule}) made
this pattern particularly clear.  A chain rule identifies two sampling procedures: choosing the required
subflags directly in the host, or making the same choice in stages through an
intermediate $\sigma$-flag.  On paper, these procedures are readily seen as two
descriptions of the same experiment.  In Lean, after unfolding the
definitions, the proof came down to a bijection between the finite sets of
choices made by the two procedures.  This bijection yielded equality of the
corresponding counts and hence the chain rule.

The same method proves permutation invariance, invariance under inserting an
empty flag, the downward-operator counting identities, and the adequacy of
the computational subgraph representation.  Although these obligations arise
in different parts of the formalization, each reduces to an explicit
correspondence between two finite sets.  The bijection contains the
mathematical argument, while the instance discipline of
\Cref{sec:fintype} handles the resulting cardinalities.  We do not claim
that bijections are the only possible approach, but their effectiveness
across these different obligations made them our default method for finite
counting identities.

\section{A Meta-Theory of Ensemble Semantics}
\label{sec:metatheory}

The definition of \lean{forbidLE} in \Cref{sec:forbidden} was a deliberate
departure from the usual presentation of constrained flag algebras.  Rather
than building a new algebra for each forbidden family, it retains the same
ambient algebra and imposes the forbidden-graph condition only when
interpreting its elements.  It considers only unlabeled positive
homomorphisms satisfying that condition, applies the corresponding
random-extension measure to each one, and requires the desired inequality to
hold with probability one under that measure.
This engineering choice exposed a mathematical problem in its own right:
when does this \emph{ensemble semantics} validate exactly the same typed
inequalities as a flag algebra in which the constraint is built in from the
beginning?  This led us to develop a separate mathematical meta-theory.
This section gives the definitions and main results needed to state that answer;
a forthcoming paper will present the full theory and its applications.

\paragraph{The constrained positive-homomorphism space.}
Let \(\mathcal G\) be the class of all finite simple graphs, let~\(\mathcal K\subseteq\mathcal G\) be a hereditary class, and fix a flag type
\(\sigma\).  Here \emph{hereditary} means closed under taking induced
subgraphs: if~\(G\in\mathcal K\) and \(S\subseteq V(G)\), then
\(G[S]\in\mathcal K\).  In particular, the \(H\)-free classes used earlier are
hereditary, because taking an induced subgraph cannot create a new
subgraph isomorphic to \(H\).

To make the underlying graph class explicit, we write
\(\FlagAlg{\sigma}[\mathcal G]\) for the ambient
\(\sigma\)-typed flag algebra and
\(\FlagAlg{\sigma}[\mathcal K]\) for the algebra constructed using only graphs
in \(\mathcal K\).  Only \(\sigma\)-flags whose underlying unlabeled graphs
belong to~\(\mathcal K\) occur in this constrained algebra.  Because
\(\mathcal K\) is hereditary, the remaining \(\sigma\)-flags, whose underlying
graphs lie outside \(\mathcal K\), span an ideal,\footnote{Here an ideal \(I\)
is a linear subspace of the flag algebra that is also closed under
multiplication by arbitrary algebra elements.  Thus, \(I\) is closed under
addition and real scalar multiplication, and if \(x\in I\) and \(a\) belongs
to the flag algebra, then \(ax\in I\).  This property makes the quotient by
\(I\) an algebra in which every element of \(I\) is identified with zero.} and the usual constrained
algebra is canonically isomorphic to the quotient of the ambient algebra by
this ideal.  We therefore have a canonical quotient map
\[
  q_\sigma:
  \FlagAlg{\sigma}[\mathcal G]
  \twoheadrightarrow
  \FlagAlg{\sigma}[\mathcal K]
\]
that sends every \(\sigma\)-flag outside \(\mathcal K\) to zero and every
remaining \(\sigma\)-flag to its counterpart in the constrained algebra.

Let
\[
  X_\sigma
  :=
  \mathrm{Hom}^{+}\bigl(\FlagAlg{\sigma}[\mathcal G],\R\bigr)
\]
be the compact space of positive homomorphisms on the ambient
algebra, with the topology of pointwise convergence on \(\sigma\)-flags.
Pullback along \(q_\sigma\) embeds the positive homomorphisms of the constrained
algebra into \(X_\sigma\).  Its image is
\begin{equation}\label{eq:Q-sigma}
  Q_\sigma
  :=
  \left\{
    \psi\circ q_\sigma
    \;\middle|\;
    \psi\in
    \mathrm{Hom}^{+}\bigl(\FlagAlg{\sigma}[\mathcal K],\R\bigr)
  \right\}
  \subseteq X_\sigma.
\end{equation}
Thus, \(Q_\sigma\) is the space of all \(\sigma\)-typed positive homomorphisms
admitted by the constrained quotient.

There is also an intrinsic description:
\begin{equation}\label{eq:Q-sigma-intrinsic}
  Q_\sigma
  =
  \left\{
    \chi\in X_\sigma
    \;\middle|\;
    \chi([F])=0
    \text{ for every \(\sigma\)-flag \(F\) whose underlying graph lies outside
    \(\mathcal K\)}
  \right\}.
\end{equation}
The inclusion from left to right is immediate.  If
\(\chi=\psi\circ q_\sigma\) belongs to \(Q_\sigma\), then
\(q_\sigma([F])=0\) for every \(\sigma\)-flag \(F\) outside \(\mathcal K\),
and hence \(\chi([F])=0\).

For the reverse inclusion, suppose that \(\chi\in X_\sigma\) vanishes on
every \(\sigma\)-flag outside \(\mathcal K\).  Here heredity ensures that the
span of these flags is an ideal.  Every \(\sigma\)-flag \(F\) appearing in
the product of a \(\sigma\)-flag \(F'\) outside \(\mathcal K\) with another
\(\sigma\)-flag \(F''\) contains \(F'\) as an induced subflag.  The flag \(F\)
must therefore also lie outside \(\mathcal K\), since otherwise heredity would
imply that the underlying graph of \(F'\) lies in \(\mathcal K\).  Thus, the
span of the flags outside \(\mathcal K\) is precisely the ideal killed by
\(q_\sigma\).  The homomorphism \(\chi\) vanishes on this ideal and therefore
factors through \(q_\sigma\) as a positive homomorphism on the constrained
algebra.  Hence \(\chi\in Q_\sigma\), proving the reverse inclusion.
Equation~\eqref{eq:Q-sigma-intrinsic} also shows
that \(Q_\sigma\) is closed in \(X_\sigma\), since it is an intersection of
zero sets of continuous evaluation maps.  We write
\(Q_0:=Q_{\emptyset}\) for the corresponding space of constrained empty-type
positive homomorphisms.

\paragraph{Random extensions and the root-planting space.}
Assume that \(\sigma\) is \emph{non-degenerate} for \(\mathcal K\): that is,
there is some \(\phi_0\in Q_0\) for which
\(\phi_0(\langle\sigma\rangle_0)>0\).  By Razborov's random-extension result
recalled in \Cref{sec:background-downward}, every such \(\phi_0\) determines a
unique probability measure \(\Ext_\sigma(\phi_0)\) on \(X_\sigma\).  It is
obtained by choosing a uniformly random occurrence of \(\sigma\) in each
member of a graph sequence converging to \(\phi_0\) and passing to the
limit.  For a probability
measure~\(\mu\) on~\(X_\sigma\), its
support \(\operatorname{supp}(\mu)\) consists of those points
\(\chi\in X_\sigma\) for which every open neighborhood of \(\chi\) has positive
\(\mu\)-measure.
Define
\begin{equation}\label{eq:S-sigma}
  S_\sigma
  :=
  \overline{
    \bigcup_{\substack{\phi_0\in Q_0\\
                        \phi_0(\langle\sigma\rangle_0)>0}}
    \operatorname{supp}\bigl(\Ext_\sigma(\phi_0)\bigr)
  }
  \subseteq X_\sigma.
\end{equation}
Thus, to form \(S_\sigma\), we first collect the supports of
\(\Ext_\sigma(\phi_0)\) over all \(\phi_0\in Q_0\) for which \(\sigma\) has
positive density.  We then take the closure in \(X_\sigma\), adding every point
that can be approximated arbitrarily closely by points in those supports.

The support of each random-extension measure is contained in the constrained
space:
\[
  \operatorname{supp}\bigl(\Ext_\sigma(\phi_0)\bigr)
  \subseteq Q_\sigma
  \qquad
  \bigl(\phi_0\in Q_0,
        \phi_0(\langle\sigma\rangle_0)>0\bigr).
\]
Indeed, let \(F\) be a \(\sigma\)-flag outside \(\mathcal K\).  Applying the
downward operator to \([F]\) produces a non-negative multiple of
\([F|_\emptyset]\), whose underlying graph also lies outside \(\mathcal K\).
Every \(\phi_0\in Q_0\) evaluates \([F|_\emptyset]\) to zero.  Therefore, if~\(\phi\sim\Ext_\sigma(\phi_0)\), the random-extension
identity in Equation~\eqref{eq:ensemble} gives
\[
  \mathbb{E}[\phi([F])]=0.
\]
Since every positive homomorphism \(\phi'\) satisfies \(\phi'([F])\geq0\), it
follows that for \(\phi\sim\Ext_\sigma(\phi_0)\), we have
\(\phi([F])=0\) with probability one.  Applying this argument to
the countable collection of \(\sigma\)-flags outside \(\mathcal K\) shows that a positive
homomorphism sampled from \(\Ext_\sigma(\phi_0)\) belongs to \(Q_\sigma\) with
probability one.  Since \(Q_\sigma\) is closed, Equation~\eqref{eq:S-sigma} yields the
fundamental inclusion
\begin{equation}\label{eq:S-subset-Q}
  S_\sigma\subseteq Q_\sigma.
\end{equation}
We say that \((\mathcal G,\mathcal K,\sigma)\) is
\emph{root-plantable}\footnote{The vertices distinguished by the flag type
\(\sigma\) are often called the \emph{roots}, which explains the first part of
the term \emph{root-plantable}.  A finite \(\sigma\)-flag \((G,\theta)\)
records a profile of \(\sigma\)-flag densities relative to a specified embedding
\(\theta:\sigma\to G\), whose image consists of the labeled vertices.  To
\emph{plant} this profile in a larger graph \(H\in\mathcal K\), without
specifying labels in advance, is to produce many embeddings
\(\widehat\theta:\sigma\to H\) whose profiles of \(\sigma\)-flag densities approximate the
profile specified by \((G,\theta)\).  Here ``many'' means that these embeddings
form a fraction bounded away from zero of all embeddings of \(\sigma\) in
\(H\).  Consequently, a uniformly chosen occurrence of \(\sigma\) reproduces
the prescribed profile with non-vanishing probability.  The equality
\(S_\sigma=Q_\sigma\) says that every constrained \(\sigma\)-typed positive
homomorphism can be approximated by finite constructions of this kind.} when
\(S_\sigma=Q_\sigma\).

\begin{definition}[Quotient and ensemble orders]\label{def:two-constrained-orders}
For \(f,g\in\FlagAlg{\sigma}[\mathcal G]\), define
\[
  f\leq^{\mathrm{quot}}_{\mathcal K,\sigma}g
  \quad\Longleftrightarrow\quad
  \chi(f)\leq\chi(g)
  \quad\text{for every }\chi\in Q_\sigma,
\]
and
\[
\begin{aligned}
  f\leq^{\mathrm{ens}}_{\mathcal K,\sigma}g
  \quad\Longleftrightarrow\quad
  \Ext_\sigma(\phi_0)
  \bigl(\{\chi\in X_\sigma\mid \chi(f)\leq\chi(g)\}\bigr)=1
  \quad\text{for every }\phi_0\in Q_0\text{ with }
  \phi_0(\langle\sigma\rangle_0)>0.
\end{aligned}
\]
\end{definition}

The first is the ordinary semantic order in the constrained quotient algebra:
equivalently, every positive homomorphism
\(\psi:\FlagAlg{\sigma}[\mathcal K]\to\R\) satisfies
\(\psi(q_\sigma(f))\leq\psi(q_\sigma(g))\).  The second quantifies over each
\(\phi_0\in Q_0\) on its right-hand side.  For such a
\(\phi_0\), the measure \(\Ext_\sigma(\phi_0)\) is the probability distribution
on \(X_\sigma\) obtained by taking a sequence of finite graphs converging
to \(\phi_0\), choosing uniformly at random an embedding of \(\sigma\) into
each of them, using the embedded vertices as labels, and passing to the
limit.  Thus, the second display requires
that the random positive homomorphism \(\chi\sim\Ext_\sigma(\phi_0)\) satisfy
\(\chi(f)\leq\chi(g)\) with probability one.
For the \(H\)-free class, consisting of graphs containing no subgraph
isomorphic to \(H\),
\(\lean{forbidLE H f g}\) in our Lean formalization implements this ensemble order.  The corresponding
notions of non-negativity are \(0\leq^{\mathrm{quot}}_{\mathcal K,\sigma}f\) and
\(0\leq^{\mathrm{ens}}_{\mathcal K,\sigma}f\).

\begin{theorem}[Support-closure criterion]\label{thm:support-closure-summary}
Let \(\mathcal K\) be hereditary and let \(\sigma\) be non-degenerate for
\(\mathcal K\).  For every \(f,g\in\FlagAlg{\sigma}[\mathcal G]\),
\[
  f\leq^{\mathrm{quot}}_{\mathcal K,\sigma}g
  \quad\Longrightarrow\quad
  f\leq^{\mathrm{ens}}_{\mathcal K,\sigma}g.
\]
Moreover, the two orders agree for every pair \((f,g)\) if and only if the class
is root-plantable at \(\sigma\):
\[
  \left(
    \forall f,g,\quad
    f\leq^{\mathrm{quot}}_{\mathcal K,\sigma}g
    \ \Longleftrightarrow\
    f\leq^{\mathrm{ens}}_{\mathcal K,\sigma}g
  \right)
  \quad\Longleftrightarrow\quad
  S_\sigma=Q_\sigma.
\]
\end{theorem}

The proof explains why the equality \(S_\sigma=Q_\sigma\) is exactly the
condition under which the two orders agree.  Evaluation of
\(g-f\) defines a continuous function
\(\chi\mapsto\chi(g-f)\) on \(X_\sigma\).  Almost-sure non-negativity under a
measure is equivalent to non-negativity on its support; continuity then shows
that non-negativity in the ensemble semantics is exactly non-negativity on
\(S_\sigma\).  Non-negativity in the quotient semantics is, by definition,
non-negativity on \(Q_\sigma\).
The inclusion in Equation~\eqref{eq:S-subset-Q} proves the forward implication, and
\(S_\sigma=Q_\sigma\) proves equivalence.

The necessity of root-plantability requires more than this observation.  If
\(S_\sigma\subsetneq Q_\sigma\), choose
\(\chi_*\in Q_\sigma\setminus S_\sigma\).  Topological separation gives a
continuous function that is strictly positive on \(S_\sigma\) and strictly
negative at \(\chi_*\).  The Stone--Weierstrass theorem then approximates this
function by evaluation of a flag-algebra element.  The resulting element is
non-negative in the ensemble semantics but not in the quotient semantics,
so the two orders
cannot agree for all elements.  Thus, \Cref{thm:support-closure-summary} is a completeness theorem:
quotient proofs are always sound for ensemble semantics, and they capture all
ensemble-valid typed inequalities exactly under root-plantability.

\paragraph{When root-plantability holds.}
The criterion in \Cref{thm:support-closure-summary} reduces semantic completeness to a structural problem about the
graph class.  Our main positive result gives a broad answer.  Given
\(G\in\mathcal K\), a vertex~\(v\in V(G)\), and a graph \(J\), write
\(G[v\mapsto J]\) for the graph obtained by replacing \(v\) by \(J\), keeping
the edges inside \(J\), and joining every vertex of \(J\) to every former
neighbor of \(v\).  Call \(\mathcal K\) \emph{blow-up-closed} if, for every
\(G\in\mathcal K\), \(v\in V(G)\), and \(N \in \mathbb{N}\), some graph \(J\) on exactly
\(N\) vertices satisfies \(G[v\mapsto J]\in\mathcal K\).  For a hereditary
class, this is equivalent to asking only for \(\lvert J\rvert\geq N\):
choose an induced \(N\)-vertex subgraph \(J'\) of \(J\), and observe that
\(G[v\mapsto J']\) is an induced subgraph of \(G[v\mapsto J]\).  The interior
graph \(J\) may depend on \(G\), \(v\), and \(N\); the definition requires the
existence of a suitable interior, not that every interior works.

\begin{theorem}[Blow-up closure implies root-plantability]
\label{thm:blowup-root-plantability-summary}
Every blow-up-closed hereditary graph class is root-plantable at every
nonempty flag type that is non-degenerate for the class.  The empty flag type is
root-plantable for every hereditary graph class, independently of blow-up
closure.  Consequently, quotient and ensemble orders agree at every
non-degenerate flag type in a blow-up-closed class.
\end{theorem}

The proof starts from a target \(\chi\in Q_\sigma\), represented by a sequence
of finite \(\sigma\)-flags whose underlying graphs lie in \(\mathcal K\).  For
each such \(\sigma\)-flag \((G,\theta)\), the proof applies the preceding
replacement construction once for each vertex of \(G\): the vertices
\(v\in V(G)\) are replaced one at a time by graphs \(J_v\), and at every
step blow-up closure supplies an interior for which the current graph
remains in \(\mathcal K\), so the final blow-up also lies in \(\mathcal K\).  The resulting blow-up has a natural projection back to \(G\): every
vertex of \(J_v\) is mapped to \(v\).  We call \(V(J_v)\), the inverse image of
\(v\) under this projection, the \emph{fiber over~\(v\)}.

The fibers over the labeled vertices \(\theta(i)\) are made large enough that a
uniformly chosen embedding of \(\sigma\) places each label \(i\) in the
corresponding fiber with probability bounded away from zero.  When a fixed
finite \(\sigma\)-flag is sampled from the blow-up, its sampled vertices lie in distinct
fibers with probability tending to one.  On this event, their adjacencies are
determined by the corresponding vertices of \(G\), so the internal structure
of the replacement graphs \(J_v\) does not affect the sampled \(\sigma\)-flag.

These estimates connect the finite blow-up construction to the definition of
\(S_\sigma\).  Let \(U\) be any open neighborhood of \(\chi\) in \(X_\sigma\).
Because \(X_\sigma\) has the topology of pointwise convergence on \(\sigma\)-flags,
membership in \(U\) can be ensured by approximating the values of finitely many
\(\sigma\)-flags within a prescribed tolerance.  Choose~\((G,\theta)\)
sufficiently far along the sequence representing \(\chi\), and then choose its
blow-up sufficiently large.  The construction ensures that a positive fraction
of the embeddings of \(\sigma\) into the blow-up produce \(\sigma\)-typed
density profiles lying in \(U\): the labeled vertices enter the designated
fibers with probability bounded away from zero, while the distinct-fiber
estimate preserves the required \(\sigma\)-flag densities.

Passing to a limit of these unlabeled blow-ups gives some \(\phi_0\in Q_0\) for
which
\[
  \Ext_\sigma(\phi_0)(U)>0.
\]
Therefore, \(U\) intersects \(\operatorname{supp}(\Ext_\sigma(\phi_0))\).  Since
every neighborhood \(U\) of \(\chi\) has this property, \(\chi\) lies in the
closure of the union of these supports, namely \(S_\sigma\).  Hence
\(Q_\sigma\subseteq S_\sigma\), while the opposite inclusion is
Equation~\eqref{eq:S-subset-Q}.

Independent blow-ups, in which every replacement graph \(J_v\) is an
independent set, cover all \(K_r\)-free classes
for \(r\geq3\), including triangle-free graphs.  Complete blow-ups, in which
every \(J_v\) is a clique, cover, for example, graphs whose connected
components are cliques.  Both are instances of
\Cref{thm:blowup-root-plantability-summary}.

\paragraph{Why the criterion is nontrivial.}
Heredity alone does not imply root-plantability.  Let \(\mathcal K\) be the
class of \(C_4\)-free graphs, those containing no \(C_4\) as a subgraph, and
let \(\sigma_1\) be the one-vertex type.
Every large \(C_4\)-free graph has \(o(n^2)\) edges, so a vertex chosen uniformly
at random has normalized degree tending to zero.  
Let~\(d\in\FlagAlg{\sigma_1}[\mathcal G]\) be the flag-algebra element represented by
the following two-vertex \(\sigma_1\)-flag:
\[
  d :=
  \begin{tikzpicture}[scale=0.75,baseline=-0.55ex]
    \draw[graph edge] (0,0) -- (1.2,0);
    \node[graph root] at (0,0) {};
    \node[graph vertex] at (1.2,0) {};
    \node[font=\scriptsize] at (0,-0.32) {\(1\)};
  \end{tikzpicture}
\]
The filled vertex carries label \(1\), and the open vertex is unlabeled.  Thus,
the density of \(d\) in a finite \(\sigma_1\)-flag is the normalized degree of
its labeled vertex: the probability that one additional uniformly sampled
vertex is adjacent to it.  The preceding observation implies that every
\(\chi\in S_{\sigma_1}\) satisfies \(\chi(d)=0\), so \(-d\) is non-negative
in the ensemble semantics.  On the other hand, arbitrarily large stars are
\(C_4\)-free: the star \(K_{1,t}\) consists of one vertex, its
\emph{center}, adjacent to \(t\) pairwise non-adjacent leaves, and contains
no cycle at all.  Label the center of \(K_{1,t}\) to obtain a
\(\sigma_1\)-flag, and let \(\chi_*\in Q_{\sigma_1}\) be the positive
homomorphism arising as the limit of this sequence of \(\sigma_1\)-flags.
Every leaf is adjacent to the center, so \(\chi_*(d)=1\).  Since
every point of \(S_{\sigma_1}\) evaluates \(d\) to zero,
\(\chi_*\notin S_{\sigma_1}\), and hence
\begin{equation}\label{eq:c4-strict-inclusion}
  S_{\sigma_1}\subsetneq Q_{\sigma_1}.
\end{equation}
Moreover, \(\chi_*(-d)=-1<0\), so \(-d\) fails to be non-negative in the
quotient semantics.
The example exhibits a genuine gap between the two typed orders: quotient
semantics permits a specially chosen vertex to carry the label, whereas
ensemble semantics obtains the labeled vertex by uniform random sampling.

The gap does not invalidate the final unlabeled density bounds of this
paper.  The
certificate pipeline consumes its typed inequalities only through the
downward operator, which preserves non-negativity in the ensemble
semantics
(\lean{downward\_\allowbreak preserve\_\allowbreak semanticCone},
\Cref{sec:formal-downward}).  Its conclusions land at the empty flag type,
which is root-plantable for every hereditary class by
\Cref{thm:blowup-root-plantability-summary}.  The final conversion
(\lean{generalizedTuranDensity\_\allowbreak le\_\allowbreak of\_\allowbreak forbidLE},
\Cref{sec:forbidden}) then turns these conclusions into classical
statements about Tur\'an densities.  Root-plantability is therefore not a
hypothesis of any of these results.  It concerns only the completeness of
the intermediate typed calculus, whether typed ensemble reasoning proves
exactly the typed inequalities of the constrained algebra.

\paragraph{Formalization and scope.}
The Lean formalization of the meta-theory presented in this section is included
in the public release of our flag-algebra development linked in
\Cref{sec:intro}, under the module \lean{LeanFlagAlgebras.MetaTheory}.  The
constrained space \(Q_\sigma\) in Equation~\eqref{eq:Q-sigma}, the root-planting space
\(S_\sigma\) in Equation~\eqref{eq:S-sigma}, and the condition
\(S_\sigma=Q_\sigma\) are formalized as \lean{Q$\sigma$}, \lean{S$\sigma$}, and
\lean{RootPlantable}, respectively.  The special cases
\(0\leq^{\mathrm{quot}}_{\mathcal K,\sigma}f\) and
\(0\leq^{\mathrm{ens}}_{\mathcal K,\sigma}f\) of the two orders in
\Cref{def:two-constrained-orders} are formalized as \lean{QuotientNonneg} and
\lean{EnsembleNonneg}, respectively; the general comparison \(f\leq g\) is
expressed by applying these predicates to \(g-f\).  At the theorem level,
\lean{quotient\_implies\_ensemble} formalizes the unconditional implication in
\Cref{thm:support-closure-summary}, while \lean{support\_criterion} formalizes
its characterization of when the two orders agree;
\lean{blowupClosed\_root\_plantable} and
\lean{heredClass\_emptyType\_rootPlantable} formalize, respectively, the
nonempty-type and empty-type claims of
\Cref{thm:blowup-root-plantability-summary}; and
\lean{c4free\_not\_rootPlantable}, together with
\lean{S$\sigma$\_subset\_Q$\sigma$}, formalizes the strict inclusion in
Equation~\eqref{eq:c4-strict-inclusion}.

These Lean files were generated by Claude Code from an already-developed
mathematical account and its theorem targets.  Lean checks the resulting proof
terms, and the released development builds without incomplete proofs.

This autoformalization was feasible because it extended, rather than rebuilt,
the manual formalization reported in the preceding sections.  The manual
library had already fixed the representations and semantic interfaces and
provided the ambient flag algebra, positive-homomorphism spaces, random
extensions, the downward operator, and the density and representation
theorems on which the meta-theory depends.  The generated meta-theory files
could therefore focus on the new constrained layer, beginning with the
forbidden ideal and constrained quotient, while reusing these established
foundations.

The present section isolates the core meta-theoretic consequences of
\lean{forbidLE}: quotient inequalities always imply ensemble inequalities,
root-plantability characterizes when the converse holds, and blow-up closure
is a broad sufficient condition for root-plantability.  A forthcoming paper
will give the complete mathematical and formal account and pursue three
extensions.  First, it will develop the finite planting criterion.  For every
fixed flag size and error tolerance, this criterion asks that each sufficiently
large finite \(\sigma\)-flag in the class admit a larger graph in the class
with a positive-density set of \(\sigma\)-embeddings that reproduce, to the
specified accuracy, all of its \(\sigma\)-flag densities up to that size.
This finite condition implies root-plantability.  Second, the forthcoming
work will study further graph classes for which quotient and ensemble
semantics agree or disagree.  A representative source of disagreement is a
pinning obstruction: every admissible random-extension measure forces the
density of a particular \(\sigma\)-flag to one value almost surely, while the
constrained quotient contains a positive homomorphism assigning it a
different value.  Third, it will develop relative ensemble semantics for
arbitrary sets of admissible unlabeled positive homomorphisms, allowing
non-hereditary limit constraints such as fixing the edge density.

\section{Open Design Questions}
\label{sec:unsure}

The formalization and end-to-end certificate proofs presented above are fully
checked in Lean.  This section instead concerns design questions whose broader
consequences cannot be settled by the present implementation alone: when
forbidden-graph constraints should enter the theory, which engineering costs
are specific to Lean, whether specification and computation should use
separate representations, and how far the current computational layer can
scale.  For each question, we distinguish what the present development
demonstrates from what would require alternative implementations or larger
computations to determine.

\paragraph{When should a forbidden-graph constraint enter the formalization?}
The predicate \lean{forbidLE} (\Cref{sec:forbidden}) keeps the ambient flag
algebra independent of the forbidden graph and imposes the constraint later,
when a semantic inequality is stated.  Razborov's universal-theory
presentation instead incorporates the constraint into the construction of
the flags, algebras, and positive homomorphisms.  Consequently, the two
approaches give different semantics to intermediate inequalities involving
labeled flags: the former restricts their interpretations, whereas the latter
removes inadmissible flags through the constrained algebra.  As
\Cref{thm:support-closure-summary} shows, this distinction disappears
precisely when the underlying hereditary class is root-plantable; in that
case, the two semantics validate the same typed inequalities.  Even without
root-plantability, every quotient inequality implies the corresponding
ensemble inequality.  This one-way implication is sufficient for
transferring the formal flag-algebra inequalities used here to the final
unlabeled Tur\'an bounds.

These semantic results do not by themselves determine which design is
preferable in a proof assistant.  Even when root-plantability ensures that the
two semantics validate the same inequalities, the implementations may have
different proof-engineering costs.  Deferring the constraint allows core
definitions and algebraic lemmas to be reused across forbidden graphs;
incorporating it at the outset may streamline arguments specific to a fixed
extremal problem.  Evaluating this tradeoff would require formalizing
comparable certificate proofs under both designs.

\paragraph{Which engineering obstacles are Lean-specific?}
\Cref{sec:obstacles} includes both obligations common to quotient-based
formalizations and difficulties tied more closely to Lean.  If flags are
represented as isomorphism classes, a construction first defined on labeled
graphs must be proved invariant under isomorphism before it can be applied to
flags.  The same well-definedness obligation would arise in any proof
assistant using this representation.  The remaining difficulties have causes
more closely tied to Lean.  Binding a particular \lean{Fintype}
instance is needed because Lean's elaborator may synthesize different instance
terms for different occurrences of the same carrier type.  The use of
\lean{HEq} reflects Lean's dependent type theory: pointwise-equal type families
need not be definitionally equal, so terms indexed by them cannot immediately
be compared using ordinary equality.
Other proof assistants handle finite enumeration, implicit arguments, and
dependent equality differently.  We therefore do not know which of these
Lean-specific proof-engineering techniques would be needed in a port to Rocq,
Agda, or Isabelle/HOL\@.  The underlying combinatorial identities and the
well-definedness obligations created by quotient representations would remain,
but their formal expression and proof cost could differ.

\paragraph{Should specification and computation use separate representations?}
Our development maintains two representations
(\Cref{sec:abstract,sec:reflection}).  The specification layer uses
quotients and definitions close to the mathematics, allowing the abstract
theory to ignore choices of graph representatives.  The reflection layer uses
the concrete \lean{Sym2Graph} representation to evaluate finite densities,
downward coefficients, and multiplication tables.  Adequacy theorems connect
the two, turning these evaluations into facts about the specification-level
definitions.  The density and multiplication facts used throughout our case
studies are obtained in this way.

This separation lets us keep the abstract development simple while retaining
an efficient executable representation.  It also creates a substantial
bridge of adequacy proofs.  One alternative is to state the theory directly
on a concrete, canonical representation: this would reduce the bridge, but
representation and canonicalization choices would become visible in the
statements and proofs of the abstract theory.  We have not carried this
alternative through comparable case studies, so we do not know which
architecture minimizes the total proof and computation cost.

\paragraph{Can the computational layer scale beyond five vertices?}
Our seven end-to-end certificates (\Cref{sec:results}) expand their
calculations over graphs on three, four, or five vertices.  At this scale,
the optimized enumeration of \Cref{sec:reflection-generation}, size-inductive,
deduplicated, and pruned under the forbidden graph, keeps the required
families of finite identities manageable.  For the Erd\H{o}s pentagon
certificate, whose calculations expand over five-vertex graphs, the three
flag types require 1800, 672, and 360 pair-density theorems, respectively.
Lean checks all 2832 identities by \lean{decide +kernel} and then uses the
resulting theorems in the complete certificate proof.  This confirms that
certificates at the five-vertex scale do go through the current reflection
and automation layers, though only just: the Erd\H{o}s pentagon is the
largest certificate carried entirely on the kernel-checked path.

Going substantially beyond five vertices remains a major engineering
challenge.  As the number of vertices grows, so do the families of flags,
the density and multiplication identities generated from their pairs, and the
expressions that the final algebraic proof must normalize.  The cost of
kernel-checked evaluation grows with them: already at five vertices, the two
largest certificates, with 10,332 and 8,172 pair-density theorems, could not
be checked by \lean{decide +kernel} in acceptable time and fall back to
\lean{native\_decide} (\Cref{tab:compile-times}), and eliminating that
fallback is part of the same scaling problem.  Scaling to larger
certificates will therefore likely require more efficient executable graph
representations and counting algorithms in the reflection layer, as well as
improvements to theorem generation and algebraic normalization.  This
ceiling is not unique to formal verification: certificate generation itself
faces the same combinatorial growth, and Flagmatic is in practice limited to
expanding its calculations in the basis of graphs on about seven or eight
vertices.  The formal pipeline reaches its
limit earlier, and we have not yet determined which of its costs dominates
or how much of the gap such optimizations can close.

\section{Related Work}
\label{sec:related}

\paragraph{Formalizations of combinatorics.}
Substantial combinatorial arguments have already been formalized in Lean.
Dillies and Mehta formalized Szemer\'edi's regularity lemma and the
qualitative form of Roth's theorem~\cite{dillies2022szemeredi}, and Mehta
the Kruskal--Katona theorem~\cite{mehta2022kruskal}.  Subercaseaux et
al.~\cite{subercaseaux2024hexagon} verified the empty-hexagon number against
a SAT certificate.  The collaborative PFR project~\cite{pfr2023} formalized
the Gowers--Green--Manners--Tao proof~\cite{gowers2025marton} of the
polynomial Freiman--Ruzsa conjecture in characteristic two.  These developments lean on heavy
computation or on a large supporting library.  A flag-algebra argument needs
both at once: it combines quotient and measure-theoretic semantics with
thousands of exact finite-density computations and an externally computed
semidefinite certificate.

\paragraph{Proof by reflection in proof assistants.}
Proof by reflection is a classical technique in the Rocq (previously called
Coq) community; Chlipala
gives an extended treatment~\cite{chlipala2013cpdt}, and Lean supports
verified computation through tactics such as \lean{decide} and
\lean{norm\_num}.  Cohen et al.~\cite{cohen2013refinements} develop
a systematic methodology for building efficiently computable representations
alongside their abstract counterparts in Rocq.  Our
\lean{Sym2Graph}/\lean{SimpleGraph} pair with its adequacy theorems follows
the same philosophy, adapted to Lean~4's type-class and instance-search
mechanism.

\paragraph{Computer-assisted flag algebra arguments.}
The Flagmatic software~\cite{vaughan2013flagmatic} automates the computation
of exact and approximate flag-algebra bounds, but it does not produce
proof-assistant proof terms; Razborov's survey~\cite{razborov2013flag}
organizes many early applications of the method.  Our work adds a
proof-producing reconstruction stage: certificate data proposes a
calculation, while Lean verifies the finite identities, the
positive-semidefinite matrices, and the final combinatorial statement.

\paragraph{A concurrent formalization of local flag algebras.}
Concurrently with and independently of our work, Davey, Hurley, de Joannis
de Verclos, Kang, and Volec formalized in Lean~4 the theory of \emph{local
flag algebras} introduced in their accompanying paper~\cite{davey2026local}
and applied in its companion~\cite{davey2026strong}.
The local calculus normalizes densities by the maximum degree $\Delta$
rather than by the number of vertices, so that the method still yields
nontrivial bounds for sparse graphs of small maximum degree.  Under this
normalization, no quotient is taken: their algebra is defined directly on
formal combinations of isomorphism classes, evaluation at a finite graph is
multiplicative only up to an error bounded by a constant multiple of
$1/\Delta$, and positivity is again captured by a semantic cone.  Their formalization
machine-certifies the new bounds of both papers, on pentagon counts in
triangle-free graphs~\cite{davey2026local} and on the strong chromatic
index~\cite{davey2026strong}.  Independently,
their design arrives at the reflection pattern of \Cref{sec:reflection}: a
noncomputable specification mirrored by a computable Boolean-adjacency
representation with adequacy lemmas.  Our development formalizes the
classical theory instead (the quotient algebra, positive homomorphisms,
and the ensemble semantics) and is organized as a reusable library and
proof-generating pipeline rather than around individual bounds, so the two
developments are complementary in both scope and design.

\paragraph{A concurrent formalization of computational flag algebras.}
A further concurrent development is Spiegel's \emph{computational flag
algebras} in Lean~\cite{spiegel2026computational}.  It is an independent
ongoing work on the Lean formalization of the flag algebras, which focuses
on carrying out flag-algebra inequality proofs fully and automatically
inside Lean, without relying on an external tool or solver.

\paragraph{SDP certificate verification.}
The verification of SDP certificates has been studied for sum-of-squares
proofs, which relax polynomial optimization problems to semidefinite
programs, as studied by Parrilo and
Sturmfels~\cite{parrilo2003minimizing}.  Harrison~\cite{harrison2007verifying}
verifies sum-of-squares certificates for polynomial inequalities in HOL
Light, with numerical methods finding the certificates and the proof
assistant checking them.  Morrison's \lean{sos} tactic, recently added
under the official \lean{leanprover} organization~\cite{morrison2026sos},
implements Harrison's procedure in Lean~4: it finds a certificate with an
external semidefinite solver, rounds it to rational form, and verifies it
through an LDL$^\top$ factorization.  We use the same
discovery-versus-verification separation, applied to flag-algebra
certificates rather than polynomial inequalities.

\paragraph{Graph limits and Turán densities.}
The mathematical foundation of flag algebras (convergent graph sequences,
positive homomorphisms, the graphon limit) is developed in depth by
Lov\'asz and Szegedy and by Lov\'asz~\cite{lovasz2006limits,lovasz2012large}.
Freer has independently formalized a
substantial graphon theory in Lean~4, including cut distance, weak
regularity,
counting and inverse-counting lemmas, compactness, and the equivalence of two
notions of convergence~\cite{freer2026graphon}.  Our development works with
convergent flag sequences, positive homomorphisms, and random extensions
rather than graphons; connecting the two formalizations into a common formal
account of the flag-algebra/graphon correspondence remains future work.

\section{Conclusion}
\label{sec:conclusion}

We have developed a Lean~4 formalization of Razborov's flag algebra method for
simple graphs.  Its specification layer represents flags as quotients modulo
isomorphism, constructs the quotient algebra of density-expansion relations,
defines positive homomorphisms and the semantic non-negativity cone, and
introduces the ensemble semantic order, which imposes forbidden-subgraph
constraints only when algebra elements are compared.  On top of this
specification, an executable reflection layer computes finite flag data and
connects these computations to the abstract definitions through adequacy
theorems, while a certificate-to-proof compiler uses the resulting identities
to reconstruct Lean proofs from externally generated flag-algebra
certificates.  We evaluated the compiler on seven upper-bound certificates:
Mantel's theorem, bounds on the \(P_3\)- and \(C_4\)-densities of triangle-free
graphs, the Erd\H{o}s pentagon theorem, and edge-density bounds for
\(K_4\)-free, \(K_5\)-free, and \(C_5\)-free graphs.  The formalization also
supports proofs written directly against its mathematical interface: we used
it to prove the matching lower bounds for Mantel's theorem and the
Erd\H{o}s pentagon theorem, as well as two inequalities of Goodman.  Finally,
we summarized a separately developed mathematical meta-theory comparing the
ensemble and quotient semantics.  Non-negativity in the quotient semantics
always implies non-negativity in the ensemble semantics, while the converse
holds for every typed expression exactly when the constrained class is
root-plantable at that flag type.

We do not claim to have found the uniquely right way to formalize flag
algebras.  Our aim has instead been to make the development's design choices
and their consequences explicit: the engineering obstacles they create, the
executable and automation layers needed to address those obstacles, the proof
patterns that may transfer to other formalizations, and the questions that
remain open.  Some engineering decisions reflect what is computationally
feasible in the current implementation rather than a principled design
preference.  Most notably, the two largest certificates still use
\lean{native\_decide}; we have stated this trust tradeoff explicitly.

We expect the architecture to support additional parts of Razborov's
flag-algebra calculus, especially the differential structure developed in
Section~4.3 of his original paper~\cite{razborov2007flag}.  Formalizing this
structure would test how much of the existing quotient-level semantics and
finite-counting infrastructure can be reused beyond the certificate arguments
considered here.  We have also not tested whether the architecture extends to
richer combinatorial structures, such as hypergraphs, directed graphs, and
oriented graphs.  Other directions for future work include building a formal
bridge between our positive-homomorphism semantics and graphon theory.  In
particular, the representation of every positive homomorphism by a
graphon would connect our development more directly to the Lov\'asz
theory~\cite{lovasz2012large} and to the existing Lean graphon
formalization~\cite{freer2026graphon}.  Automating the discovery of SDP
certificates within Lean, rather than importing them from external solvers, is
a longer-term goal.  More immediately, extending the \lean{decide +kernel}
path from the five certificates it currently covers to the two certificates
that still use \lean{native\_decide}, or replacing the latter with a formally
verified external checker, would eliminate the remaining dependency on Lean's
native compiler.  This dependency is a feasibility constraint of the current
encoding, not a logical limitation of the method.

\subsection*{Use of AI Assistance}
The layers of this development used AI assistance to different degrees.  The
specification and reflection layers of \Cref{sec:abstract,sec:reflection},
including the quotient constructions, their interface theorems, and the
adequacy proofs, were formalized manually by the authors.  In the
certificate-to-proof compiler of \Cref{sec:flagmatic}, most of the required
metaprogramming, including the elaboration-time generation commands, the custom
tactics, and the certificate compiler, was developed with the help of
Claude Code.  Finally, the mathematical results of the meta-theory
summarized in \Cref{sec:metatheory} were obtained with the help of GPT~Pro,
and their Lean development was autoformalized entirely by Claude Code.  We
also used Codex and Claude Code to assist with writing and revising this paper.

\section*{Acknowledgments}
We would like to thank Jineon Baek, Taeyoung Kim, Hyunwoo Lee, Joonkyung Lee, Christoph Spiegel, 
and Jan Volec for helpful discussions on Razborov's flag algebras and their Lean formalization.  
We are also grateful to
Ross Kang and Sidharth Hariharan for encouraging us to work harder to remove
\lean{native\_decide} in our Lean formalization.
This work was supported by the National Research Foundation of Korea (NRF)
grant funded by the Korean Government (MSIT) (No.~RS-2023-00279680) and by the
Institute for Basic Science (IBS-R029-C1).

\commentout{
\appendix
\section{Lessons for Future Formalizations}
\label{sec:lessons}

This section distills six lessons from concrete design choices, successful
proof patterns, and failed approaches in the development.  The first four
concern representation and automation, the fifth concerns constrained
semantics, and the sixth concerns reporting the trusted base.  They are
case-study observations intended to inform future formalizations, not claims
that the same choices are optimal in every proof assistant or application.

\paragraph{Lesson 1 (Keeping quotient objects as quotients localizes
representation costs).}
Flag algebras are naturally quotient objects: graphs up to isomorphism, and
flag-algebra elements modulo the density-expansion relations.  A
formalization could pick canonical representatives in the specification
layer to make more operations directly executable.  Our specification instead
keeps the quotients, while the reflection layer uses the concrete
\lean{Sym2Graph} representation.  Density adequacy theorems connect the two,
factoring through the count equality
\lean{labeledGraphListCount\_eq\_sym2InducedSubgraphListCount}.

This design concentrates quotient and representation-specific subgraph
reasoning in the adequacy proofs instead of repeating it in thousands of
generated density lemmas.  The adequacy proofs are themselves substantial,
so this is a tradeoff rather than a free abstraction.  The lesson supported by
our case study is narrow: when computation is supplied by a concrete
refinement of a quotient specification, an adequacy boundary can localize
representation costs while leaving later algebraic statements independent of
the chosen representation.

\paragraph{Lesson 2 (Computable is not the same as usable).}
The naive reflective design is to leave the abstract quotient definitions in
place, give the relevant finite types \lean{Fintype} instances, and ask
\lean{native\_decide} to evaluate density goals by unfolding
\lean{flagDensity}.  This is computable in principle but not viable for the
pentagon proof: each density check repeatedly constructs abstract subgraphs,
compares quotient classes by searching for label-preserving isomorphisms, and
carries proof-valued fields through the evaluator.  The graphs are small, but
the computation induced by their specification is not.

What worked was to treat reflection as verified compilation.  The abstract
definitions remain the specification; the reflected isomorphism and density
procedures form a compilation target with deliberate optimizations (edge-count
pruning before permutation enumeration; type-aware permutation that fixes
labeled vertices and enumerates only the remainder).  Adequacy theorems certify
that the executable procedures preserve the specification's meaning.  The
lesson is that decidability alone says little about feasibility: the data
representation presented to the evaluator can determine whether reflection is
practical.  A separate verified implementation can improve that representation
without replacing the mathematical specification.

\paragraph{Lesson 3 (Generated names can carry metadata for automation).}
Every flag and flag-algebra constant in the development is named according
to the scheme \lean{Flag\_n\_k\_m\_i} and \lean{FlagAlgebra\_n\_k\_m\_i},
where $n$ is the vertex count, $k$ and $m$ identify the flag type, and $i$
is a canonical index.  The convention is not merely organizational.  While
Lean processes a proof script and turns tactic output into typed terms, the
sorting tactics inspect each constant's syntactic \lean{Name} and parse its
trailing index.  The tactic \lean{ac\_sort\_at\_pipeline} uses that index as a
sort key without unfolding the flag or comparing graph structures.

A recorded counterexample from the pentagon proof is a commented-out call that
reads \texttt{-- sort\_at -- takes more than 10 minutes}.
The generic \lean{sort\_at}, which uses \lean{simp} without index guidance,
does not finish promptly on that expression.  The index-guided
\lean{ac\_sort\_at\_pipeline} completes the same step by exploiting the name
encoding to bypass that search.  Names are not inherently better than user
attributes or a custom expression type.  The narrower lesson is that, for
generated declarations with a stable naming scheme, names can provide a
lightweight metadata channel for proof automation.

\paragraph{Lesson 4 (External computation is a candidate generator, not a
proof).}
The rational certificate produced by the SDP workflow, and the Python
compiler script that turns it into Lean source, are both unverified.  Neither is
trusted for the logical validity of a theorem accepted by Lean.  The
script emits flags, matrices, factorizations, and a target inequality;
Lean then checks the resulting PSD obligations, density and multiplication
identities, algebraic calculation, and final statement.  Candidate data that
do not support that statement cause proof construction to fail.

The PSD data is checked separately: the rational matrix and its claimed
$LDL^\top$ factorization are written explicitly, and Lean verifies both the
non-negativity of the diagonal and the matrix identity before using the
quadratic form.  There remains a separate fidelity boundary: Lean does not
prove that the generated declarations faithfully transcribe the source
certificate, so a translation error could produce a valid theorem different
from the intended one.  The lesson is to separate these questions.  External
tools can propose witnesses and proof structure without becoming authorities
for logical validity, while fidelity to their source data must be audited or
verified separately (\Cref{sec:flagmatic-trust}).

\paragraph{Lesson 5 (Moving a constraint can change intermediate semantics).}
Razborov's presentation fixes a universal theory of forbidden graphs before
constructing the constrained algebra.  Our development instead keeps an
ambient algebra and introduces the forbidden-graph assumption when an
inequality is stated.  This deferral makes unconstrained algebraic lemmas
reusable, and the transfer theorem recovers the final unlabeled Tur\'an bounds
used in our case studies.

The change is not merely organizational.  As \Cref{sec:forbidden} explains,
the two approaches need not give the same meaning to intermediate statements
about labeled flags.  The lesson is not that constraints should always be
deferred, but that moving a constraint across an abstraction boundary can
change the meaning of such statements.  A formalization should define the
resulting comparison explicitly and prove the relationship it needs rather
than treating deferral as a semantics-preserving refactoring.

\paragraph{Lesson 6 (Report the trust boundary for each proof obligation).}
Our formalization uses two evaluators with different trusted-computing-base
profiles:
\lean{native\_decide} compiles and runs a decidable proposition using Lean's
native compiler, whereas \lean{decide +kernel} performs the evaluation inside
the kernel.  This difference matters even though both commands close Lean
goals.

Five of the seven certificate proofs, the pentagon among them, discharge
every finite computation by \lean{decide +kernel}; for the pentagon this
means 15 batch equalities checked by the kernel, from which Lean derives the
2,832 individual pair-density theorems by projection.  In the $K_5$-free and
$C_5$-free edge-density proofs, kernel evaluation of the generated
computations is not practical at the current scale, so those files use
\lean{native\_decide}; their SDP matrix equalities are small enough for
kernel evaluation and remain discharged by \lean{decide +kernel}.

This asymmetry is a feasibility outcome, not a general recommendation about
where trust should lie.  Reporting it lets a reader distinguish the
kernel-checked PSD step from the bulk density computations that additionally
trust the native compiler.  More generally, a trust profile should identify
the mechanism used for each class of obligation, even when performance rather
than principle determined the placement.
}

\bibliographystyle{ACM-Reference-Format}
\begin{thebibliography}{99}

\bibitem{razborov2007flag}
A.~A.~Razborov.
\newblock Flag algebras.
\newblock \textit{Journal of Symbolic Logic}, 72(4):1239--1282, 2007.
  \DOI{10.2178/jsl/1203350785}.

\bibitem{razborov2013flag}
A.~A.~Razborov.
\newblock Flag algebras: An interim report.
\newblock In R.~L.~Graham, J.~Ne\v{s}et\v{r}il, and S.~Butler, editors,
  \textit{The Mathematics of Paul Erd\H{o}s II}, 2nd ed., pages 207--232.
  Springer, 2013. \DOI{10.1007/978-1-4614-7254-4\_16}.

\bibitem{grzesik2012}
A.~Grzesik.
\newblock On the maximum number of five-cycles in a triangle-free graph.
\newblock \textit{Journal of Combinatorial Theory, Series B},
  102(5):1061--1066, 2012. \DOI{10.1016/j.jctb.2012.04.001}.

\bibitem{hatami2012}
H.~Hatami, J.~Hladk\'{y}, D.~Kr\'{a}l', S.~Norine, and A.~Razborov.
\newblock On the number of pentagons in triangle-free graphs.
\newblock \textit{Journal of Combinatorial Theory, Series A},
  120(3):722--732, 2013. \DOI{10.1016/j.jcta.2012.12.008}.

\bibitem{razborov2008triangles}
A.~A.~Razborov.
\newblock On the minimal density of triangles in graphs.
\newblock \textit{Combinatorics, Probability and Computing}, 17(4):603--618,
  2008. \DOI{10.1017/S0963548308009085}.

\bibitem{razborov2010tetrahedron}
A.~A.~Razborov.
\newblock On 3-hypergraphs with forbidden 4-vertex configurations.
\newblock \textit{SIAM Journal on Discrete Mathematics}, 24(3):946--963, 2010.
  \DOI{10.1137/090747476}.

\bibitem{baber2012turan}
R.~Baber and J.~Talbot.
\newblock New Tur\'{a}n densities for 3-graphs.
\newblock \textit{Electronic Journal of Combinatorics}, 19(2):\#P22, 2012.
  \DOI{10.37236/2360}.

\bibitem{balogh2016induced5}
J.~Balogh, P.~Hu, B.~Lidick\'{y}, and F.~Pfender.
\newblock Maximum density of induced 5-cycle is achieved by an iterated blow-up
  of 5-cycle.
\newblock \textit{European Journal of Combinatorics}, 52:47--58, 2016.
  \DOI{10.1016/j.ejc.2015.08.006}.

\bibitem{balogh2017rainbow}
J.~Balogh, P.~Hu, B.~Lidick\'{y}, F.~Pfender, J.~Volec, and M.~Young.
\newblock Rainbow triangles in three-colored graphs.
\newblock \textit{Journal of Combinatorial Theory, Series B}, 126:83--113,
  2017. \DOI{10.1016/j.jctb.2017.04.002}.

\bibitem{dillies2022szemeredi}
Y.~Dillies and B.~Mehta.
\newblock Formalising Szemer\'{e}di's regularity lemma in Lean.
\newblock In \textit{Proc.\ ITP 2022}, LIPIcs~237, pages 9:1--9:19, 2022.
  \DOI{10.4230/LIPIcs.ITP.2022.9}.

\bibitem{mantel}
W.~Mantel, Problem 28, Wiskundige Opgaven, 10:60--61, 1907.

\bibitem{goodman1959}
A.~W.~Goodman.
\newblock On sets of acquaintances and strangers at any party.
\newblock \textit{The American Mathematical Monthly}, 66(9):778--783, 1959.
  \DOI{10.2307/2310464}.

\bibitem{mehta2022kruskal}
B.~Mehta.
\newblock Formalising the Kruskal--Katona theorem in Lean.
\newblock In \textit{Intelligent Computer Mathematics (CICM 2022)},
  LNCS~13467, pages 75--91, 2022.
  \DOI{10.1007/978-3-031-16681-5\_5}.

\bibitem{subercaseaux2024hexagon}
B.~Subercaseaux, W.~Nawrocki, J.~Gallicchio, C.~Codel, M.~Carneiro, and
  M.~J.~H.~Heule.
\newblock Formal verification of the empty hexagon number.
\newblock In \textit{Proc.\ ITP 2024}, LIPIcs~309, pages 35:1--35:19, 2024.
  \DOI{10.4230/LIPIcs.ITP.2024.35}.

\bibitem{gowers2025marton}
W.~T.~Gowers, B.~Green, F.~Manners, and T.~Tao.
\newblock On a conjecture of Marton.
\newblock \textit{Annals of Mathematics}, 201(2):515--549, 2025.
  \DOI{10.4007/annals.2025.201.2.5}.

\bibitem{pfr2023}
A.~Anderson et al.
\newblock Formalization of the polynomial Freiman--Ruzsa conjecture of Marton.
\newblock Formalization project, version 0.1.0, 2023.
  \url{https://teorth.github.io/pfr/}.

\bibitem{vaughan2012flagmaticguide}
E.~R.~Vaughan.
\newblock Flagmatic user's guide, version 1.5, 2012.
\newblock Distributed with the Flagmatic software; archived copy of the
  distribution site at
  \url{https://web.archive.org/web/20160112235442/http://flagmatic.org/}.

\bibitem{vaughan2013flagmatic}
E.~R.~Vaughan.
\newblock Flagmatic 2.0.
\newblock Computer software, 2013. Archived copy available at
  \url{https://github.com/jsliacan/flagmatic-2.0}.

\bibitem{chlipala2013cpdt}
A.~Chlipala.
\newblock \textit{Certified Programming with Dependent Types: A Pragmatic
  Introduction to the Coq Proof Assistant}.
\newblock MIT Press, 2013.

\bibitem{cohen2013refinements}
C.~Cohen, M.~D\'{e}n\`{e}s, and A.~M\"{o}rtberg.
\newblock Refinements for free!
\newblock In \textit{Proc.\ CPP 2013}, LNCS~8307, pages 147--162, 2013.
  \DOI{10.1007/978-3-319-03545-1\_10}.

\bibitem{harrison2007verifying}
J.~Harrison.
\newblock Verifying nonlinear real formulas via sums of squares.
\newblock In \textit{Proc.\ TPHOLs 2007}, LNCS~4732, pages 102--118, 2007.
  \DOI{10.1007/978-3-540-74591-4\_9}.

\bibitem{parrilo2003minimizing}
P.~A.~Parrilo and B.~Sturmfels.
\newblock Minimizing polynomial functions.
\newblock In \textit{Algorithmic and Quantitative Real Algebraic Geometry},
  DIMACS Series in Discrete Mathematics and Theoretical Computer Science,
  volume~60, pages 83--99, 2003.

\bibitem{morrison2026sos}
K.~Morrison.
\newblock Harrison's sum-of-squares decision procedure for nonlinear real
  arithmetic in Lean~4.
\newblock Computer software, 2026.
  \url{https://github.com/leanprover/sos}.

\bibitem{ebner2017structured}
G.~Ebner, S.~Ullrich, J.~Roesch, J.~Avigad, and L.~de~Moura.
\newblock A metaprogramming framework for formal verification.
\newblock \textit{Proc.\ ACM Program.\ Lang.}, 1(ICFP):34:1--34:29, 2017.
  \DOI{10.1145/3110278}.

\bibitem{lovasz2012large}
L.~Lov\'{a}sz.
\newblock \textit{Large Networks and Graph Limits}.
\newblock American Mathematical Society Colloquium Publications, volume~60,
  2012.

\bibitem{lovasz2006limits}
L.~Lov\'{a}sz and B.~Szegedy.
\newblock Limits of dense graph sequences.
\newblock \textit{Journal of Combinatorial Theory, Series~B}, 96(6):933--957,
  2006.
  \DOI{10.1016/j.jctb.2006.05.002}.

\bibitem{freer2026graphon}
C.~Freer.
\newblock Graphon theory in Lean~4.
\newblock Computer software, 2026.
  \url{https://github.com/cameronfreer/graphon}.

\bibitem{davey2026local}
E.~Davey, E.~Hurley, R.~de~Joannis~de~Verclos, R.~J.~Kang, and J.~Volec.
\newblock Local flag algebras.
\newblock arXiv:2607.12461, 2026.

\bibitem{davey2026strong}
E.~Davey, E.~Hurley, R.~de~Joannis~de~Verclos, R.~J.~Kang, and J.~Volec.
\newblock Strong edge-colouring via local flag algebras.
\newblock arXiv:2607.17421, 2026.

\bibitem{spiegel2026computational}
C.~Spiegel.
\newblock Computational flag algebras in Lean.
\newblock Talk at the workshop \textit{Flags in the Mountains}, CSU Mountain
  Campus, June 2026.  Slides at
  \url{https://christophspiegel.berlin/assets/slides/FlagsMountains_2026.pdf}.

\end{thebibliography}

\appendix
\section{Compilation Times}
\label{sec:compile-times}

Table~\ref{tab:compile-times} reports end-to-end compilation time and peak
memory for all seven compiler examples of \Cref{sec:results}, compiled with
Lean~4 (v4.27.0) on an
Intel Core i7-14700K (20 cores) with 64\,GB of RAM under Windows~11.  Each
measurement covers the complete source file: flag enumeration, pair-density and
multiplication lemma generation, PSD certificate checking, the objective
expansion, and the main theorem.  The generated files differ only in how their
bridging lemmas are discharged.  Under \lean{decide +kernel} (the
\lean{flagGen.kernelDecide} option), the file introduces no compiled-evaluation
axiom; under \lean{native\_decide}, it is faster but additionally depends on the
\lean{Lean.ofReduceBool} and \lean{Lean.trustCompiler} axioms.  We report both
modes for the five examples where kernel evaluation is tractable; for the two
largest, the $K_5$- and $C_5$-free edge-density examples, only
\lean{native\_decide} was measured.

Each figure is the mean over five runs of a single-file compilation against the
prebuilt library, with the sample standard deviation; memory is the peak
resident set, taken as the maximum over the five runs.  Wall-clock times are
highly repeatable (standard deviation below $0.5\%$ of the mean in every
configuration except the Erd\H{o}s pentagon under kernel evaluation, at
$3.9\%$).

\begin{table}[h]
  \centering
  \caption{End-to-end Lean compilation time (mean of five runs, $\pm$ sample
    standard deviation) and peak resident memory (maximum over five runs) for
    all seven examples, under \lean{decide +kernel} and \lean{native\_decide}.
    ``Lemmas'' is the total count of generated \lean{flagDensity}$_2$
    lemmas.  A dash marks a configuration not measured under kernel evaluation.}
  \label{tab:compile-times}
  \setlength{\tabcolsep}{4.5pt}
  \begin{tabular}{lccr rr rr}
    \toprule
    & & & & \multicolumn{2}{c}{\lean{decide +kernel}}
        & \multicolumn{2}{c}{\lean{native\_decide}} \\
    \cmidrule(lr){5-6}\cmidrule(lr){7-8}
    Case & Forbid & $N$ & Lemmas
      & Time (s) & RAM (GB) & Time (s) & RAM (GB) \\
    \midrule
    Mantel (edge density)              & $K_3$ & 3 &    15 & $14.1\pm0.3$   &  2.4 & $11.5\pm0.1$  & 2.0 \\
    $P_3$ density                      & $K_3$ & 3 &    15 & $13.3\pm0.1$   &  2.4 & $11.1\pm0.0$  & 2.0 \\
    $C_4$ density                      & $K_3$ & 4 &   210 & $74.4\pm0.1$   &  8.2 & $19.8\pm0.2$  & 2.3 \\
    Edge density                       & $K_4$ & 4 &   390 & $133.3\pm0.3$  & 11.2 & $27.9\pm0.1$  & 2.5 \\
    Erd\H{o}s pentagon ($C_5$ density) & $K_3$ & 5 &  2832 & $2786\pm108$   & 42.3 & $123.4\pm0.5$ & 4.2 \\
    Edge density                       & $K_5$ & 5 & 10332 & ---            & ---  & $342.6\pm0.6$ & 8.1 \\
    Edge density                       & $C_5$ & 5 &  8172 & ---            & ---  & $539.6\pm0.5$ & 8.6 \\
    \bottomrule
  \end{tabular}
\end{table}

\end{document}
